\chapter{Logarithmic vertex algebras: the triplet and doublet algebras}\label{ch:L18}
\begingroup\setlength{\emergencystretch}{3em}

\providecommand{\leighteenvac}{\mathbf{1}}
\providecommand{\leighteenV}{\mathcal{V}}
\providecommand{\leighteenW}{\mathcal{W}}
\providecommand{\leighteenA}{\mathcal{A}}
\providecommand{\leighteenVL}{V_{L_p}}
\providecommand{\leighteenQ}{Q}
\providecommand{\leighteenO}{O}
\providecommand{\leighteenAtwo}{A_2}
\providecommand{\leighteenbar}[1]{\overline{#1}}
\providecommand{\leighteenSF}{\mathrm{SF}}
\providecommand{\leighteenHS}{T}

This chapter documents the release of the LeanVOA library named in the front matter. Checks that were run at an earlier state of the library and not re-run for this edition are
marked as such. Every Lean quotation below is taken from the files of the release and is
machine-checked by \texttt{scripts/textbook/sync.py} to occur verbatim (up to whitespace and
comments) in the source; proof bodies are not reproduced, and a quotation that ends in \texttt{:=}
or \texttt{:= by} is the statement only. The chapter covers the files assigned to it in
\texttt{Notes/Textbook/assignment.txt}: in \leanpath{VOA/\allowbreak Examples/} the files
\leanpath{TripletLattice.lean}, \leanpath{TripletScreening.lean}, \leanpath{TripletTTNull.lean},
\leanpath{TripletIrrational.lean}, \leanpath{TripletLogModular.lean}, \leanpath{TripletW2*.lean},
\leanpath{TripletWp*.lean}, \leanpath{Doublet*.lean} and \leanpath{ZhuTripletW2*.lean}; in
\leanpath{VOA/\allowbreak Modularity/} the files \leanpath{LogarithmicPseudoTrace*.lean},
\leanpath{LogarithmicQSeries.lean} and \leanpath{PseudoTrace.lean}. That is 182 files and roughly
thirty-three thousand lines, but about 140 of the files are the pieces of two large exact computations
(a sparse polynomial identity in a Fock space and a sparse certificate in Zhu's algebra), each of
which is consumed by a single theorem; we explain what those computations are, and quote their
interfaces, but do not tour them. (The release audit of 2026-10-27 re-elaborated the headline statements on branch
\texttt{deligne-lean} at an earlier commit, where the \leanpath{VOA/} sources were unchanged since the audit's start;
they still match the quotations below. The quotations were re-verified at that commit (not re-run for this edition);
see Section~\ref{l18:sec:verification}.)

A few results of this chapter were checked in Lean by short checks written against the library. These
checks are not part of the library or of this release; every theorem in them depends only on
\lean{[propext, Classical.choice, Quot.sound]}.

Chapter~\ref{ch:01} is the Lean toolkit of the book; generic Lean idioms (\lean{Submodule},
\lean{finrank}, \lean{↥}, \lean{→ₗ[ℂ]}, structures and their fields, \lean{Prop}-valued
definitions used as named hypotheses) are explained there and are not explained again. We rely on
the following accepted chapters and do not repeat them: Chapter~\ref{ch:L04} (subalgebras
generated by a set, screening operators), Chapter~\ref{ch:L05} (lattice vertex algebras,
background charge, quadratic refinements), Chapter~\ref{ch:L06} (the Virasoro algebra, Kac
determinant, screenings), Chapter~\ref{ch:L11} ($C_2$, presentations), Chapter~\ref{ch:L12}
(modules of all kinds, Zhu's algebra, $C_2$-cofiniteness, the rationality hierarchy),
Chapter~\ref{ch:L13} (the Dong--Jiang induced module), Chapters~\ref{ch:L14}
and~\ref{ch:L15} (the modular toolkit and Zhu's modularity theorem) and Chapter~\ref{ch:L17} (super
vertex algebras, the even part, symplectic fermions).

%======================================================================
\section{What this chapter is about, and what it found}\label{l18:sec:intro}
%======================================================================

The rational vertex operator algebras --- lattice VOAs, affine VOAs at positive integer level,
the Virasoro minimal models, the Moonshine module --- have a semisimple representation theory: a
finite list of irreducible modules, every module a direct sum of them, and characters that span a
representation of $SL_2(\Z)$ (Zhu's theorem, Chapter~\ref{ch:L15}). The simplest vertex operator
algebras that are \emph{not} of this kind, yet are still ``finite'' in the sense of
$C_2$-cofiniteness, are the \emph{triplet algebras} $W(p)$, $p\ge2$, of Kausch (1991). They are
the chiral algebras of the logarithmic conformal field theories of central charge
\begin{equation}\label{l18:eq:cp}
c_{1,p}=1-\frac{6(p-1)^2}{p},\qquad c_{1,2}=-2,\quad c_{1,3}=-7,\quad c_{1,4}=-\tfrac{25}{2},\ \dots
\end{equation}
This chapter documents what the library proves about them, and about two neighbouring objects:
the \emph{doublet algebra} $\leighteenA(p)$, and the \emph{pseudo-trace} construction which replaces
traces in the modular theory of non-semisimple module categories.

\paragraph{Two different Lean objects called ``$W(2)$''.} The single most important thing a reader
must keep in mind is that the library contains \emph{two} constructions of the $p=2$ triplet
algebra, on different carrier types, and \emph{no isomorphism between them is proved}:
\begin{itemize}
\item \lean{VOA.LatGen.Wp.tripletWVOA 2 hp} --- the \emph{lattice model}: the vertex operator
  subalgebra, generated by four explicit states, of the rank-one lattice VOA $V_{L_2}$,
  $L_2=\Z\alpha$, $\langle\alpha,\alpha\rangle=4$, deformed by a background charge. It exists for
  every $p\ge1$ as \lean{tripletWVOA p hp} (Section~\ref{l18:sec:lattice}).
\item \lean{VOA.SympFerm.tripletW2} --- the \emph{Abe model}: the even part of the
  symplectic-fermion super vertex operator algebra of Chapter~\ref{ch:L17}
  (Section~\ref{l18:sec:abe}). The name ``Abe model'' refers to T.~Abe's 2007 paper, which studies
  $W(2)$ in exactly this realization $\leighteenSF(1)^+$. Concretely: the symplectic-fermion
  algebra is generated by two fermionic fields of weight one; its \emph{even part} is the span of
  the states containing an even number of fermion modes, and that even part, with its own vertex
  operator algebra structure, is the Abe model. The formal definition, with its Lean source, is
  Definition~\ref{l18:def:tripletW2}; until then, ``Abe model'' means exactly this sentence.
\end{itemize}
In the literature these are the same algebra (Kausch; Gaberdiel--Kausch; Abe; Adamovi\'c--Milas).
In Lean they are unrelated, and \emph{the theorems proved about them are different}. The only formal
link is the equality of the two central charges, $-2=-2$
(\lean{tripletW_centralCharge_eq_tripletW2}, Theorem~\ref{l18:thm:ccmatch}). Table~\ref{l18:tab:two}
summarises what is proved about each.

\paragraph{Vocabulary used in the table and the headlines.} The table and the list of headline
results below use a number of notions whose home is an earlier chapter. So that this section can be
read on its own, here is a one-sentence definition of each, with a pointer to the place where it is
developed in full (with its Lean source and fidelity check). Throughout, $V$ is a vertex operator
algebra with conformal vector $\omega$ and $L_0=\omega_{(1)}$.
\begin{itemize}
\item \emph{Generalized $L_0$-eigenspace.} For a linear operator $L_0$ on a space $M$ and
  $h\in\C$, the generalized eigenspace of eigenvalue $h$ is
  $M_{[h]}=\{v\in M\mid (L_0-h)^k v=0\text{ for some }k\ge1\}$. It contains the ordinary eigenspace
  $\{v\mid L_0v=hv\}$ and equals it exactly when $L_0$ is diagonalisable on $M_{[h]}$. On
  $M_{[h]}$ one splits $L_0=h+L_0^{\mathrm{nil}}$ with $L_0^{\mathrm{nil}}=L_0-h$ nilpotent.
\item \emph{Jordan block.} A Jordan block of size two at eigenvalue $h$ is a pair of vectors
  $\phi,\psi$ on which $L_0$ acts by the matrix $\left(\begin{smallmatrix}h&1\\0&h\end{smallmatrix}\right)$,
  i.e.\ $L_0\phi=h\phi$ and $L_0\psi=h\psi+\phi$; then $\psi$ lies in the generalized eigenspace but
  is not an eigenvector. Any nilpotent operator on a finite-dimensional space is, in a suitable basis,
  a direct sum of such blocks (of various sizes); $L_0^{\mathrm{nil}}\neq0$ means at least one block
  has size $\ge2$.
\item \emph{Ordinary module} (Chapter~\ref{ch:L12}, Definition~\ref{l12:def:ordmodule}). A
  $V$-module on which $L_0$ is \emph{diagonalisable}, with finite-dimensional eigenspaces and
  eigenvalues bounded below in real part. These are the modules of rational CFT.
\item \emph{Admissible module} (Chapter~\ref{ch:L12}, Definition~\ref{l12:def:admissible}; also
  Chapter~\ref{ch:L09}, Section~\ref{l09:sec:admissible} for the affine use of the word). A module
  graded by $\N$ (``by the mode degree'') so that every mode $u_{(n)}$ of a homogeneous $u$ shifts the
  degree by $\wt u-n-1$; $L_0$ is \emph{not} required to be diagonalisable. Every ordinary module
  carries such a grading (Chapter~\ref{ch:L12}, Theorem~\ref{l12:thm:ordadm}).
\item \emph{Admissibly rational} (Chapter~\ref{ch:L12}, Definition~\ref{l12:def:admrat}). $V$ is
  admissibly rational if \emph{every admissible module is completely reducible}, i.e.\ a direct sum of
  irreducible ones. A logarithmic module (below) is admissible but not a direct sum of irreducibles,
  so its existence is exactly the failure of this property.
\item \emph{Gannon-rational} (Chapter~\ref{ch:L12}, Definition~\ref{l12:def:gannon}). The variant of
  rationality used in Gannon's book: every ordinary module is completely reducible, there are finitely
  many irreducible ordinary modules up to isomorphism, and every irreducible admissible module is
  ordinary.
\item \emph{$C_2$-cofinite} (Chapter~\ref{ch:L12}, Definition~\ref{l12:def:c2}). The quotient
  $R_V=V/C_2(V)$, where $C_2(V)$ is the span of all $u_{(-2)}v$, is finite-dimensional. This is the
  ``finiteness'' that replaces rationality in the logarithmic world.
\item \emph{Strongly generated} (made precise in Section~\ref{l18:sec:conv}). $V$ is strongly
  generated by a set $S$ if every state is a linear combination of ``words''
  $s^1_{(-m_1)}\cdots s^k_{(-m_k)}\leighteenvac$ with $s^j\in S$ and all $m_j\ge1$; physically, every
  field is a normally ordered polynomial in the fields of $S$ and their derivatives.
\item \emph{Strongly finite} (Chapter~\ref{ch:L12}, Definition~\ref{l12:def:strong}). $V$ is simple,
  self-dual (isomorphic to its contragredient module) and $C_2$-cofinite (CFT type being automatic in
  the library's setting): the hypotheses under which the
  modern non-semisimple (``logarithmic'') theory of Creutzig--Gannon and McRae works.
\item \emph{Logarithmic module} (Chapter~\ref{ch:L12}, Definition~\ref{l12:def:log}). A module that
  is the direct sum of its generalized $L_0$-eigenspaces, each finite-dimensional, with weights bounded
  below. The definition allows $L_0^{\mathrm{nil}}=0$ (then the module is ordinary); a module is
  \emph{genuinely} logarithmic when $L_0^{\mathrm{nil}}\neq0$, i.e.\ when it has at least one Jordan
  block of size $\ge2$, and that is what Theorem~\ref{l18:thm:logmodule} produces.
\item \emph{Super vertex operator algebra} (Chapter~\ref{ch:L17}, Definitions~\ref{l17:def:sva}
  and~\ref{l17:def:svoa}). A vertex algebra whose state space is $\Z/2$-graded into an even (bosonic)
  and an odd (fermionic) part, with locality and the Jacobi identity holding up to the Koszul sign
  $(-1)^{|u||v|}$ when two odd fields are exchanged. Its \emph{even part} is an ordinary vertex
  (operator) algebra. Symplectic fermions (Section~\ref{l18:sec:physics} below) are the example used
  here.
\end{itemize}

\begin{table}[ht]
\centering\small
\begin{tabular}{p{4.3cm}p{5.1cm}p{5.1cm}}
\toprule
 & lattice model \lean{tripletWVOA 2 hp} & Abe model \lean{tripletW2}\\
\midrule
VOA structure, $c=-2$ & proved & proved\\
no weight-one states & proved (needs strong generation) & proved (elementary)\\
strongly generated by $\omega,E,H,F$ & \textbf{proved} & open (named \lean{Prop})\\
$C_2$-cofinite & \textbf{proved}, $\dim A_2\le80$ & conditional on one generation gate\\
nilpotency of generators in $A_2$ & proved ($\omega$ via the Adamovi\'c--Milas identity) & proved ($\bar\omega^6=0$)\\
simple & \textbf{proved} (Zhu certificate) & not addressed\\
strongly finite & \textbf{proved} & not addressed\\
not admissibly rational & \textbf{proved} & conditional on \lean{hC2}, \lean{hfive} (plus the vacuously true \lean{hblocks}); \lean{hzhu} discharged\\
not Gannon-rational & \textbf{proved} & not proved\\
$A(V)$ not semisimple & not addressed & conditional on citations\\
logarithmic module exists & not addressed & conditional on \emph{four} hypotheses, one of which (\lean{OmegaJordanAt}) essentially assumes the Jordan block\\
formal link to the other column & \multicolumn{2}{p{10.2cm}}{only the equality of central charges $-2=-2$; no map between the two models}\\
\bottomrule
\end{tabular}
\caption{What the library proves about its two realizations of $W(2)$. ``Proved'' means
\texttt{\#print axioms}-clean (standard three axioms only) with no hypothesis; ``conditional''
means a theorem with explicit named hypotheses. Neither column uses \lean{VOA.pbw}. The unconditional
results are theorems about the \emph{lattice} model only; nothing transfers to the Abe model.}
\label{l18:tab:two}
\end{table}

\paragraph{The headline results, stated honestly.}
\begin{enumerate}
\item \emph{The lattice triplet at $p=2$ is $C_2$-cofinite, simple, strongly finite, and neither
  admissibly rational nor Gannon-rational} --- all unconditionally, \emph{for the lattice model}
  \lean{tripletWVOA 2 hp} only (Theorems~\ref{l18:thm:c2two}, \ref{l18:thm:irrtwo}, \ref{l18:thm:simple}). The
  rationality notions are the library's (Chapter~\ref{ch:L12}): \lean{AdmissiblyRational} asks that every weak submodule
  of every admissible module have a weak complement, which is not literally the complete reducibility of
  Dong--Li--Mason; the witness used here (the ambient $V_{L_2}$, in which $W(2)$ has no complement even among weak
  submodules) refutes the DLM notion as well, on paper. This is the
  library's first non-rational, strongly finite vertex operator algebra; it is the benchmark that
  the logarithmic literature (Creutzig--Gannon, McRae) calls a ``strongly finite but non-rational''
  VOA. The $C_2$ proof rests on an exact polynomial identity in the Fock space of weight ten
  (the $p=2$ case of an identity of Adamovi\'c and Milas, with leading coefficient
  $B_2=1024/1575$), and the simplicity proof on an exact membership certificate in Zhu's algebra.
\item \emph{For general $p$}, $C_2$-cofiniteness of \lean{tripletWVOA p hp} is proved conditionally
  on two named finite computations: a generation statement (\lean{AdjoinGenWp}, or the finite OPE
  table \lean{LoweringPairsWp}) and the nilpotency of $\bar\omega$ in the form
  \lean{OmegaPowMem p hp m} (or the identity \lean{AMOmegaIdentity p hp (3p-1)})
  (Theorem~\ref{l18:thm:c2general}). A Lean check (not part of this release) shows these hypotheses are satisfiable (they
  hold at $p=2$), but for $p\ge3$ they are not discharged: ``$W(p)$ is $C_2$-cofinite'' for general $p$ is a
  \textbf{conditional} reduction, not a theorem. In the course of this the library found that the printed argument of
  Adamovi\'c--Milas (2008, Theorem~3.1) for the nilpotency of $\bar\omega$ is a non-sequitur, and
  routes around it through their own Corollary~4.6 (Warning in Section~\ref{l18:sec:c2general}).
\item \emph{For the Abe model} (a separate Lean object, linked to the lattice model only by $c=-2$),
  $C_2$-cofiniteness is reduced to one named generation statement
  (\lean{AbeQuotientGeneration}) and is \textbf{conditional}; all four generators are proved nilpotent in $A_2$; non-rationality
  and non-semisimplicity of Zhu's algebra are proved only \textbf{conditionally} on explicitly cited literature
  theorems (\lean{C2Cofinite}, a Wedderburn block count, and $\dim\ge5$ of a block space, or $\dim A(V)=11$ in place of
  the $C_2$/block-space pair; the fourth input of the non-rationality theorem, the semisimplicity criterion \lean{hzhu},
  is discharged for every VOA by \lean{zhuAdmissibleSemisimple_holds}, Chapter~\ref{ch:L12}); and a logarithmic module (one on which $L_0$ has a nonzero nilpotent
  part) is constructed conditionally on \textbf{four} hypotheses, not one: two finiteness hypotheses on Zhu's algebras
  (\lean{hAnm}, \lean{hA}), the Dong--Jiang gate \lean{DJAnnEqAn tripletW2 0}, and \lean{OmegaJordanAt tripletW2 h}. The
  last literally asserts that $[\omega]-h$ is not semisimple in $A(V)$, so the Jordan block of the module is a formal
  transcription of that hypothesis, which is proved for no VOA and no $h$.
\item \emph{Pseudo-traces} (Miyamoto 2004, in Arike's Hattori--Stallings form) are constructed for
  arbitrary finitely generated projective modules over arbitrary algebras and validated against the
  ordinary trace (Section~\ref{l18:sec:pseudo}); the logarithmic one-point $q$-series built from
  them is polynomial in $\tau$, which is the precise source of the $\tau$-linear terms in
  logarithmic modular transformations. No VOA instance of the required ``commutant data'' is
  constructed, so the genus-one layer applies to no concrete VOA yet.
\item \emph{The doublet} $\leighteenA(p)$ is constructed as a generated super subalgebra of a
  concrete half-lattice ambient (a super near-VOA with states of negative weight), for $4\mid p$ only,
  and proved only bounded below; its identification
  with the literature's doublet and the relation $\leighteenA(p)^{\mathrm{even}}=W(p)$ are not proved
  (Section~\ref{l18:sec:doublet}). ``The doublet algebra is formalized'' would overclaim.
\end{enumerate}

\paragraph{The findings of the fidelity checks, in brief.} Each is explained in a red box where it
occurs.
\begin{itemize}
\item The docstrings of \leanpath{TripletWp.lean} and \leanpath{TripletWpC2.lean} list ``$E\ne0$''
  and ``$W(p)\ne V_{L_p}$'' as unproved. At $p=2$ both follow in a few lines from library theorems;
  we checked them in Lean (Theorem~\ref{l18:thm:Ene0};
  the check is not part of this release). For $p\ge3$ they remain unproved.
\item The library calls $Q=(e^{\alpha})_{(0)}$ the ``short screening'' in
  \leanpath{TripletScreening.lean} and \leanpath{TripletWp.lean}, but the ``long screening'' in
  \leanpath{TripletWpC2.lean}. The literature's name is the latter ($\alpha$ is the longer of the two
  screening momenta $\alpha$ and $-\alpha/p$).
\item The name \lean{tripletW p hp} is mathematically faithful (the subalgebra of $V_{L_p}$
  generated by $\omega,E,H,F$ is the triplet algebra), but the library never proves that it equals
  the kernel of the short screening $\widetilde Q=(e^{-\alpha/p})_{(0)}$, which is the definition
  most papers use.
\item Several hypotheses of the Abe-model irrationality theorems are redundant or vacuously true
  (\lean{hfin} beside \lean{hdim}; \lean{WedderburnBlockCount ℂ (A tripletW2) 4} is true for any
  non-semisimple algebra); none is harmful.
\item \leanpath{TripletWpPairClosure.lean} still says the three mixed coordinates (the three
  explicit products $F_{(0)}H$, $F_{(1)}E$, $F_{(3)}E$ of the lowering generator $F$ with $H$ and $E$,
  expressed as polynomials in the generators; they are the only entries of the generator OPE table
  not handled by a general argument, see Section~\ref{l18:sec:mixed} and Theorem~\ref{l18:thm:mixed})
  and the Adamovi\'c--Milas identity are ``not claimed here''; they are now proved in later files
  (stale docstring).
\item The logarithmic pseudo-trace layer (\lean{LogPseudoTraceCommutantData}) has no instance in
  the library; it is not vacuous (an instance with auxiliary algebra $\C$ exists on paper, recovering
  ordinary traces), but no theorem of that layer is about a concrete VOA.
\end{itemize}

%======================================================================
\section{The physical picture: logarithmic conformal field theory}\label{l18:sec:physics}
%======================================================================

This section is motivation; nothing in it is a claim about the Lean library. The theorems that
formalise pieces of it are pointed to.

\subsection{Why logarithms}
In a two-dimensional conformal field theory the dilatation operator on the plane is $L_0+\bar L_0$.
In unitary theories, and in rational theories generally, $L_0$ is diagonalisable on every module, so
two-point functions of $L_0$-eigenstates of weight $h$ are pure powers,
$\langle\phi(z)\phi(0)\rangle\propto z^{-2h}$. Nothing forces diagonalisability in a non-unitary
theory. If $L_0$ has a Jordan block,
\begin{equation}\label{l18:eq:jordan}
L_0\,\phi=h\,\phi,\qquad L_0\,\psi=h\,\psi+\phi,
\end{equation}
then $L_0$ acts on the span of $\phi,\psi$ by the $2\times2$ \emph{Jordan block}
$\left(\begin{smallmatrix}h&1\\0&h\end{smallmatrix}\right)$ (in the basis $\phi,\psi$, the columns being
the images): $\phi$ is an honest eigenvector, while $\psi$ is only a \emph{generalized} eigenvector,
$(L_0-h)\psi=\phi\neq0$ but $(L_0-h)^2\psi=0$. One calls $\psi$ the \emph{logarithmic partner} of
$\phi$: ``partner'' because $\psi$ is tied to $\phi$ by the off-diagonal entry and the two cannot be
separated into independent eigenstates, ``logarithmic'' because of what this does to correlation
functions. Indeed, the dilatation $z\mapsto\lambda z$ acts on a field by $\lambda^{L_0}$, and on the
Jordan block $\lambda^{L_0}=\lambda^h(1+\log\lambda\,(L_0-h))$, so $\lambda^{L_0}\psi=\lambda^h(\psi+\log\lambda\,\phi)$:
the logarithm of the scale enters the transformation law. Imposing global conformal covariance with this
transformation law (together with the vanishing of $\langle\phi\phi\rangle$ which it implies) forces
\begin{equation*}
\langle\phi(z)\phi(0)\rangle=0,\qquad \langle\phi(z)\psi(0)\rangle=\frac{b}{z^{2h}},\qquad
\langle\psi(z)\psi(0)\rangle=\frac{\theta-2b\log z}{z^{2h}} .
\end{equation*}
This is the origin of the name: correlators acquire logarithms, the term $-2b\log z$ in
$\langle\psi\psi\rangle$, which never occurs in a theory where $L_0$ is diagonalisable. Such theories occur in statistical
mechanics at the points where a model is non-unitary and non-generic: dense polymers and the
two-dimensional sandpile ($c=-2$), critical percolation ($c=0$), the spanning-tree and loop-erased
random-walk models, and they occur in string theory as the $bc$ and $\beta\gamma$ ghost sectors and
as worldsheet theories on supergroups. The constant $b$ is a genuine invariant (the ``logarithmic
coupling'' of Gaberdiel, Mathieu, Ridout, Vasseur and others).

In vertex-algebra language, \eqref{l18:eq:jordan} says that a module $M$ is a \emph{logarithmic}
(grading-restricted generalized) module in the sense of Huang--Lepowsky--Zhang: $M$ is the direct
sum of generalized $L_0$-eigenspaces, and the nilpotent part $L_0^{\mathrm{nil}}=L_0-L_0^{ss}$ is
nonzero. The library's notion is \lean{LogVModule} (Chapter~\ref{ch:L12},
Definition~\ref{l12:def:log}); Theorem~\ref{l18:thm:logmodule} below constructs such a module for
the Abe model, conditionally.

\subsection{The triplet algebra}
Take one free boson $\varphi$ with background charge, so that the stress tensor is
$T=\frac12(\partial\varphi)^2+\frac{\alpha_0}{2}\partial^2\varphi$ and $c=1-3\alpha_0^2$. The vertex
operators $V_\beta=e^{\beta\varphi}$ have weight $\frac12\beta(\beta-\alpha_0)$. There are exactly two
momenta $\alpha_\pm$ with weight one, $\alpha_+\alpha_-=-2$, $\alpha_++\alpha_-=\alpha_0$; their
contour integrals $Q_\pm=\oint V_{\alpha_\pm}$ are the \emph{screening charges}, which commute with
the Virasoro algebra (Chapter~\ref{ch:L06}, Section~\ref{l06:sec:screening}). For
$\alpha_+=\sqrt{2p}$, $\alpha_-=-\sqrt{2/p}$ one gets $c=c_{1,p}$ of \eqref{l18:eq:cp}. Normalise
$\alpha:=\alpha_+$, so that $\langle\alpha,\alpha\rangle=2p$ and $\alpha_-=-\alpha/p$. The lattice
$L_p=\Z\alpha$ is even, and the lattice vertex algebra $V_{L_p}$ (Chapter~\ref{ch:L05}) with the
deformed stress tensor is a vertex operator algebra of central charge $c_{1,p}$; in it,
\begin{equation}\label{l18:eq:weights}
\wt\bigl(e^{m\alpha}\bigr)=pm^2-(p-1)m,\qquad \wt e^{\alpha}=1,\qquad \wt e^{-\alpha}=2p-1 .
\end{equation}
The \emph{long screening} $Q=e^{\alpha}_{(0)}$ acts on $V_{L_p}$; the \emph{short screening}
$\widetilde Q=e^{-\alpha/p}_{(0)}$ maps $V_{L_p}$ to the module $V_{L_p-\alpha/p}$. Following
Kausch, Feigin--Gainutdinov--Semikhatov--Tipunin and Adamovi\'c--Milas, the triplet algebra is
\begin{equation}\label{l18:eq:Wp}
W(p)=\ker_{V_{L_p}}\widetilde Q,
\end{equation}
and it is strongly generated by $\omega$ and the three states
\begin{equation}\label{l18:eq:EHF}
F=e^{-\alpha},\qquad H=QF,\qquad E=Q^2F,\qquad Q^3F=0,
\end{equation}
of weight $2p-1$, which transform as an $\mathfrak{sl}_2$-triplet under the screening algebra (hence
the name). All three are \emph{Virasoro primary}: $L_nE=L_nH=L_nF=0$ for all $n\ge1$ (Chapter~\ref{ch:L04},
Definition~\ref{l04:def:primary}; proved in Theorem~\ref{l18:thm:pairs}), i.e.\ they are highest-weight
vectors for the Virasoro algebra, the fields of a $W$-algebra beyond $T$. For $p=2$ the three generators have weight $3$.

The representation theory of $W(p)$ (Adamovi\'c--Milas 2008, Nagatomo--Tsuchiya) has $2p$
irreducible modules, with lowest weights among the entries of the extended Kac table
\begin{equation}\label{l18:eq:kac}
h_{r,s}=\frac{(pr-s)^2-(p-1)^2}{4p}.
\end{equation}
At $p=2$ these are $h=0,\ 1,\ -\tfrac18,\ \tfrac38$ (the vacuum module; the odd part of the symplectic fermions, whose ground states $\chi^\pm_{-1}\leighteenvac$ have weight $1$; and two ``twisted'' modules). There are finitely many irreducibles,
but the category is \emph{not semisimple}: there are indecomposable projective modules on which
$L_0$ has Jordan blocks of size two. Correspondingly Zhu's algebra $A(W(p))$ has dimension $6p-1$
but is not semisimple ($11=1+4+4+2$ at $p=2$, the last summand being $\C[\varepsilon]/(\varepsilon^2)$).

\subsection{Symplectic fermions}
At $p=2$ there is a second, fermionic, realization. A pair of weight-one fermions $\chi^\pm$ with
$\chi^+(z)\chi^-(w)\sim(z-w)^{-2}$ generates the symplectic-fermion super VOA $\leighteenSF$ at
$c=-2$ (Kausch 2000; Chapter~\ref{ch:L17}, Section~\ref{l17:sec:sf}). It violates spin--statistics
(fermions of integer weight), which is why it is a super VOA and not a VOA. Its even part
$\leighteenSF^+$ (the $\Z/2$-orbifold by fermion number) is $W(2)$; the three triplet generators are
the fermion bilinears $\chi^a_{-2}\chi^b_{-1}|0\rangle$, symmetrised. In statistical mechanics this is
the theory of dense polymers and of the abelian sandpile. Free fields make $W(2)$ computable by
hand, which is why Abe's 2007 analysis works in this model.

\subsection{Characters, modular transformations and pseudo-traces}
The vacuum character of $W(2)$ is
\begin{equation}\label{l18:eq:char}
\chi_0(\tau)=q^{1/12}\cdot\frac12\Bigl[\prod_{n\ge1}(1+q^n)^2+\prod_{n\ge1}(1-q^n)^2\Bigr]
=q^{1/12}\bigl(1+q^2+4q^3+5q^4+8q^5+10q^6+\cdots\bigr),
\end{equation}
and under $S:\tau\mapsto-1/\tau$ it transforms (Flohr; Fuchs--Hwang--Semikhatov--Tipunin;
Flohr--Gaberdiel) as
\begin{equation}\label{l18:eq:Schi}
\chi_0(-1/\tau)=\tfrac14\bigl[\chi_{-1/8}(\tau)-\chi_{3/8}(\tau)\bigr]
+\Bigl(-\frac{i\tau}{2}\Bigr)\bigl[\chi_0(\tau)-\chi_1(\tau)\bigr].
\end{equation}
The term proportional to $\tau$ cannot be a combination of characters, since each character is
$q^h$ times a power series and so is an eigenfunction of $\tau\mapsto\tau+1$. The $SL_2(\Z)$-closure
of the characters therefore needs extra functions with explicit $\tau$-dependence; Miyamoto (2004)
showed that for any $C_2$-cofinite VOA these come from \emph{pseudo-traces}: functionals
$\varphi$ on an auxiliary algebra $P$ commuting with the action, such that
$\mathrm{tr}^{\varphi}_M\,o(v)q^{L_0-c/24}$ makes sense even when $L_0$ is not diagonalisable. On a
generalized eigenspace, $q^{L_0}=q^h e^{2\pi i\tau L_0^{\mathrm{nil}}}=q^h\sum_k\frac{(2\pi i\tau)^k}{k!}
(L_0^{\mathrm{nil}})^k$ is a \emph{polynomial in $\tau$} times $q^h$. Section~\ref{l18:sec:pseudo}
formalises exactly this. Section~\ref{l18:sec:logmodular} formalises the converse direction: a
$\tau$-linear term in $\chi_0(-1/\tau)$ is incompatible with rationality.

\subsection{The doublet}
Halving the lattice, $\frac12L_p=\Z\frac\alpha2$, gives two states $a^-=e^{-\alpha/2}$ and
$a^+=Qa^-$ of weight $\frac{3p-2}{4}$, which generate the \emph{doublet} algebra $\leighteenA(p)$
(Adamovi\'c--Milas 2009). Its even part for the charge-parity automorphism is $W(p)$. For $p=2$
the doublet generators have weight $1$ and $\leighteenA(2)$ is the symplectic-fermion algebra itself;
for $4\mid p$ the weight is a half-integer and the half lattice is even, so $\leighteenA(p)$ is a
vertex algebra with half-integral weights (another spin--statistics violation).

%======================================================================
\section{Conventions recalled}\label{l18:sec:conv}
%======================================================================

\begin{itemize}
\item \textbf{Modes.} \lean{mode f n = f (-n-1)}: if $Y(u,z)=\sum_n u_{(n)}z^{-n-1}$ then
  \lean{mode (W.Y u) n} is $u_{(n)}$, and \lean{W.smul u n v} is $u_{(n)}v$. Virasoro modes are
  \lean{W.L n} $=\omega_{(n+1)}$, so $L_{-2}=\omega_{(-1)}$, $L_{-1}=T=\omega_{(0)}$,
  $L_0=\omega_{(1)}$. A mode $u_{(n)}$ of a homogeneous $u\in V_k$ maps $V_l$ to $V_{k+l-n-1}$.
\item \textbf{$C_2$ and Zhu.} \lean{O₂ W} is $C_2(V)=\mathrm{span}\{u_{(-2)}v\}$,
  \lean{A₂ W} $=V/C_2(V)$ with its commutative product induced by $u_{(-1)}v$ (written
  \lean{poissonMul}), \lean{mk₂ W u} $=\bar u$ is the class of $u$, and
  \lean{C2Cofinite W} means \lean{FiniteDimensional ℂ (A₂ W)} (Chapter~\ref{ch:L12},
  Definition~\ref{l12:def:c2}). \lean{O W}, \lean{A W} $=V/O(V)$, \lean{zhuMul} and \lean{mk W} are
  Zhu's subspace, algebra, product and quotient map (Section~\ref{l12:sec:zhu}).
\item \textbf{Strong generation.} \lean{iterProdSpan W S} is the span of the words
  $s^1_{(-m_1)}\cdots s^k_{(-m_k)}\leighteenvac$ with $s^j\in S$ and all $m_j\ge1$
  (Chapter~\ref{ch:L12}, Theorem~\ref{l12:thm:fingen}). $V$ is \emph{strongly generated} by $S$ when
  this span is everything.
\item \textbf{Generated subalgebras.} \lean{VOASubalgebra.generated W S} is the smallest vertex
  operator subalgebra containing $S$: it contains $\leighteenvac$ and $\omega$, is closed under all
  products $u_{(n)}v$, $n\in\Z$, and under the homogeneous projections
  (Chapter~\ref{ch:L04}, Definition~\ref{l04:def:generated}; the defect that it always contains
  $\omega$ is harmless here since $\omega$ is one of our generators).
  \lean{generatedVOA W S} is the VOA structure on its carrier.
\item \textbf{Lattices.} \lean{LatModuleGen L} is the state space $\mathcal F\otimes\C[L]$ of the
  lattice vertex algebra of an integral lattice $L$, \lean{latChargeGen L β} is $e^\beta$,
  \lean{latStateGen L q β} is $q\otimes e^\beta$ for a Fock polynomial $q$, \lean{latVacuumGen L}
  is $\leighteenvac=1\otimes e^0$, and \lean{latticeCocycle L β γ} $=\varepsilon(\beta,\gamma)\in\{\pm1\}$
  is the cocycle (Chapter~\ref{ch:L05}, Section~\ref{l05:sec:general}). \lean{YfullGen} is the
  reconstructed vertex map.
\item \textbf{Background charge.} For $\lambda\in\C^{\mathrm{rank}}$, \lean{BgA.omegaBg L lam}
  $=\omega_\lambda=\omega+\lambda(-2)\leighteenvac$ and \lean{BgA.LBg L lam n} $=L^\lambda_n=L_n-(n+1)\lambda(n)$;
  \lean{BgB.LBG} is the same operator in a second file's conventions, and \lean{bgPair L lam β}
  $=\langle\lambda,\beta\rangle$ (Chapter~\ref{ch:L05}, Definition~\ref{l05:def:latbg}). The weight of
  $q\otimes e^\beta$ for $L^\lambda_0$ is $\deg q+\frac12\langle\beta,\beta\rangle-\langle\lambda,\beta\rangle$.
\item \textbf{Hypotheses \lean{hp : 1 ≤ p}.} Every lattice-model declaration carries
  \lean{(p : ℕ) (hp : 1 ≤ p)}; at $p=2$ the argument \lean{(hp : 1 ≤ 2)} is a proof of a true
  numeral inequality and carries no information. The literature's range is $p\ge2$; at $p=1$ the
  construction still type-checks but describes $V_{L_1}=V_{A_1}$ (see Section~\ref{l18:sec:lattice}).
\end{itemize}

%======================================================================
\section{The lattice model I: the ambient and the generators}\label{l18:sec:lattice}
%======================================================================

\subsection{The deformed rank-one lattice VOA}

\begin{definition}[The lattice $L_p$ and its background charge]\label{l18:def:latP}
For $p\in\N$ let $L_p=\Z\alpha$ be the rank-one lattice with $\langle m\alpha,n\alpha\rangle=2pmn$.
The \emph{background charge} is $\lambda_p=c_p\,\alpha$ with $c_p=\frac{p-1}{2p}$, so that
$\langle\lambda_p,m\alpha\rangle=(p-1)m$ and $\langle\lambda_p,\lambda_p\rangle=\frac{(p-1)^2}{2p}$.
The \emph{deformed weight} of $e^{m\alpha}$ is $q_p(m)=\frac12\langle m\alpha,m\alpha\rangle-\langle\lambda_p,m\alpha\rangle=pm^2-(p-1)m$.
\end{definition}

\noindent File \leanpath{VOA/\allowbreak Examples/\allowbreak TripletLattice.lean}.
\begin{leancode}
noncomputable def latP (p : ℕ) : IntegralLattice where
  carrier := ℤ
  form := ((2 * p : ℕ) : ℤ) • LinearMap.mul ℤ ℤ

def alphaP (p : ℕ) : (latP p).carrier := (1 : ℤ)

noncomputable def cP (p : ℕ) : ℂ := ((p : ℂ) - 1) / (2 * (p : ℂ))

noncomputable def lamP (p : ℕ) : Fin (rank (latP p)) → ℂ :=
  BgA.latBg (latP p) (cP p) (alphaP p)

def qZP (p : ℕ) (m : ℤ) : ℤ := (p : ℤ) * m ^ 2 - ((p : ℤ) - 1) * m
\end{leancode}
\begin{leancode}
theorem bgPair_lamP (p : ℕ) (hp : 1 ≤ p) (m : ℤ) :
    BgA.bgPair (latP p) (lamP p) m = ((p : ℂ) - 1) * (m : ℂ) := by

theorem qZP_eq_halfNorm_sub_bgPair (p : ℕ) (hp : 1 ≤ p) (m : ℤ) :
    ((qZP p m : ℤ) : ℂ)
      = (1 / 2 : ℂ) * (((latP p).form m m : ℤ) : ℂ) - BgA.bgPair (latP p) (lamP p) m := by

theorem qZP_pos (p : ℕ) (hp : 1 ≤ p) {m : ℤ} (hm : m ≠ 0) : 0 < qZP p m := by

theorem centralCharge_latP (p : ℕ) (hp : 1 ≤ p) :
    ((rank (latP p) : ℂ)) - 12 * BgA.bgForm (latP p) (lamP p) (lamP p)
      = 1 - 6 * ((p : ℂ) - 1) ^ 2 / (p : ℂ) := by
\end{leancode}
\begin{dictionary}
\lean{latP p} has carrier $\Z$ (the integer $m$ stands for $m\alpha$, and \lean{alphaP p} $=1$ is
$\alpha$) and Gram form $\langle m,n\rangle=2pmn$ (\lean{latP_form}). It is even for every $p$ and
positive definite for $p\ge1$ (\lean{latP_even}, \lean{latP_posDef}); at $p=1$ it is the $A_1$
lattice (\lean{latP_one_form}). \lean{BgA.latBg L c α} is the frame-coordinate vector of $c\cdot\alpha$
in the orthonormal frame used for background charges (Chapter~\ref{ch:L05},
Definition~\ref{l05:def:latbg}); \lean{bgPair_lamP} confirms $\langle\lambda_p,m\alpha\rangle=(p-1)m$
and \lean{qZP_eq_halfNorm_sub_bgPair} confirms that the integer \lean{qZP p m} is exactly
$\frac12\langle\beta,\beta\rangle-\langle\lambda_p,\beta\rangle$ at $\beta=m\alpha$. The integrality of
$q_p$ is what allows a $\Z$-grading (it is wrapped as a $\Q$-valued quadratic refinement
\lean{qrP p}, Definition~\ref{l05:def:quadrefine}). \lean{centralCharge_latP} is $c=\mathrm{rank}-12\langle\lambda,\lambda\rangle$
evaluated: $1-12\cdot\frac{(p-1)^2}{2p}=c_{1,p}$. The hypothesis \lean{hp : 1 ≤ p} is needed only to
divide by $p$ and for positive definiteness. Everything matches
\eqref{l18:eq:cp}--\eqref{l18:eq:weights}; there is no hidden convention.
\end{dictionary}
\begin{physics}
$L_p$ is the momentum lattice of a compact boson at the radius where the screening momentum
$\alpha_+=\sqrt{2p}$ is a lattice vector; $\lambda_p$ is the background charge
$\frac{\alpha_0}{2}$ of the Coulomb gas, and $q_p(m)$ is the Feigin--Fuchs conformal weight
$\frac12\beta(\beta-\alpha_0)$ at $\beta=m\alpha_+$. Positivity of $q_p$ off $m=0$
(\lean{qZP_pos}) is the reason the deformed lattice theory is still of CFT type, although it is not
unitary.
\end{physics}
\begin{status}
All proved, no hypotheses beyond \lean{hp}; \texttt{\#print axioms} of the instance below is
\lean{[propext, Classical.choice, Quot.sound]}.
\end{status}

\begin{theorem}[The deformed lattice VOA $V_{L_p}$]\label{l18:thm:latVOA}
For $p\ge1$ the lattice vertex algebra $V_{L_p}$, graded by $L^{\lambda_p}_0$ and with conformal
vector $\omega_{\lambda_p}=\omega_{L_p}+\lambda_p(-2)\leighteenvac$, is a vertex operator algebra of
central charge $c_{1,p}$, with weight spaces $(V_{L_p})_k=\bigoplus_{m}\mathcal F_{k-q_p(m)}\otimes e^{m\alpha}$.
In it $e^{\alpha}$ has weight $1$, $e^{-\alpha}$ has weight $2p-1$, and $e^{m\alpha}$ has weight $q_p(m)$.
\end{theorem}
\noindent File \leanpath{VOA/\allowbreak Examples/\allowbreak TripletLattice.lean}.
\begin{leancode}
noncomputable def tripletLatticeVOA (p : ℕ) (hp : 1 ≤ p) :
    VertexOperatorAlgebra (LatModuleGen (latP p)) where
  grading := gradTrip p
  isInternal :=
    gradeCollapse_isInternal (bigradα_iSupIndep (latP p)) (bigradα_iSup_eq_top (latP p)) (wtZP p)
  finrank := fun n => by
    rw [gradTrip_eq]
    exact gradingLatGenQ_finrank (latP p) (qFinite_qrP p hp) _
  Y := YfullGen (latP p) (latP_even p) (latP_posDef p hp)
  vacuum := latVacuumGen (latP p)
...
  conformalVector := BgA.omegaBg (latP p) (lamP p)
...
  centralCharge := 1 - 6 * ((p : ℂ) - 1) ^ 2 / (p : ℂ)
  virasoro := virasoro_gradTrip p hp
\end{leancode}
\begin{leancode}
noncomputable def gradTrip (p : ℕ) (k : ℤ) : Submodule ℂ (LatModuleGen (latP p)) :=
  gradeCollapse (bigradα (latP p)) (wtZP p) k

def wtZP (p : ℕ) (x : ℤ × (latP p).carrier) : ℤ := x.1 + qZP p x.2

theorem tripletLatticeVOA_L (p : ℕ) (hp : 1 ≤ p) (n : ℤ) :
    (tripletLatticeVOA p hp).L n = BgA.LBg (latP p) (lamP p) n :=

theorem latChargeGen_alpha_mem_grade_one (p : ℕ) (hp : 1 ≤ p) :
    latChargeGen (latP p) (alphaP p) ∈ (tripletLatticeVOA p hp).grading 1 := by

theorem latChargeGen_neg_alpha_mem_grade (p : ℕ) (hp : 1 ≤ p) :
    latChargeGen (latP p) (-alphaP p) ∈ (tripletLatticeVOA p hp).grading (2 * (p : ℤ) - 1) := by

theorem latChargeGen_mem_gradTrip (p : ℕ) (m : ℤ) :
    latChargeGen (latP p) m ∈ gradTrip p (qZP p m) := by
\end{leancode}
\begin{dictionary}
The instance has all nineteen fields of \lean{VertexOperatorAlgebra} (Chapter~\ref{ch:01}); the
quotation shows the ones that carry the content, and the ellipses mark the remaining fields, each a
named lemma (\lean{gradTrip_grading_mode}, \lean{L0_gradTrip}, \lean{translation_gradTrip},
\lean{virasoro_gradTrip}, \dots). The decisive points:
\begin{itemize}
\item \lean{Y} is the \emph{undeformed} lattice vertex map \lean{YfullGen} of
  Chapter~\ref{ch:L05}, Theorem~\ref{l05:thm:latticeVOAGen}: a background charge changes the
  stress tensor and the grading, not the operator product. So locality, vacuum and creation axioms
  are inherited verbatim.
\item \lean{grading k} $=$ \lean{gradTrip p k} collects the bigraded pieces \lean{bigradα (w, β)}
  ($=\mathcal F_w\otimes e^\beta$) with $w+q_p(\beta)=k$ (\lean{wtZP}); \lean{gradTrip_eq} identifies
  it with the refined grading \lean{gradingLatGenQ (latP p) (qrP p) k} of Chapter~\ref{ch:L05}.
  Finite-dimensionality of each piece uses \lean{qFinite_qrP} ($|m|\le q_p(m)$).
\item \lean{conformalVector} is $\omega_{\lambda_p}$, and \lean{tripletLatticeVOA_L} says that the
  Virasoro modes of the instance are exactly the deformed $L^{\lambda_p}_n$. The Virasoro relation
  (field \lean{virasoro}) is proved with central charge $1-6(p-1)^2/p$
  (\lean{virasoro_gradTrip}), so the field \lean{centralCharge} is not a free label.
\item Grading axiom (VA1): the deformed weight differs from the undeformed one by $-\langle\lambda_p,\beta\rangle$, a
  linear function of the charge $\beta$; charge is additive under products, so the correction cancels
  in $\wt u+\wt v-n-1$. Translation (VOA3): $L^\lambda_{-1}=L_{-1}$ because the correction
  $-(n+1)\lambda(n)$ vanishes at $n=-1$.
\end{itemize}
The three membership theorems are \eqref{l18:eq:weights} as statements about the grading of the
instance. The ``\lean{hp}'' in \lean{tripletLatticeVOA p hp} matters only through positive
definiteness; at $p=1$ the instance is the undeformed $V_{A_1}$ ($\lambda_1=0$, $c=1$).
\end{dictionary}
\begin{physics}
This is the free-field ``Coulomb gas'' of central charge $c_{1,p}$: one compact boson with
background charge, compactified so that the long screening current $e^{\alpha}$ is a local field. All
the logarithmic theories of this chapter live inside it as kernels of screening charges. Note that
$V_{L_p}$ itself is a perfectly rational lattice theory \emph{as a VOA with its undeformed stress
tensor}; it is the change of $\omega$ that makes the triplet subalgebra appear.
\end{physics}
\begin{status}
Proved, unconditional (only \lean{hp}). \texttt{\#print axioms VOA.LatGen.Triplet.tripletLatticeVOA}:
\lean{[propext, Classical.choice, Quot.sound]}. No \lean{VOA.pbw}.
\end{status}

\subsection{The screening and the three generators}

\begin{definition}[Screening, $F$, $H$, $E$]\label{l18:def:screen}
Let $L$ be an even positive-definite lattice, $\alpha\in L$, and $\lambda$ a background charge. Put
$Q=(e^\alpha)_{(0)}=\Res_z Y(e^\alpha,z)$, and
\begin{equation*}
F=e^{-\alpha},\qquad H=QF,\qquad E=QH=Q^2F .
\end{equation*}
For $L=L_p$, $\alpha=\alpha_p$, these are the states \lean{Ftrip p}, \lean{Htrip p}, \lean{Etrip p}, and
$\omega_{\lambda_p}$ is \lean{omegaTrip p}.
\end{definition}
\noindent Files \leanpath{VOA/\allowbreak Examples/\allowbreak TripletScreening.lean} (first block) and
\leanpath{VOA/\allowbreak Examples/\allowbreak TripletWp.lean} (second block).
\begin{leancode}
noncomputable def Qscr (α : L.carrier) : End ℂ (LatModuleGen L) := mode (Yeαvec L α) 0

noncomputable def Fscr (α : L.carrier) : LatModuleGen L := latChargeGen L (-α)

noncomputable def Hscr (α : L.carrier) : LatModuleGen L := Qscr L α (Fscr L α)

noncomputable def Escr (α : L.carrier) : LatModuleGen L := Qscr L α (Hscr L α)
\end{leancode}
\begin{leancode}
noncomputable def omegaTrip (p : ℕ) : LatModuleGen (latP p) :=
  BgA.omegaBg (latP p) (lamP p)

noncomputable def Ftrip (p : ℕ) : LatModuleGen (latP p) :=
  Screen.Fscr (latP p) (alphaP p)

noncomputable def Htrip (p : ℕ) : LatModuleGen (latP p) :=
  Screen.Hscr (latP p) (alphaP p)

noncomputable def Etrip (p : ℕ) : LatModuleGen (latP p) :=
  Screen.Escr (latP p) (alphaP p)
\end{leancode}
\begin{dictionary}
\lean{Yeαvec L α} is the vertex operator of $e^\alpha$ written out as a normal-ordered exponential
(\lean{Qscr_eq_mode_YfullGen} proves it is the zero mode of the reconstructed \lean{YfullGen}, so
$Q$ really is $(e^\alpha)_{(0)}$ in the VOA $V_{L_p}$). In \leanpath{TripletWpC2.lean} the same operator is
re-presented as the generic screening of Chapter~\ref{ch:L04}, Definition~\ref{l04:def:screening}:
\begin{leancode}
noncomputable def Qwp (p : ℕ) (hp : 1 ≤ p) : End ℂ (LatModuleGen (latP p)) :=
  Screening.screening (tripletLatticeVOA p hp) (latChargeGen (latP p) (alphaP p))
\end{leancode}
with \lean{Qwp_eq_Qscr}. (Chapter~\ref{ch:L04} remarks that the triplet screening fits the generic
definition only when $\lambda=0$; that remark predates \lean{tripletLatticeVOA}, and \lean{Qwp}
shows the fit is now exact.)
\end{dictionary}
\begin{warning}
\textbf{Naming: ``short'' versus ``long'' screening.} \leanpath{TripletScreening.lean} and
\leanpath{TripletWp.lean} call $Q=(e^\alpha)_{(0)}$ the ``short screening operator''. In the
literature (Adamovi\'c--Milas; FGST) the two screening momenta are $\alpha$ with
$\langle\alpha,\alpha\rangle=2p$ and $-\alpha/p$ with norm $2/p$; $e^{\alpha}_{(0)}$ is the \emph{long}
screening $Q$, and the short one $\widetilde Q=e^{-\alpha/p}_{(0)}$, which defines $W(p)$ as a kernel,
does not act on $V_{L_p}$ at all. \leanpath{TripletWpC2.lean} uses the correct name. Docstring defect
(documentation only); the mathematics is unaffected. The library never constructs $\widetilde Q$.
\end{warning}
\begin{physics}
$F=e^{-\alpha}$ is the vertex operator of the lowest charge in the triplet; applying the screening
charge once and twice produces the neutral member $H$ (a Schur polynomial in the boson oscillators,
of charge $0$) and the top member $E$ (charge $+\alpha$). Because $Q$ commutes with the deformed
Virasoro algebra, all three have the same conformal weight; the three fields together with $T$ are
Kausch's $W$-algebra generators.
\end{physics}
\begin{status}
Definitions.
\end{status}

\begin{theorem}[Properties of the screening and the generators]\label{l18:thm:screen}
Let $L$ be positive definite, $\alpha\in L$ with $\langle\alpha,\alpha\rangle=2p$, $p\ge1$.
\begin{enumerate}
\item If $\frac12\langle\alpha,\alpha\rangle-\langle\lambda,\alpha\rangle=1$ (i.e.\ $e^\alpha$ has deformed weight $1$),
  then $[L^\lambda_0,Q]=0$.
\item $H=\varepsilon(\alpha,-\alpha)\,S_{2p-1}(\alpha)\otimes e^0$, where $S_N(\alpha)$ is the $N$-th
  \emph{vector Schur polynomial}: the polynomial in the oscillator (creation) variables
  $\alpha(-1),\alpha(-2),\dots$ of the Fock space (Chapter~\ref{ch:L05}, Definition~\ref{l05:def:fock})
  defined by the generating series
  \begin{equation}\label{l18:eq:schur}
  \sum_{N\ge0}S_N(\alpha)\,z^N=\exp\Bigl(\sum_{m\ge1}\frac{\alpha(-m)}{m}\,z^m\Bigr),
  \end{equation}
  so $S_0=1$, $S_1=\alpha(-1)$, $S_2=\frac12\alpha(-1)^2+\frac12\alpha(-2)$,
  $S_3=\frac16\alpha(-1)^3+\frac12\alpha(-1)\alpha(-2)+\frac13\alpha(-3)$; and $S_N(\alpha)\ne0$ for
  $\alpha\ne0$, hence $H\ne0$.
\item $Q^3F=0$.
\item If moreover $\langle\lambda,\alpha\rangle=p-1$, then $F$, $H$, $E$ are $L^\lambda_0$-eigenvectors of eigenvalue $2p-1$.
\end{enumerate}
\end{theorem}
\noindent File \leanpath{VOA/\allowbreak Examples/\allowbreak TripletScreening.lean}.
\begin{leancode}
theorem LBG_zero_comm_Qscr (hPos : L.PosDef) (lam : Fin (rank L) → ℂ) (α : L.carrier)
    (hwt : (1 / 2 : ℂ) * ((L.form α α : ℤ) : ℂ) - bgPair L lam α = 1) :
    LBG L lam 0 * Qscr L α = Qscr L α * LBG L lam 0 := by

theorem Hscr_eq (α : L.carrier) (p : ℕ) (hp : 1 ≤ p) (hαα : L.form α α = 2 * (p : ℤ)) :
    Hscr L α = latticeCocycle L α (-α) • latStateGen L (vSchur L α (2 * p - 1)) 0 := by

theorem vSchur_ne_zero (hPos : L.PosDef) (α : L.carrier) (hα : α ≠ 0) (N : ℕ) :
    vSchur L α N ≠ 0 := by

theorem Hscr_ne_zero (hPos : L.PosDef) (α : L.carrier) (hα : α ≠ 0) (p : ℕ) (hp : 1 ≤ p)
    (hαα : L.form α α = 2 * (p : ℤ)) : Hscr L α ≠ 0 := by

theorem Qscr_cube_Fscr (α : L.carrier) (p : ℕ) (hp : 1 ≤ p) (hαα : L.form α α = 2 * (p : ℤ)) :
    Qscr L α (Qscr L α (Qscr L α (Fscr L α))) = 0 :=

theorem LBG_zero_Fscr (hPos : L.PosDef) (lam : Fin (rank L) → ℂ) (α : L.carrier) (p : ℕ)
    (hαα : L.form α α = 2 * (p : ℤ)) (hlam : bgPair L lam α = (p : ℂ) - 1) :
    LBG L lam 0 (Fscr L α) = (2 * (p : ℂ) - 1) • Fscr L α := by

theorem LBG_zero_Escr (hPos : L.PosDef) (lam : Fin (rank L) → ℂ) (α : L.carrier) (p : ℕ)
    (hp : 1 ≤ p) (hαα : L.form α α = 2 * (p : ℤ)) (hlam : bgPair L lam α = (p : ℂ) - 1) :
    LBG L lam 0 (Escr L α) = (2 * (p : ℂ) - 1) • Escr L α := by
\end{leancode}
\begin{dictionary}
\lean{LBG L lam 0} is $L^\lambda_0$ (the \lean{BgB} copy of \lean{BgA.LBg}); the product \lean{*} of
endomorphisms is composition, so the first statement is the commutation $L^\lambda_0Q=QL^\lambda_0$.
Its hypothesis \lean{hwt} is exactly ``$e^\alpha$ has deformed weight one'' and for $L_p$ it is
\lean{triplet_wt_one} ($p-(p-1)=1$). \lean{vSchur L α N} is the polynomial $S_N(\alpha)$ in the
creation variables defined by $\sum_NS_N(\alpha)z^N=\exp\bigl(\sum_{m\ge1}\frac{\alpha(-m)}{m}z^m\bigr)$
(Chapter~\ref{ch:L05}); \lean{latStateGen L q 0} is $q\otimes e^0$; \lean{latticeCocycle} is the sign
$\varepsilon$. The weight statement for $H$ (\lean{LBG_zero_Hscr}) is analogous. The general-$L$
form is stronger than needed (the triplet uses $L=L_p$); there is no weakening. What the file does
\emph{not} prove is $E\ne0$ (see Theorem~\ref{l18:thm:Ene0} below for $p=2$).
\end{dictionary}
\begin{proof}[Proof on paper]
(1) On $q\otimes e^\gamma$ with $\deg q=d$, $Q$ produces states of charge $\gamma+\alpha$ and Fock degree
$d-1-\langle\alpha,\gamma\rangle$ (the operator $e^\alpha_{(0)}$ contains $z^{\langle\alpha,\gamma\rangle}$ and the coefficient
of $z^{-1}$ is extracted). The deformed weight changes by
$(-1-\langle\alpha,\gamma\rangle)+\bigl(\tfrac12\langle\alpha+\gamma,\alpha+\gamma\rangle-\langle\lambda,\alpha+\gamma\rangle\bigr)-\bigl(\tfrac12\langle\gamma,\gamma\rangle-\langle\lambda,\gamma\rangle\bigr)
=-1+\tfrac12\langle\alpha,\alpha\rangle-\langle\lambda,\alpha\rangle=0$.
(2) By the lattice vertex operator formula (Chapter~\ref{ch:L05}), acting on the Fock vacuum
$e^{-\alpha}$ only the creation half of $Y(e^\alpha,z)$ survives, and
$Y(e^\alpha,z)e^{-\alpha}=\varepsilon(\alpha,-\alpha)z^{\langle\alpha,-\alpha\rangle}\exp\bigl(\sum_{m\ge1}\frac{\alpha(-m)}mz^m\bigr)e^0
=\varepsilon(\alpha,-\alpha)z^{-2p}\exp\bigl(\sum_{m\ge1}\frac{\alpha(-m)}mz^m\bigr)e^0$. By the generating
series \eqref{l18:eq:schur} this is $\varepsilon\sum_NS_N(\alpha)z^{N-2p}e^0$. Since
$Q=e^\alpha_{(0)}$ is the coefficient of $z^{-1}$, we need $N-2p=-1$, i.e.\ $N=2p-1$, and
$H=QF=\varepsilon\,S_{2p-1}(\alpha)e^0$. Non-vanishing of $S_N$: differentiating
with respect to the variable $\alpha_b(-1)$ gives $\partial S_{N+1}=\alpha_bS_N$, so by induction
$S_N\ne0$ once some coordinate $\alpha_b\ne0$. (3) $Q$ preserves weight $2p-1$ and raises charge by $\alpha$;
$Q^3F$ would have charge $2\alpha$, whose lowest weight is $q_p(2)=2p+2>2p-1$, so it is zero (Fock
degree $-3$). (4) $\wt F=\frac12\cdot2p+(p-1)=2p-1$, and $Q$ preserves weight by (1).
\end{proof}
\begin{physics}
(1) is the statement that a screening charge is a conformal invariant: the integral of a weight-one
field commutes with $L_0$ (indeed with all $L_n$; the library proves the full primarity of $E,H,F$
in Theorem~\ref{l18:thm:pairs}). (3) says the screening string starting at $F$ has length three:
the triplet is a spin-one representation of the $\mathfrak{sl}_2$ generated by the screenings.
\end{physics}
\begin{status}
All proved, unconditional. \texttt{\#print axioms} for \lean{LBG_zero_comm_Qscr}, \lean{Hscr_ne_zero},
\lean{Qscr_cube_Fscr}: \lean{[propext, Classical.choice, Quot.sound]}.
\end{status}

\subsection{The triplet algebra as a generated subalgebra}

\begin{definition}[The lattice triplet $W(p)$]\label{l18:def:tripletW}
$W(p)$ is the vertex operator subalgebra of $V_{L_p}$ generated by
$\{\omega_{\lambda_p},E,H,F\}$, with the induced VOA structure.
\end{definition}
\noindent File \leanpath{VOA/\allowbreak Examples/\allowbreak TripletWp.lean}.
\begin{leancode}
noncomputable def tripGens (p : ℕ) : Set (LatModuleGen (latP p)) :=
  {omegaTrip p, Etrip p, Htrip p, Ftrip p}

noncomputable def tripletW (p : ℕ) (hp : 1 ≤ p) : VOASubalgebra (tripletLatticeVOA p hp) :=
  VOASubalgebra.generated (tripletLatticeVOA p hp) (tripGens p)

noncomputable def tripletWVOA (p : ℕ) (hp : 1 ≤ p) :
    VertexOperatorAlgebra (tripletW p hp).carrier :=
  VOASubalgebra.generatedVOA (tripletLatticeVOA p hp) (tripGens p)
\end{leancode}
\begin{dictionary}
The carrier of \lean{tripletWVOA p hp} is the subtype \lean{(tripletW p hp).carrier} of
$V_{L_p}$; its vertex operators, vacuum and conformal vector are the ambient ones
(\lean{tripletWVOA_Y_coe}, \lean{tripletWVOA_vacuum_coe}, \lean{tripletWVOA_conformalVector_coe},
all \lean{rfl}). \lean{VOASubalgebra.generated} is closure under \emph{all} modes and homogeneous
projections (Section~\ref{l18:sec:conv}); this is \emph{weak} generation. It is \emph{not} the
definition \eqref{l18:eq:Wp} as a kernel, and it does not by itself say that $W(p)$ is
\emph{strongly} generated by the four states; the latter is a theorem of Adamovi\'c--Milas
(Proposition~2.3) which the library proves at $p=2$ (Theorem~\ref{l18:thm:mixed}) and carries as
a hypothesis for general $p$.

\textbf{Is this the triplet algebra?} Mathematically yes: by Adamovi\'c--Milas (2008, Prop.~2.3)
the kernel $\ker\widetilde Q$ is strongly generated by $\omega,E,H,F$, so it equals the subalgebra they
generate. Formally, the library proves neither $\ker\widetilde Q\subseteq W(p)$ nor the converse ---
$\widetilde Q$ does not exist in the library. So every theorem about \lean{tripletWVOA} is a theorem
about the subalgebra generated by four explicit states; that this subalgebra is ``the'' triplet
algebra is a citation.

\textbf{Could it be degenerate?} Three degeneracies would make the name misleading: $W(p)=0$,
$W(p)=V_{L_p}$, or $E=0$. The library proves $W(p)\ne0$ for all $p$
(\lean{tripletW_carrier_ne_bot}, via $H\ne0$) and $\omega,H,F\ne0$. The docstring lists $E\ne0$ and
$W(p)\ne V_{L_p}$ as open. At $p=2$ both are settled by a Lean check outside the library
(Theorem~\ref{l18:thm:Ene0} and Remark~\ref{l18:cor:proper}); for $p\ge3$ they remain open in
Lean, although nothing proved depends on them (the relations below are true whether or not $E=0$).
At $p=1$ the generators generate all of $V_{L_1}$ (the docstring of \leanpath{TripletWpC2.lean}
records this), so the $p=1$ instance is not a triplet algebra.
\end{dictionary}
\begin{physics}
$W(p)$ is the chiral algebra of the $(1,p)$ logarithmic minimal model: the symmetry algebra
of the $c_{1,p}$ theory whose Virasoro representations are the logarithmic extensions of the
first row of the Kac table. Its elements are the fields of the Coulomb gas that are ``screened'',
i.e.\ that have well-defined correlators with the screening charge inserted.
\end{physics}
\begin{status}
Definition; \texttt{\#print axioms VOA.LatGen.Wp.tripletWVOA}: \lean{[propext, Classical.choice,
Quot.sound]}.
\end{status}

\begin{theorem}[Central charge; non-degeneracy; consistency with the Abe model]\label{l18:thm:ccmatch}
$W(p)$ has central charge $c_{1,p}$; it contains $\leighteenvac,\omega,E,H,F$; it is nonzero; and at
$p=2$ its central charge equals that of \lean{SympFerm.tripletW2}.
\end{theorem}
\noindent File \leanpath{VOA/\allowbreak Examples/\allowbreak TripletWp.lean}.
\begin{leancode}
@[simp] theorem tripletW_centralCharge (p : ℕ) (hp : 1 ≤ p) :
    (tripletWVOA p hp).centralCharge = 1 - 6 * ((p : ℂ) - 1) ^ 2 / (p : ℂ) := rfl

theorem tripletW_carrier_ne_bot (p : ℕ) (hp : 1 ≤ p) :
    (tripletW p hp).carrier ≠ ⊥ := by

theorem omegaTrip_ne_zero (p : ℕ) (hp : 1 ≤ p) : omegaTrip p ≠ 0 := by

theorem tripletW_centralCharge_eq_tripletW2 (hp : 1 ≤ 2) :
    (tripletWVOA 2 hp).centralCharge = SympFerm.tripletW2.centralCharge := by
\end{leancode}
\begin{dictionary}
The central charge is \lean{rfl}: a subalgebra keeps the ambient conformal vector. The last
statement compares two numbers, $-2=-2$; it is \emph{not} an isomorphism and the library says so
explicitly. \lean{omegaTrip_ne_zero} is proved by observing that $L_0$ acts by $1$ on $e^\alpha\ne0$.
\end{dictionary}
\begin{physics}
Equality of central charges is the weakest possible consistency check between two free-field
realizations of the same CFT (bosonic Coulomb gas and symplectic fermions); the real
identification, due to Kausch, is a bosonization statement that the library does not formalise.
\end{physics}
\begin{status}
Proved, unconditional; standard axioms only.
\end{status}

\begin{theorem}[$E\ne0$ and $W(2)\ne V_{L_2}$; Lean check outside the library]\label{l18:thm:Ene0}
At $p=2$, $E\ne0$, and $e^\alpha\notin W(2)$; hence $W(2)$ is a proper subalgebra of $V_{L_2}$.
\end{theorem}
\begin{proof}[Proof on paper]
The library proves $F_{(3)}E=-48\,\varepsilon(\alpha,-\alpha)\varepsilon(-\alpha,\alpha)\,\omega$
(\lean{Ftrip_three_Etrip_two_explicit}, Section~\ref{l18:sec:c2two}); the right side is nonzero, so
$E\ne0$. The library also proves that $W(2)$ has no nonzero state of weight one
(\lean{eq_zero_of_mem_tripletW_two_and_grading_one}, Theorem~\ref{l18:thm:irrtwo}); $e^\alpha$ is a nonzero
state of weight one, so it is not in $W(2)$.
\end{proof}
\begin{physics}
$E\ne0$ says the top member of the screening triplet is a genuine field: the $\mathfrak{sl}_2$-string
$F\mapsto H\mapsto E$ has full length three, as the name ``triplet'' requires. Properness says the chiral algebra is
smaller than the free-field theory: the screening current $e^\alpha$ itself is not a symmetry generator.
\end{physics}
\begin{status}
Checked in Lean against the library (the check is not part of this release); not a library theorem. The three statements depend on
\lean{[propext, Classical.choice, Quot.sound]}. These close, at $p=2$, two items listed as ``not
proved'' in the docstrings of \leanpath{TripletWp.lean} and \leanpath{TripletWpC2.lean}
(documentation defect).
\end{status}
\begin{remark}\label{l18:cor:proper}
So \lean{tripletW 2 hp} is neither $0$ nor $V_{L_2}$, and its three triplet generators are nonzero: none of the
degeneracies that would make the name misleading occurs at $p=2$.
\end{remark}

%======================================================================
\section{The lattice model II: \texorpdfstring{$C_2$}{C2}-cofiniteness for general \texorpdfstring{$p$}{p}}\label{l18:sec:c2general}
%======================================================================

$C_2$-cofiniteness is the finiteness condition that survives in the logarithmic world: it gives
finitely many irreducible modules, finite-dimensional Zhu algebras and (Miyamoto) modular invariance
of pseudo-traces, but it does not give semisimplicity. Recall (Chapter~\ref{ch:L12},
Section~\ref{l12:sec:c2}) that $R_V=V/C_2(V)$ is a commutative Poisson algebra with product
$\bar u\cdot\bar v=\overline{u_{(-1)}v}$, and that \emph{if $V$ is strongly generated by $S$ then $R_V$ is
generated as a commutative algebra by the classes $\bar s$, $s\in S$} (every word containing a mode
$s_{(-m)}$ with $m\ge2$ lies in $C_2$). A commutative algebra generated by finitely many nilpotent
elements is finite-dimensional. So:

\begin{center}
\emph{strong generation by $S$ finite}\ +\ \emph{each $\bar s$ nilpotent in $R_V$}\ $\Longrightarrow$\ $C_2$-cofinite.
\end{center}

This is the route of Abe (2007) and Adamovi\'c--Milas (2008, Theorem~3.1), and of the library.

\subsection{The engine}

\begin{theorem}[Nilpotent generators give $C_2$; dimension bound]\label{l18:thm:engine}
Let $V$ be a VOA and $g:\iota\to V$ a finite family. If the classes $\bar g_i$ generate $R_V$ as a
commutative $\C$-algebra and each $\bar g_i$ is nilpotent, then $V$ is $C_2$-cofinite; if
$\bar g_i^{m_i}=0$ then $\dim R_V\le\prod_im_i$.
\end{theorem}
\noindent File \leanpath{VOA/\allowbreak Examples/\allowbreak TripletW2Nilpotent.lean}.
\begin{leancode}
theorem c2Cofinite_of_adjoin_of_isNilpotent (W : VertexOperatorAlgebra 𝒱) {ι : Type*}
    [Finite ι] (g : ι → 𝒱)
    (htop : Algebra.adjoin ℂ (Set.range (fun i => mk₂ W (g i))) = ⊤)
    (hnil : ∀ i, IsNilpotent (mk₂ W (g i))) : C2Cofinite W :=

theorem finrank_a₂_le_of_adjoin (W : VertexOperatorAlgebra 𝒱) {ι : Type*}
    [Fintype ι] (g : ι → 𝒱) (m : ι → ℕ)
    (htop : Algebra.adjoin ℂ (Set.range (fun i => mk₂ W (g i))) = ⊤)
    (hm : ∀ i, (mk₂ W (g i)) ^ m i = 0) :
    Module.finrank ℂ (A₂ W) ≤ ∏ i, m i :=
\end{leancode}
\begin{dictionary}
\lean{Algebra.adjoin ℂ s = ⊤} is Mathlib's ``the subalgebra generated by $s$ is everything'';
the algebra structure on \lean{A₂ W} is the instance \lean{a₂Algebra} built in the same file from the
descended product (Chapter~\ref{ch:L12}, Theorem~\ref{l12:thm:poisson}). \lean{IsNilpotent x} is
$\exists n,\ x^n=0$. The bound is the number of monomials $\prod_i\bar g_i^{k_i}$, $k_i<m_i$, which span. The
statement is exactly the textbook lemma; nothing is hidden. A companion,
\lean{adjoin_range_mk₂_of_iterProdSpan}, derives \lean{htop} from strong generation.
\end{dictionary}
\begin{physics}
$R_V$ is the ring of functions on the ``associated variety'' of the chiral algebra, the classical
limit in which only the leading singular terms of OPEs survive. Nilpotency of all generators means
the associated variety is a point: there are no ``flat directions'' in the space of classical field
configurations, which is the algebraic shadow of having finitely many primaries.
\end{physics}
\begin{status}
Proved, unconditional; standard axioms.
\end{status}

\subsection{The relations in \texorpdfstring{$R_{W(p)}$}{R\_W(p)}}

\begin{theorem}[Adamovi\'c--Milas relations]\label{l18:thm:AMrel}
For every $p\ge1$ and every $n\ge-2p$, in $V_{L_p}$:
\begin{equation*}
F_{(n)}F=0,\qquad E_{(n)}E=0,\qquad E_{(n)}F+F_{(n)}E+2H_{(n)}H=0,
\end{equation*}
and $Q$ is a derivation of all products: $Q(u_{(n)}v)=(Qu)_{(n)}v+u_{(n)}(Qv)$. Consequently, in
$R_{W(p)}$,
\begin{equation*}
\bar E^2=0,\qquad \bar F^2=0,\qquad \bar H^2=-\bar E\bar F,\qquad \bar H^4=0 .
\end{equation*}
\end{theorem}
\noindent File \leanpath{VOA/\allowbreak Examples/\allowbreak TripletWpC2.lean}.
\begin{leancode}
theorem smul_Ftrip_Ftrip (p : ℕ) (hp : 1 ≤ p) {n : ℤ} (hn : -(2 * (p : ℤ)) ≤ n) :
    (tripletLatticeVOA p hp).smul (Ftrip p) n (Ftrip p) = 0 := by

theorem Qwp_leibniz (p : ℕ) (hp : 1 ≤ p) (u : LatModuleGen (latP p)) (n : ℤ)
    (v : LatModuleGen (latP p)) :
    Qwp p hp ((tripletLatticeVOA p hp).smul u n v)
      = (tripletLatticeVOA p hp).smul (Qwp p hp u) n v
        + (tripletLatticeVOA p hp).smul u n (Qwp p hp v) :=

theorem EF_add_FE_add_two_HH (p : ℕ) (hp : 1 ≤ p) {n : ℤ} (hn : -(2 * (p : ℤ)) ≤ n) :
    (tripletLatticeVOA p hp).smul (Etrip p) n (Ftrip p)
      + (tripletLatticeVOA p hp).smul (Ftrip p) n (Etrip p)
      + (2 : ℂ) • (tripletLatticeVOA p hp).smul (Htrip p) n (Htrip p) = 0 := by

theorem smul_Etrip_Etrip (p : ℕ) (hp : 1 ≤ p) {n : ℤ} (hn : -(2 * (p : ℤ)) ≤ n) :
    (tripletLatticeVOA p hp).smul (Etrip p) n (Etrip p) = 0 := by
\end{leancode}
\begin{leancode}
noncomputable def genWp (p : ℕ) (hp : 1 ≤ p) : Fin 4 → (tripletW p hp).carrier
  | 0 => genOmega p hp
  | 1 => genE p hp
  | 2 => genH p hp
  | 3 => genF p hp

theorem mk₂_genE_sq (p : ℕ) (hp : 1 ≤ p) :
    (mk₂ (tripletWVOA p hp) (genE p hp)) ^ 2 = 0 := by

theorem mk₂_genF_sq (p : ℕ) (hp : 1 ≤ p) :
    (mk₂ (tripletWVOA p hp) (genF p hp)) ^ 2 = 0 := by

theorem mk₂_genH_sq (p : ℕ) (hp : 1 ≤ p) :
    (mk₂ (tripletWVOA p hp) (genH p hp)) ^ 2
      = -(mk₂ (tripletWVOA p hp) (genE p hp) * mk₂ (tripletWVOA p hp) (genF p hp)) := by

theorem mk₂_genH_pow_four (p : ℕ) (hp : 1 ≤ p) :
    (mk₂ (tripletWVOA p hp) (genH p hp)) ^ 4 = 0 := by
\end{leancode}
\begin{dictionary}
\lean{genOmega p hp} $=\langle\omega_{\lambda_p},\ldots\rangle$, \lean{genE}, \lean{genH}, \lean{genF} are the four
generators as elements of the subtype \lean{(tripletW p hp).carrier}; \lean{smul_coe} says products in
$W(p)$ are the ambient ones, so the $V_{L_p}$-relations transport. The relations hold in $V_{L_p}$
for $n\ge-2p$, in particular at $n=-1$, which is what the $R_{W(p)}$ relations use. The $R$-relations
are formed in \lean{A₂ (tripletWVOA p hp)}, i.e.\ with the subalgebra's \emph{own} $C_2$-space; this
matters, since $C_2(W)\subseteq C_2(V)\cap W$ can be strict (Section~\ref{l18:sec:abe}). If $E$ were
zero the relations would degenerate to $\bar H^2=0$, which still implies $\bar H^4=0$; so the
conclusions are not hostage to the open question $E\ne0$ for $p\ge3$.
\end{dictionary}
\begin{proof}[Proof on paper]
$Y(e^\beta,z)e^\gamma=\varepsilon\,z^{\langle\beta,\gamma\rangle}(\cdots)e^{\beta+\gamma}$ with $(\cdots)$ a power series in $z$;
for $\beta=\gamma=-\alpha$, $\langle\beta,\gamma\rangle=2p$, so $F_{(n)}F=0$ for $n\ge-2p$
(\lean{smul_latChargeGen_eq_zero}). The derivation property of a zero mode is the $\ell=0$ slice of
the Borcherds identity (Chapter~\ref{ch:L04}, Theorem~\ref{l04:thm:scrder}; here
\lean{Core.ScreeningKernel.screening_leibniz} needs no hypothesis on $Q$). Applying $Q$, $Q^2$, $Q^3$,
$Q^4$ to $F_{(n)}F=0$ and using $QF=H$, $QH=E$, $QE=0$ gives successively $H_{(n)}F+F_{(n)}H=0$,
$E_{(n)}F+2H_{(n)}H+F_{(n)}E=0$, $E_{(n)}H+H_{(n)}E=0$, and $E_{(n)}E=0$. In $R_{W(p)}$ the product is
commutative, so the $n=-1$ relations read $\bar F^2=\bar E^2=0$ and $2\bar E\bar F+2\bar H^2=0$. Then
$\bar H^4=(\bar E\bar F)^2=\bar E^2\bar F^2=0$.
\end{proof}
\begin{physics}
These are statements about the singular parts of the OPEs $F(z)F(w)$ etc.: the like-charge OPE
$e^{-\alpha}(z)e^{-\alpha}(w)\sim(z-w)^{2p}$ is so regular that even its normal-ordered product
vanishes, and the screening symmetry rotates this into the full set of triplet OPE relations. In the
classical (associated-variety) limit the triplet fields become nilpotent coordinates.
\end{physics}
\begin{status}
Proved, unconditional; \texttt{\#print axioms VOA.LatGen.WpC2.mk₂\_genH\_pow\_four}:
\lean{[propext, Classical.choice, Quot.sound]}.
\end{status}

\subsection{The two residual hypotheses, and the theorem}

\begin{definition}[Generation and $\bar\omega$-nilpotency hypotheses]\label{l18:def:residuals}
\begin{itemize}
\item \lean{StrongGenWp p hp}: $W(p)$ is strongly generated by $\omega,E,H,F$.
\item \lean{AdjoinGenWp p hp}: $R_{W(p)}$ is generated as a commutative algebra by
  $\bar\omega,\bar E,\bar H,\bar F$ (implied by the previous one).
\item \lean{OmegaNilpotentWp p hp}: $\bar\omega$ is nilpotent in $R_{W(p)}$.
\item \lean{OmegaPowMem p hp m}: $L_{-2}^m\leighteenvac\in C_2(W(p))$, i.e.\ $\bar\omega^m=0$.
\end{itemize}
\end{definition}
\noindent File \leanpath{VOA/\allowbreak Examples/\allowbreak TripletWpC2.lean}.
\begin{leancode}
def StrongGenWp (p : ℕ) (hp : 1 ≤ p) : Prop :=
  iterProdSpan (tripletWVOA p hp) (Set.range (genWp p hp)) = ⊤

def AdjoinGenWp (p : ℕ) (hp : 1 ≤ p) : Prop :=
  Algebra.adjoin ℂ (Set.range (fun i => mk₂ (tripletWVOA p hp) (genWp p hp i))) = ⊤

theorem adjoinGenWp_of_strongGenWp (p : ℕ) (hp : 1 ≤ p) (h : StrongGenWp p hp) :
    AdjoinGenWp p hp :=

def OmegaNilpotentWp (p : ℕ) (hp : 1 ≤ p) : Prop :=
  IsNilpotent (mk₂ (tripletWVOA p hp) (genOmega p hp))

def OmegaPowMem (p : ℕ) (hp : 1 ≤ p) (m : ℕ) : Prop :=
  TripletNilp.pMulPow (tripletWVOA p hp) (genOmega p hp) m ∈ O₂ (tripletWVOA p hp)
\end{leancode}
\begin{dictionary}
\lean{TripletNilp.pMulPow W u m} is the $m$-fold normal-ordered power $u_{(-1)}(u_{(-1)}(\cdots u_{(-1)}\leighteenvac))$,
and for $u=\omega$ it is $L_{-2}^m\leighteenvac$ (\lean{pMulPow_genOmega_succ}); \lean{mk₂_genOmega_pow_eq_zero_iff}
proves \lean{OmegaPowMem p hp m} $\iff\bar\omega^m=0$. \lean{OmegaNilpotentWp} is deliberately exponent-free
so that no guessed exponent can make a theorem vacuous. All four are genuine \lean{Prop}s about the
\emph{subalgebra's} $C_2$ quotient. They are hypotheses, never axioms.
\end{dictionary}
\begin{physics}
\lean{StrongGenWp} is the statement that the chiral algebra is a $W$-algebra with exactly four
generating fields $T,W^+,W^0,W^-$ (every field is a normally ordered polynomial in them and their
derivatives). \lean{OmegaPowMem} says that a power of the stress tensor, $:T^m:$, is ``null modulo
derivatives and composite fields'' in the classical limit --- the existence of a null vector of a
special shape.
\end{physics}
\begin{status}
Definitions. A Lean check (not part of this release) proves
\lean{AdjoinGenWp 2 hp} and \lean{OmegaPowMem 2 hp 5}: the hypotheses are \emph{satisfiable}, so the
conditional theorems below are not vacuous.
\end{status}

\begin{theorem}[$C_2$-cofiniteness of $W(p)$, conditional]\label{l18:thm:c2general}
For every $p\ge1$ and $m\in\N$: if $R_{W(p)}$ is generated by $\bar\omega,\bar E,\bar H,\bar F$ and
$L_{-2}^m\leighteenvac\in C_2(W(p))$, then $W(p)$ is $C_2$-cofinite and $\dim R_{W(p)}\le16m$.
\end{theorem}
\noindent File \leanpath{VOA/\allowbreak Examples/\allowbreak TripletWpC2.lean}.
\begin{leancode}
theorem c2Cofinite_tripletWVOA (p : ℕ) (hp : 1 ≤ p) (hgen : AdjoinGenWp p hp)
    (homega : OmegaNilpotentWp p hp) : C2Cofinite (tripletWVOA p hp) := by

theorem c2Cofinite_tripletWVOA_of_omegaPowMem (p : ℕ) (hp : 1 ≤ p) (m : ℕ)
    (hgen : AdjoinGenWp p hp) (homega : OmegaPowMem p hp m) :
    C2Cofinite (tripletWVOA p hp)
      ∧ Module.finrank ℂ (A₂ (tripletWVOA p hp)) ≤ 16 * m :=
\end{leancode}
\begin{dictionary}
Two hypotheses, both named: generation (\lean{hgen}) and $\bar\omega$-nilpotency (\lean{homega}). The
\textsc{Library Map} summary ``general $p$ conditional on \lean{OmegaPowMem}'' omits the first; it is
also needed. $16m=m\cdot2\cdot4\cdot2$ is the engine's bound with exponents $(m,2,4,2)$ for
$(\bar\omega,\bar E,\bar H,\bar F)$. With Adamovi\'c--Milas' $m=3p-1$ this is $48p-16$; the true value
of $\dim R_{W(p)}$ is expected by the library's docstring to be $6p-1$ (Abe proves $\le11$ at $p=2$), so the bound is crude but finite,
which is all $C_2$-cofiniteness asks.
\end{dictionary}
\begin{proof}[Proof on paper]
Theorem~\ref{l18:thm:engine} with $g=(\omega,E,H,F)$: generation is \lean{hgen}, nilpotency of
$\bar E,\bar F,\bar H$ is Theorem~\ref{l18:thm:AMrel}, and of $\bar\omega$ is \lean{homega}.
\end{proof}
\begin{physics}
This is ``the $(1,p)$ triplet theory is a finite (lisse) logarithmic CFT'', modulo the two
statements that its $W$-algebra closes on four generators and that $:T^{3p-1}:$ is classically null.
\end{physics}
\begin{status}
Proved, \textbf{conditional on \lean{AdjoinGenWp p hp} and \lean{OmegaPowMem p hp m}} (or
\lean{OmegaNilpotentWp}); \texttt{\#print axioms}: standard three. No \lean{VOA.pbw}. The hypotheses are discharged only at
$p=2$ (in a Lean check outside the library, not part of this release); for $p\ge3$ this is a conditional reduction, not ``$W(p)$ is $C_2$-cofinite''.
\end{status}

\begin{warning}
\textbf{The printed argument for $\bar\omega$-nilpotency in Adamovi\'c--Milas (2008), Theorem~3.1, is
a non-sequitur} (recorded in the docstring of \leanpath{TripletWpC2.lean}; we have checked the
reasoning). The paper says: by [Adamovi\'c 2003], $\bar H^2=C_p\bar\omega^{2p-1}$ with $C_p\ne0$; since
$\bar H^4=0$, $\bar\omega^{4p-2}=0$. But the cited relation holds in \emph{Zhu's algebra} $A(V)=V/O(V)$
--- and there it is $[H]*[H]=C_pP([\omega])$ for a non-homogeneous polynomial $P$ of degree $2p-1$,
not the monomial. The only comparison map between the two quotients is the epimorphism
$R_V\twoheadrightarrow\mathrm{gr}\,A(V)$, which goes the wrong way: a relation in $A(V)$ pulls back to
nothing in $R_V$. \emph{The repair} is inside the same paper: Corollary~4.6 proves
$\bar\omega^{3p-1}=0$ from an identity \emph{in $V_{L_p}$ itself},
\begin{equation}\label{l18:eq:AMid}
Q^2e^{-2\alpha}=A\cdot H_{(-2p-1)}H+\bigl(B_p\,L_{-2}^{3p-1}+w\bigr)\leighteenvac,\qquad B_p\ne0,
\end{equation}
with $w$ a combination of Virasoro monomials each containing some $L_{-n}$, $n\ge3$; every term other
than $B_pL_{-2}^{3p-1}\leighteenvac$ is visibly in $C_2$, and the left side is too, since
$Q^2e^{-2\alpha}=\varepsilon^{-1}Q^2(F_{(-2p-1)}F)=\varepsilon^{-1}(E_{(-2p-1)}F+2H_{(-2p-1)}H+F_{(-2p-1)}E)$.
Corollary~4.6 does not depend on Theorem~3.1, so there is no circularity. The library formalises
exactly this repair (next theorem). (This is a finding about the literature, not a Lean defect.)
\end{warning}

\begin{definition}[The Adamovi\'c--Milas identity; deep Virasoro descendants]\label{l18:def:AMid}
\lean{VirDeep p hp} is the span of the states $L_{-n}y$, $y\in W(p)$, $n\ge1$, $n\ne2$.
\lean{AMOmegaIdentity p hp m} asserts \eqref{l18:eq:AMid} with exponent $m$: there are $A,B\in\C$, $B\ne0$,
and $w\in$ \lean{VirDeep p hp} with $Q^2e^{-2\alpha}=A\cdot H_{(-2p-1)}H+B\cdot L_{-2}^m\leighteenvac+w$.
\end{definition}
\noindent File \leanpath{VOA/\allowbreak Examples/\allowbreak TripletWpAmbient.lean}.
\begin{leancode}
noncomputable def VirDeep (p : ℕ) (hp : 1 ≤ p) : Submodule ℂ (LatModuleGen (latP p)) :=
  Submodule.span ℂ {x | ∃ n : ℤ, 1 ≤ n ∧ n ≠ 2 ∧ ∃ y ∈ (tripletW p hp).carrier,
    x = (tripletLatticeVOA p hp).L (-n) y}

noncomputable def omegaPowAmb (p : ℕ) (hp : 1 ≤ p) : ℕ → LatModuleGen (latP p)
  | 0 => latVacuumGen (latP p)
  | (m + 1) => (tripletLatticeVOA p hp).smul (omegaTrip p) (-1) (omegaPowAmb p hp m)

def AMOmegaIdentity (p : ℕ) (hp : 1 ≤ p) (m : ℕ) : Prop :=
  ∃ A B : ℂ, B ≠ 0 ∧ ∃ w ∈ VirDeep p hp,
    Qwp p hp (Qwp p hp (latChargeGen (latP p) (-alphaP p + -alphaP p)))
      = A • (tripletLatticeVOA p hp).smul (Htrip p) (-(2 * (p : ℤ)) - 1) (Htrip p)
        + B • omegaPowAmb p hp m + w
\end{leancode}
\begin{leancode}
theorem smul_Htrip_Htrip_bottom_mem_O₂amb (p : ℕ) (hp : 1 ≤ p) :
    (tripletLatticeVOA p hp).smul (Htrip p) (-(2 * (p : ℤ)) - 1) (Htrip p) ∈ O₂amb p hp :=

theorem Qwp_sq_latChargeGen_mem_O₂amb (p : ℕ) (hp : 1 ≤ p) :
    Qwp p hp (Qwp p hp (latChargeGen (latP p) (-alphaP p + -alphaP p))) ∈ O₂amb p hp := by

theorem omegaPowMem_of_amOmegaIdentity (p : ℕ) (hp : 1 ≤ p) (m : ℕ) (h : AMOmegaIdentity p hp m) :
    OmegaPowMem p hp m := by

theorem c2Cofinite_tripletWVOA_of_amCor44 (p : ℕ) (hp : 1 ≤ p) (m : ℕ)
    (hgen : AdjoinGenWp p hp) (h : AMOmegaIdentity p hp m) :
    C2Cofinite (tripletWVOA p hp)
      ∧ Module.finrank ℂ (A₂ (tripletWVOA p hp)) ≤ 16 * m :=
\end{leancode}
\begin{dictionary}
\lean{O₂amb p hp} is the image of $C_2(W(p))$ in $V_{L_p}$ under the inclusion; \lean{mem_O₂amb_iff}
proves that ambient membership of a state of $W(p)$ is equivalent to intrinsic membership, so nothing
is lost by computing in $V_{L_p}$. \lean{omegaPowAmb p hp m} is $L_{-2}^m\leighteenvac$ (since
$\omega_{(-1)}=L_{-2}$). The latticed vector \lean{-alphaP p + -alphaP p} is $-2\alpha$. The library's
\lean{VirDeep} is larger than Adamovi\'c--Milas' $w$-space (it allows $L_{-n}$ applied to any state of
$W(p)$, not only Virasoro monomials on the vacuum), so \lean{AMOmegaIdentity} is \emph{weaker} (easier to
satisfy) than the paper's identity; that is harmless because every element of \lean{VirDeep} is in
$C_2$ (\lean{virDeep_le_O₂amb}: $L_{-1}=T$ and $L_{-n}=\omega_{(-n+1)}$ with $-n+1\le-2$ for $n\ge3$),
which is all the argument needs. The two membership theorems are \emph{unconditional}, for every $p$:
they are the two ``visibly in $C_2$'' facts of the warning. The identity cannot hold trivially: the
grading forces $m=3p-1$ (\lean{omegaPowAmb_mem_grading}, everything lives in weight $6p-2$).
\end{dictionary}
\begin{proof}[Proof on paper]
Given \eqref{l18:eq:AMid}, $B\,L_{-2}^m\leighteenvac=Q^2e^{-2\alpha}-A\,H_{(-2p-1)}H-w\in C_2(W(p))$ by
the two membership theorems and \lean{virDeep_le_O₂amb}; divide by $B$.
\end{proof}
\begin{physics}
\eqref{l18:eq:AMid} is a null-vector statement in the Coulomb gas: the doubly screened vertex operator
$Q^2e^{-2\alpha}$ is a descendant of the vacuum module plus a composite of $H$ with itself, and its
purely-$T$ part has a nonzero coefficient. It is the free-field reason for the polynomial relation
satisfied by $T$ in the triplet theory.
\end{physics}
\begin{status}
Proved; the last theorem is conditional on \lean{AdjoinGenWp} and \lean{AMOmegaIdentity}; the two
membership theorems are unconditional. Standard axioms.
\end{status}

\subsection{Reducing generation to a finite OPE table}

\begin{theorem}[Generation from the generator OPE table]\label{l18:thm:pairs}
Let $T_p$ be the span of the strictly-negative-mode words in $\omega,E,H,F$. If $g_{(j)}h\in T_p$ for
all generators $g,h$ and $0\le j<4p+2$, then $W(p)$ is strongly generated by $\omega,E,H,F$. The
$\omega$-row of the table holds for every $p$, because $E,H,F$ are Virasoro primary. The remaining
mixed pairs reduce to the two ``lowering rows'' $F_{(j)}H$ and $F_{(j)}E$.
\end{theorem}
\noindent File \leanpath{VOA/\allowbreak Examples/\allowbreak TripletWpPairClosure.lean}.
\begin{leancode}
def GenPairClosedWp (p : ℕ) (hp : 1 ≤ p) : Prop :=
  ∀ g ∈ tripGens p, ∀ h ∈ tripGens p, ∀ j : ℤ, 0 ≤ j →
    (tripletLatticeVOA p hp).smul g j h ∈ iterProdSpan (tripletLatticeVOA p hp) (tripGens p)

theorem strongGenWp_of_genPairClosed (p : ℕ) (hp : 1 ≤ p) (h : GenPairClosedWp p hp) :
    StrongGenWp p hp :=

def LoweringPairsWp (p : ℕ) (hp : 1 ≤ p) : Prop :=
  (∀ j : ℤ, 0 ≤ j →
      (tripletLatticeVOA p hp).smul (Ftrip p) j (Htrip p)
        ∈ iterProdSpan (tripletLatticeVOA p hp) (tripGens p))
    ∧
  (∀ j : ℤ, 0 ≤ j →
      (tripletLatticeVOA p hp).smul (Ftrip p) j (Etrip p)
        ∈ iterProdSpan (tripletLatticeVOA p hp) (tripGens p))

theorem screenGens_primary (p : ℕ) (hp : 1 ≤ p) :
    ∀ g ∈ screenGens p, ∀ n : ℤ, 1 ≤ n → (tripletLatticeVOA p hp).L n g = 0 := by

theorem c2Cofinite_tripletWVOA_of_lowering_AM (p : ℕ) (hp : 1 ≤ p)
    (hlower : LoweringPairsWp p hp) (hAM : AMOmegaIdentity p hp (3 * p - 1)) :
    C2Cofinite (tripletWVOA p hp)
      ∧ Module.finrank ℂ (A₂ (tripletWVOA p hp)) ≤ 16 * (3 * p - 1) :=
\end{leancode}
\begin{dictionary}
\lean{GenPairClosedWp} is stated in the ambient $V_{L_p}$ with \lean{iterProdSpan} of
\lean{tripletLatticeVOA}; \lean{strongGenWp_iff} (\leanpath{TripletWpAmbient.lean}) shows the ambient span
equals $W(p)$ exactly when \lean{StrongGenWp} holds, so this is faithful. The reduction from ``closed
under all products of all elements'' to ``closed under non-negative products of generators'' is the
general \lean{iterProdSpan_smul_closed} (Chapter~\ref{ch:L12}, Theorem~\ref{l12:thm:fingen} and its companion),
which needs the generators homogeneous of positive weight; the bound $j<4p+2$ comes from
$g_{(j)}h=0$ for $j\ge\wt g+\wt h$. Both \lean{GenPairClosedWp} and \lean{LoweringPairsWp} quantify over
\emph{all} $j\ge0$ rather than $0\le j<4p+2$; the two forms are equivalent, because for larger $j$ the
product has negative weight, hence is $0$, which lies in every span. \lean{screenGens p} $=\{E,H,F\}$. \lean{screenGens_primary} is
$L_nE=L_nH=L_nF=0$ for $n\ge1$. The final theorem therefore has exactly two named hypotheses for general
$p$: the two lowering rows (a finite set of Fock-space computations, $0\le j<4p+2$) and the
Adamovi\'c--Milas identity.
\end{dictionary}
\begin{proof}[Proof on paper]
Strong generation: commuting a non-negative mode $g_{(j)}$ past a negative one,
$[g_{(j)},h_{(-m)}]=\sum_i\binom ji(g_{(i)}h)_{(j-m-i)}$, produces modes of the products $g_{(i)}h$; if
these lie in the span of words, an induction on weight shows the span is closed. Primarity: charged
vacua $e^{\beta}$ are primary for the deformed stress tensor (a Fock-vacuum computation), so $F$ and
$e^\alpha$ are; skew-symmetry gives $Q\omega=0$, so $Q$ commutes with all $L_n$, and $L_nH=L_nQF=QL_nF=0$,
likewise for $E$. Mixed pairs: applying the derivation $Q$ to $F_{(j)}H$ and $F_{(j)}E$ and using the
exact identities of Theorem~\ref{l18:thm:AMrel} recovers $H_{(j)}H$, $H_{(j)}E$, $E_{(j)}F$, \dots
\end{proof}
\begin{physics}
The table $\{g_{(j)}h\}_{j\ge0}$ is the list of singular OPE coefficients among the generating fields.
Closure of the $W$-algebra means that every singular term in $W^a(z)W^b(w)$ is a normally ordered
polynomial in $T,W^a$ and derivatives --- the classic Kausch computation, which here becomes a finite
check in the Fock space.
\end{physics}
\begin{status}
Proved; the last theorem conditional on \lean{LoweringPairsWp p hp} and
\lean{AMOmegaIdentity p hp (3 * p - 1)}. \lean{screenGens_primary} is unconditional. Standard axioms.
\end{status}

%======================================================================
\section{The lattice model III: the unconditional theorems at \texorpdfstring{$p=2$}{p=2}}\label{l18:sec:c2two}
%======================================================================

At $p=2$ both residual families are discharged by explicit computation in the Fock space of the rank-one
lattice $L_2$ ($\langle\alpha,\alpha\rangle=4$). Every Fock state of charge $0$ and given degree is a polynomial
in finitely many oscillator variables; the computations are identities between explicit polynomials with
rational coefficients, checked by Lean's \lean{ring}/\lean{norm_num} machinery, not by \lean{native_decide}.

\subsection{The three mixed coordinates}\label{l18:sec:mixed}
By Theorem~\ref{l18:thm:pairs}, strong generation reduces to showing that the products $F_{(j)}H$
and $F_{(j)}E$, $j\ge0$, are polynomials in the generators. These are ``mixed'' products (of two
\emph{different} generators, one of charge $-\alpha$ and one of charge $0$ or $+\alpha$). At $p=2$ all
but three of them vanish or are settled by short degree-zero arguments (see the dictionary after
Theorem~\ref{l18:thm:mixed}); the three that need a real computation are $F_{(0)}H$, $F_{(1)}E$ and
$F_{(3)}E$. Each lies in a finite-dimensional weight-and-charge space of the Fock space, and
``computing it'' means finding its \emph{coordinates} in a basis of that space and recognising them as
the coordinates of a polynomial in the generators. The library calls these three computations the
\emph{three mixed coordinates}; Theorem~\ref{l18:thm:mixed} states their results.
\begin{theorem}[Mixed products at $p=2$]\label{l18:thm:mixed}
With $c=\varepsilon(\alpha,-\alpha)$, $d=\varepsilon(-\alpha,\alpha)$, in $V_{L_2}$:
\begin{equation*}
F_{(0)}H=\tfrac{48c}{5}\,L_{-2}F+\tfrac{4c}{5}\,T^2F,\qquad
F_{(1)}E=-64cd\,L_{-2}\omega+12cd\,T^2\omega+10d\,TH,\qquad
F_{(3)}E=-48cd\,\omega .
\end{equation*}
Consequently all mixed generator pairs close, and $W(2)$ is strongly generated by $\omega,E,H,F$.
\end{theorem}
\noindent Files \leanpath{VOA/\allowbreak Examples/\allowbreak TripletWpTwoMixed.lean},
\leanpath{TripletWpTwoMixedClose.lean}, \leanpath{TripletWpTwoIrrational.lean}.
\begin{leancode}
noncomputable abbrev tripletTwoC : ℂ :=
  latticeCocycle (latP 2) (alphaP 2) (-alphaP 2)

noncomputable abbrev tripletTwoD : ℂ :=
  latticeCocycle (latP 2) (-alphaP 2) (alphaP 2)

def ThreeExplicitMixedIdentitiesWp (hp : 1 ≤ 2) : Prop :=
  let W := tripletLatticeVOA 2 hp
  W.smul (Ftrip 2) 0 (Htrip 2) =
      ((48 : ℂ) / 5 * tripletTwoC) • W.smul (omegaTrip 2) (-1) (Ftrip 2) +
        ((4 : ℂ) / 5 * tripletTwoC) • (W.toPreVertexAlgebra.T ^ 2) (Ftrip 2) ∧
  W.smul (Ftrip 2) 1 (Etrip 2) =
      (-64 * (tripletTwoC * tripletTwoD)) • W.smul (omegaTrip 2) (-1) (omegaTrip 2) +
        (12 * (tripletTwoC * tripletTwoD)) • (W.toPreVertexAlgebra.T ^ 2) (omegaTrip 2) +
          (10 * tripletTwoD) • W.toPreVertexAlgebra.T (Htrip 2) ∧
  W.smul (Ftrip 2) 3 (Etrip 2) =
      (-48 * (tripletTwoC * tripletTwoD)) • omegaTrip 2
\end{leancode}
\noindent File \leanpath{VOA/Examples/TripletWpTwoMixedAssembly.lean}.
\begin{leancode}
theorem Ftrip_three_Etrip_two_explicit (hp : 1 ≤ 2) :
    (tripletLatticeVOA 2 hp).smul (Ftrip 2) 3 (Etrip 2) =
      (-48 * (tripletTwoC * tripletTwoD)) • omegaTrip 2 := by

theorem threeCoordinateIdentitiesWp_two (hp : 1 ≤ 2) :
    ThreeCoordinateIdentitiesWp hp :=

theorem mixedPairsWp_two (hp : 1 ≤ 2) : MixedPairsWp 2 hp :=

theorem strongGenWp_two (hp : 1 ≤ 2) : StrongGenWp 2 hp :=
\end{leancode}
\begin{dictionary}
\lean{W.toPreVertexAlgebra.T} is the translation $T=L_{-1}$, so \lean{T ^ 2} is $T^2$, and
\lean{W.smul (omegaTrip 2) (-1) x} is $\omega_{(-1)}x=L_{-2}x$. \lean{ThreeCoordinateIdentitiesWp} (in
\leanpath{TripletWpPairClosure.lean}) is the same three equations with existential coefficients; only
membership of the right sides in the word span is used downstream. The three products not in the list
($F_{(2)}H=aF$, $F_{(5)}E=a\leighteenvac$, and the vanishing ones) are handled by degree-zero arguments.
\lean{strongGenWp_two} is the composite: mixed pairs (above) plus primarity tails plus the reduction of
Theorem~\ref{l18:thm:pairs}. The proved $F_{(3)}E\ne0$ is what gives $E\ne0$ (Theorem~\ref{l18:thm:Ene0}).
\end{dictionary}
\begin{proof}[Proof on paper]
Each left side is a single lattice product: e.g.\ $F_{(3)}E$ is the coefficient of $z^{-4}$ in
$Y(e^{-\alpha},z)$ applied to $E=\varepsilon\,(\mathrm{EE}(\alpha)_{-1}S_3(\alpha))\otimes e^{\alpha}$ (charge $\alpha$,
Fock degree $2$, weight $2+q_2(1)=3$). The result has charge $0$ and weight $3+3-3-1=2$, so it lies in the
two-dimensional space spanned by $\alpha(-1)^2\leighteenvac$ and $\alpha(-2)\leighteenvac$; expanding the normal-ordered
exponentials gives its two coordinates, which are compared with those of
$\omega_{\lambda_2}=\frac18\alpha(-1)^2\leighteenvac+\frac14\alpha(-2)\leighteenvac$. The other two identities are the same
computation in charge $-\alpha$, weight $5$, and charge $0$, weight $4$, where the spaces are larger but
still finite-dimensional.
\end{proof}
\begin{physics}
These are the most singular terms of the OPEs $W^-(z)W^0(w)$ and $W^-(z)W^+(w)$ of the $c=-2$
triplet algebra: $W^-(z)W^+(w)\sim\frac{-48cd}{(z-w)^4}T(w)+\cdots$ (the leading $(z-w)^{-6}$
term being the vacuum). They are Kausch's structure constants in the Coulomb-gas normalisation.
\end{physics}
\begin{status}
Proved, unconditional. \texttt{\#print axioms} of \lean{threeCoordinateIdentitiesWp_two} and
\lean{strongGenWp_two}: \lean{[propext, Classical.choice, Quot.sound]}.
\end{status}

\subsection{The Adamovi\'c--Milas identity at weight ten}

\begin{theorem}[The $p=2$ identity]\label{l18:thm:AMtwo}
In $V_{L_2}$, at weight $10$ and charge $0$,
\begin{equation*}
Q^2e^{-2\alpha}=3\,H_{(-5)}H+\tfrac{1024}{1575}\,L_{-2}^5\leighteenvac+R,
\end{equation*}
where $R$ is an explicit combination of eleven Virasoro words $L_{-n_1}\cdots L_{-n_k}\leighteenvac$ with
$n_1\ne2$ (so $R\in$ \lean{VirDeep}). Hence \lean{AMOmegaIdentity 2 hp 5} holds.
\end{theorem}
\noindent Files \leanpath{VOA/\allowbreak Examples/\allowbreak TripletWpTwoAM.lean},
\leanpath{TripletWpTwoAMFinal.lean}.
\begin{leancode}
noncomputable def virWordTwo (hp : 1 ≤ 2) : List ℕ → LatModuleGen (latP 2)
  | [] => latVacuumGen (latP 2)
  | n :: ns => (tripletLatticeVOA 2 hp).L (-(n : ℤ)) (virWordTwo hp ns)

noncomputable def amTwoRemainder (hp : 1 ≤ 2) : LatModuleGen (latP 2) :=
    (128 / 35 : ℂ) • virWordTwo hp [10]
  + (-27232 / 525 : ℂ) • virWordTwo hp [8, 2]
  + (-12308 / 525 : ℂ) • virWordTwo hp [7, 3]
  + (20896 / 525 : ℂ) • virWordTwo hp [6, 4]
  + (-106816 / 1575 : ℂ) • virWordTwo hp [6, 2, 2]
  + (18544 / 1575 : ℂ) • virWordTwo hp [5, 5]
  + (-59392 / 1575 : ℂ) • virWordTwo hp [5, 3, 2]
  + (-752 / 525 : ℂ) • virWordTwo hp [4, 4, 2]
  + (1688 / 525 : ℂ) • virWordTwo hp [4, 3, 3]
  + (-512 / 63 : ℂ) • virWordTwo hp [4, 2, 2, 2]
  + (832 / 315 : ℂ) • virWordTwo hp [3, 3, 2, 2]

def AMTwoCoordinateIdentity (hp : 1 ≤ 2) : Prop :=
  Qwp 2 hp (Qwp 2 hp
      (latChargeGen (latP 2) (-alphaP 2 + -alphaP 2)))
    = (3 : ℂ) • (tripletLatticeVOA 2 hp).smul (Htrip 2) (-5) (Htrip 2)
      + (1024 / 1575 : ℂ) • virWordTwo hp [2, 2, 2, 2, 2]
      + amTwoRemainder hp
\end{leancode}
\begin{leancode}
theorem amOmegaIdentity_two_of_coordinateIdentity (hp : 1 ≤ 2)
    (h : AMTwoCoordinateIdentity hp) : AMOmegaIdentity 2 hp 5 := by

theorem amTwoCoordinateIdentity_two (hp : 1 ≤ 2) : AMTwoCoordinateIdentity hp :=

theorem amOmegaIdentity_two (hp : 1 ≤ 2) : AMOmegaIdentity 2 hp 5 :=
\end{leancode}
\begin{dictionary}
\lean{virWordTwo hp [n₁,…,n_k]} is $L_{-n_1}\cdots L_{-n_k}\leighteenvac$ with the deformed Virasoro modes of
$V_{L_2}$ (the list is read left to right, outermost first). So \lean{virWordTwo hp [2,2,2,2,2]} is
$L_{-2}^5\leighteenvac$, and every word in \lean{amTwoRemainder} has first letter $\ne2$, which is why it lies in
\lean{VirDeep 2 hp} (the first letter is the outermost mode; its argument is a Virasoro descendant of the
vacuum, which lies in $W(2)$). The coefficient $B_2=\frac{1024}{1575}$ agrees with Adamovi\'c--Milas'
closed formula $(-1)^p\binom{2p}{p}(4p)^{3p-1}/\bigl(\binom{4p-1}{p}((3p-1)!)^2\bigr)$ at $p=2$:
$6\cdot8^5/(21\cdot120^2)=196608/302400=1024/1575$ (we checked the arithmetic). The twelve words are all
partitions of $10$ into parts $\ge2$, as the grading requires. The identity is an equality of two explicit
vectors of the Fock space; it is proved, not assumed.
\end{dictionary}
\begin{proof}[Proof on paper]
Write both sides as polynomials in the oscillators $a_k=\alpha(-k)$ of total degree $10$ (charge $0$):
$Q^2e^{-2\alpha}$ is $\mathrm{EE}(\alpha)_3S_7(\alpha)$ up to sign, $H=S_3(\alpha)$ and $H_{(-5)}H$ is computed by
the lattice product formula, and each Virasoro word is computed by the background-charge Virasoro action on
the Fock space at $c=-2$. There are $42$ monomials of degree $10$ that occur; comparing their coefficients
is a finite linear-algebra check. The library assembles it from a ``sparse'' representation (32-term and
42-term ledgers of monomials) in the \leanpath{TripletWpTwoAM*.lean} files.
\end{proof}
\begin{physics}
The weight-$10$ null relation of the $c=-2$ triplet model: the fifth power of the stress tensor is,
modulo $C_2$, a combination of composites involving the triplet fields. It is what bounds the ``number of
states'' in the classical limit.
\end{physics}
\begin{status}
Proved, unconditional; \texttt{\#print axioms VOA.LatGen.WpTwoAMFinal.amOmegaIdentity\_two}:
\lean{[propext, Classical.choice, Quot.sound]}.
\end{status}

\subsection{\texorpdfstring{$C_2$}{C2}-cofiniteness}

\begin{theorem}[The lattice triplet $W(2)$ is $C_2$-cofinite]\label{l18:thm:c2two}
$W(2)=$ \lean{tripletWVOA 2 hp} is $C_2$-cofinite, and $\dim W(2)/C_2(W(2))\le80$.
\end{theorem}
\noindent File \leanpath{VOA/\allowbreak Examples/\allowbreak TripletWpTwoClose.lean}.
\begin{leancode}
theorem tripletWVOA_two_c2Cofinite_and_finrank_le (hp : 1 ≤ 2) :
    C2Cofinite (tripletWVOA 2 hp)
      ∧ Module.finrank ℂ (A₂ (tripletWVOA 2 hp)) ≤ 80 :=

theorem tripletWVOA_two_c2Cofinite (hp : 1 ≤ 2) : C2Cofinite (tripletWVOA 2 hp) :=
\end{leancode}
\begin{dictionary}
The only argument is \lean{hp : 1 ≤ 2}, which is a proof of a true inequality between numerals: the
theorem is unconditional. \lean{C2Cofinite} is the standard notion (Definition~\ref{l12:def:c2}); the $C_2$
space is that of the subalgebra itself. $80=16\cdot5$. The theorem is about the \emph{lattice} model; it
says nothing about \lean{SympFerm.tripletW2} (whose $C_2$-cofiniteness is open, Section~\ref{l18:sec:abe}).
Could it be vacuous? $C_2$-cofiniteness of the zero algebra is trivial, but $W(2)\ne0$; and $W(2)$ is not
the (trivially $C_2$-cofinite) $\C\leighteenvac$, since it contains $F$ of charge $-\alpha$. A Lean check
(not part of this release) reproduces the bound from the general-$p$ theorem.
\end{dictionary}
\begin{proof}[Proof on paper]
Theorem~\ref{l18:thm:pairs} with \lean{LoweringPairsWp 2 hp} from Theorem~\ref{l18:thm:mixed} and
\lean{AMOmegaIdentity 2 hp 5} from Theorem~\ref{l18:thm:AMtwo}.
\end{proof}
\begin{physics}
The $c=-2$ triplet theory is a finite logarithmic CFT: its classical phase space is a point
($\dim R\le11$ by Abe, here bounded by $80$). This is the property that guarantees finitely many irreducible
modules and Miyamoto's modular invariance of pseudo-traces.
\end{physics}
\begin{status}
\textbf{Proved, unconditional, for the lattice model} \lean{tripletWVOA 2 hp} only. \texttt{\#print axioms VOA.LatGen.WpC2.}\allowbreak\texttt{tripletWVOA\_two\_c2Cofinite\_and\_finrank\_le}:
\lean{[propext, Classical.choice, Quot.sound]}. No \lean{VOA.pbw}.
\end{status}

\subsection{Irrationality}

\begin{theorem}[The lattice $W(2)$ is not rational]\label{l18:thm:irrtwo}
$W(2)$ has no nonzero state of weight one. The ambient $V_{L_2}$, regarded as an ordinary $W(2)$-module by
restriction, is not completely reducible: the submodule $W(2)$ has no complement. Consequently $W(2)$ is
neither Gannon-rational nor admissibly rational.
\end{theorem}
\noindent Files \leanpath{VOA/\allowbreak Examples/\allowbreak TripletWpTwoFelderStart.lean},
\leanpath{TripletWpTwoIrrational.lean}.
\begin{leancode}
noncomputable def tripletAmbientOrdModuleTwo (hp : 1 ≤ 2) :
    OrdVModule (tripletWVOA 2 hp) (LatModuleGen (latP 2)) :=
  (tripletW 2 hp).adjointRestrictOrd

noncomputable def tripletCarrierWeakSubmoduleTwo (hp : 1 ≤ 2) :
    WeakModule.Submodule' (tripletAmbientOrdModuleTwo hp).toWeakModule :=
  (tripletW 2 hp).carrierWeakSubmodule
\end{leancode}
\begin{leancode}
theorem eq_zero_of_mem_tripletW_two_and_grading_one (hp : 1 ≤ 2)
    {x : LatModuleGen (latP 2)} (hxW : x ∈ (tripletW 2 hp).carrier)
    (hx1 : x ∈ (tripletLatticeVOA 2 hp).grading 1) : x = 0 := by

theorem Ftrip_three_latChargeGen_alpha (hp : 1 ≤ 2) :
    (tripletLatticeVOA 2 hp).smul (Ftrip 2) 3
        (latChargeGen (latP 2) (alphaP 2))
      = latticeCocycle (latP 2) (-alphaP 2) (alphaP 2) •
          latChargeGen (latP 2) (-alphaP 2 + alphaP 2) := by

theorem tripletAmbientOrdModuleTwo_not_completelyReducible (hp : 1 ≤ 2) :
    ¬ OrdVModule.IsCompletelyReducible (tripletAmbientOrdModuleTwo hp) := by

theorem not_gannonRational_tripletWVOA_two (hp : 1 ≤ 2) :
    ¬ GannonRational (tripletWVOA 2 hp) := by

theorem not_admissiblyRational_tripletWVOA_two (hp : 1 ≤ 2) :
    ¬ AdmissiblyRational (tripletWVOA 2 hp) := by

theorem tripletWVOA_two_c2Cofinite_and_not_admissiblyRational (hp : 1 ≤ 2) :
    C2Cofinite (tripletWVOA 2 hp) ∧ ¬ AdmissiblyRational (tripletWVOA 2 hp) :=
\end{leancode}
\begin{dictionary}
\lean{adjointRestrictOrd} is the restriction of the adjoint module of $V_{L_2}$ to the subalgebra
(Chapter~\ref{ch:L12}, Definition~\ref{l12:def:restrict}): an \lean{OrdVModule} of $W(2)$ on the carrier
$V_{L_2}$, with weight spaces the $L_0$-eigenspaces (finite-dimensional and bounded below because
$V_{L_2}$ is a VOA of CFT type). \lean{OrdVModule.IsCompletelyReducible N} means that every weak
submodule has a complementary submodule (\lean{IsCompl} of carriers; Definition~\ref{l12:def:ordmodule}).
\lean{GannonRational} and \lean{AdmissiblyRational} are the two rationality notions of
Chapter~\ref{ch:L12} (Definitions~\ref{l12:def:gannon} and~\ref{l12:def:admrat}); the first asks complete
reducibility of ordinary modules, the second of admissible ones, and the library's lemma
\lean{isCompletelyReducible_of_admissiblyRational} transfers the second to the ordinary module at hand.
\lean{AdmissiblyRational} asks that every \emph{weak} submodule of every admissible module be complemented by a weak
submodule; this is not literally Dong--Li--Mason's complete reducibility into irreducibles (release audit 2026-10-27,
finding m5). The witness is robust to the difference: the graded submodule $W(2)\subset V_{L_2}$ has no complement even
among weak submodules, because $F_{(3)}e^\alpha=\varepsilon\leighteenvac$ and $W(2)_1=0$, so on paper it also refutes DLM
rationality. When announcing, state the notion. \lean{AdmissiblyRational} is satisfiable (it holds for the trivial
VOA, \leanpath{VOA/\allowbreak Examples/\allowbreak TrivialAdmissiblyRational.lean}), so its negation has content.
\end{dictionary}
\begin{proof}[Proof on paper]
\emph{No weight one.} By strong generation (Theorem~\ref{l18:thm:mixed}) $W(2)$ is spanned by words in
generators of weight $2$ and $3$; such a word has weight $0$ (the vacuum) or $\ge2$. \emph{Non-splitting.}
Suppose $V_{L_2}=W(2)\oplus Q$ as $W(2)$-modules. Write $e^\alpha=a+q$; since $L_0$ preserves both summands and
$L_0e^\alpha=e^\alpha$, $a$ has weight one, so $a=0$ and $e^\alpha\in Q$. Then $F_{(3)}e^\alpha\in Q$. But
$F_{(3)}e^\alpha=\varepsilon(-\alpha,\alpha)\leighteenvac$ (the lattice product $e^{-\alpha}_{(n)}e^\alpha$ is
$\varepsilon z^{-4}(1+O(z))$, so its $z^{-4}$-coefficient, $n=3$, is $\varepsilon\leighteenvac$), which is a nonzero element of
$W(2)\cap Q=0$. \emph{Irrationality.} Gannon-rationality requires every ordinary module to be completely
reducible; admissible rationality implies it for ordinary modules that are admissible, which this one is.
\end{proof}
\begin{physics}
The Coulomb-gas Fock space $V_{L_2}$ is a non-split extension of $W(2)$-modules: the weight-one field
$e^{\alpha}$ (the screening current itself) is not in the chiral algebra, but its OPE with $e^{-\alpha}$
produces the identity. This is the simplest instance of the structure FGST describe as the
``Felder complex'' of the triplet theory: free-field modules glue irreducibles into indecomposables. In a
rational theory every such extension would split.
\end{physics}
\begin{status}
\textbf{Proved, unconditional, for the lattice model only}, with the library's rationality notions. \texttt{\#print axioms} of \lean{not_gannonRational_tripletWVOA_two} and
\lean{not_admissiblyRational_tripletWVOA_two}: \lean{[propext, Classical.choice, Quot.sound]}. No Felder
exactness, no module classification and no symplectic-fermion input is used.
\end{status}

\subsection{Simplicity and strong finiteness}

\begin{definition}[The $f_2$ Zhu certificate]\label{l18:def:ftwo}
Let $f_2(X)=X^2(X+\tfrac18)(X-\tfrac38)(X-1)=X^5-\tfrac54X^4+\tfrac{13}{64}X^3+\tfrac3{64}X^2$. Its
\emph{state representative} in a VOA $V$ is $\omega^{*5}-\tfrac54\omega^{*4}+\tfrac{13}{64}\omega^{*3}+\tfrac3{64}\omega^{*2}$,
where $\omega^{*n}$ is the $n$-fold Zhu product $\omega*(\omega*\cdots*\leighteenvac)$. The \emph{certificate} is the
statement that this state lies in $O(V)$, i.e.\ $f_2([\omega])=0$ in $A(V)$.
\end{definition}
\noindent File \leanpath{VOA/\allowbreak Examples/\allowbreak TripletWpSimplicity.lean}.
\begin{leancode}
noncomputable def zhuOmegaPow (W : VertexOperatorAlgebra V) : ℕ → V
  | 0 => W.vacuum
  | n + 1 => zhuMul W W.conformalVector (zhuOmegaPow W n)

noncomputable def fTwoZhuRepresentative (W : VertexOperatorAlgebra V) : V :=
  zhuOmegaPow W 5
    - (5 / 4 : ℂ) • zhuOmegaPow W 4
    + (13 / 64 : ℂ) • zhuOmegaPow W 3
    + (3 / 64 : ℂ) • zhuOmegaPow W 2

noncomputable def fTwoScalar (h : ℂ) : ℂ :=
  h ^ 2 * (h + 1 / 8) * (h - 3 / 8) * (h - 1)

def FTwoZhuCertificate (W : VertexOperatorAlgebra V) : Prop :=
  fTwoZhuRepresentative W ∈ O W

theorem fTwoScalar_natCast_eq_zero_iff (n : ℕ) :
    fTwoScalar (n : ℂ) = 0 ↔ n = 0 ∨ n = 1 := by
\end{leancode}
\begin{dictionary}
\lean{mk_fTwoZhuRepresentative} proves that the class of the representative is
$f_2([\omega])$ evaluated in the ring $A(V)$, and \lean{fTwoPolynomial_eval} that the expanded polynomial
equals the factored \lean{fTwoScalar}; we checked the expansion by hand
($X^2(X+\frac18)(X-\frac38)(X-1)$: the product of the last three is $X^3-\frac54X^2+\frac{13}{64}X+\frac{3}{64}$). The
roots $0,1,-\frac18,\frac38$ are the lowest weights $h_{1,1},h_{2,1},h_{1,2},h_{2,2}$ of \eqref{l18:eq:kac}; the
double root at $0$ reflects the Jordan block of $A(W(2))$.
\end{dictionary}
\begin{physics}
$f_2([\omega])=0$ in Zhu's algebra says: on the ground states of any $W(2)$-module, $L_0$ satisfies
$f_2(L_0)=0$, so the possible conformal weights of ground states are $0,1,-\frac18,\frac38$ --- the $c=-2$
first-row Kac table --- and $L_0$ may have a Jordan block only at $h=0$.
\end{physics}
\begin{status}
Definitions and an arithmetic lemma; standard axioms.
\end{status}

\begin{theorem}[$W(2)$ is simple and strongly finite]\label{l18:thm:simple}
The lattice $W(2)$ satisfies the $f_2$ certificate; hence it is simple, strongly finite, and neither
admissibly nor Gannon rational (so not strongly rational).
\end{theorem}
\noindent Files \leanpath{VOA/\allowbreak Examples/\allowbreak TripletWpTwoZhuFinal.lean},
\leanpath{TripletWpStronglyFinite.lean}.
\begin{leancode}
theorem fTwoZhuCertificateWp (hp : 1 ≤ 2) : FTwoZhuCertificateWp hp :=

theorem tripletWVOA_two_isSimple (hp : 1 ≤ 2) : (tripletWVOA 2 hp).IsSimple :=

theorem tripletWVOA_two_stronglyFinite (hp : 1 ≤ 2) :
    StronglyFinite (tripletWVOA 2 hp) :=

theorem tripletWVOA_two_stronglyFinite_and_irrational (hp : 1 ≤ 2) :
    StronglyFinite (tripletWVOA 2 hp) ∧
      ¬ AdmissiblyRational (tripletWVOA 2 hp) ∧
      ¬ GannonRational (tripletWVOA 2 hp) :=
\end{leancode}
\begin{leancode}
theorem tripletWVOA_two_stronglyFinite_of_isSimple (hp : 1 ≤ 2)
    (hsimple : (tripletWVOA 2 hp).IsSimple) : StronglyFinite (tripletWVOA 2 hp) := by

theorem not_stronglyRational_tripletWVOA_two (hp : 1 ≤ 2) :
    ¬ StronglyRational (tripletWVOA 2 hp) := by
\end{leancode}
\begin{dictionary}
\lean{FTwoZhuCertificateWp hp} is \lean{FTwoZhuCertificate (tripletWVOA 2 hp)}. \lean{W.IsSimple} is
``$V\ne0$ and every ideal is $0$ or $V$'' (Chapter~\ref{ch:L04}, Definition~\ref{l04:def:simple};
Chapter~\ref{ch:L12}, Definition~\ref{l12:def:phi}). \lean{StronglyFinite W} is the structure with fields
\lean{simple}, \lean{selfDual} (a nondegenerate fully invariant bilinear form) and \lean{c2Cofinite}
(Definition~\ref{l12:def:strong}); it is the Creutzig--Gannon/McRae notion. \lean{SelfDual} here is the
genuine nondegenerate-form notion of \leanpath{StronglyRational.lean}, not the vacuous
\lean{SelfContragredient} (satisfied by every module; Chapter~\ref{ch:L12}); it is obtained from the general
recognition theorem \lean{stronglyFinite_of_L_one_grading_one_eq_zero} (simple $+$ $L_1V_1=0$ $\Rightarrow$
self-dual, Li), which applies because $V_1=0$.
\end{dictionary}
\begin{proof}[Proof on paper]
\emph{Simplicity from the certificate} (Chapter~\ref{ch:L12}, Theorem~\ref{l12:thm:zhupoly}): if $I$ is a
proper nonzero ideal, its lowest weight space $I_n$ is a module for $A(V)$ on which $[\omega]$ acts by $n$,
so $f_2(n)=0$, so $n\in\{0,1\}$; $n=0$ is impossible since $V_0=\C\leighteenvac$ and $\leighteenvac\notin I$;
$n=1$ is impossible since $V_1=0$. \emph{The certificate}: the representative is an explicit state of weight
$\le10$; it is written as a sum of $37$ ``circle'' elements $u\circ v=\Res_zY(u,z)v\frac{(1+z)^{\wt u}}{z^2}$
(which are in $O(V)$ by definition) plus a multiple of the screened state $Q^2e^{-2\alpha}$, which lies in
$O(W(p))$ for every $p$ by a separate higher-pole argument (\lean{Qwp_sq_latChargeGen_mem_O},
\leanpath{TripletWpZhuScreening.lean}); the decomposition is an identity in the Fock space verified by the
\leanpath{TripletWpTwoZhuSparse*.lean} files. \emph{Strong finiteness}: simplicity $+$ self-duality (from
$V_1=0$) $+$ $C_2$ (Theorem~\ref{l18:thm:c2two}).
\end{proof}
\begin{physics}
Strongly finite and non-rational is the defining property of a ``finite logarithmic CFT'' in the modern
sense (Creutzig--Gannon, McRae): its module category is a finite non-semisimple tensor category. $W(2)$
is the standard example, and the library now has it as a theorem.
\end{physics}
\begin{status}
\textbf{Proved, unconditional, for the lattice model only.} \texttt{\#print axioms} of \lean{tripletWVOA_two_isSimple},
\lean{tripletWVOA_two_stronglyFinite_and_irrational}, \lean{not_stronglyRational_tripletWVOA_two}:
\lean{[propext, Classical.choice, Quot.sound]}. (The docstring of \leanpath{TripletWpStronglyFinite.lean}
still says strong finiteness is ``conditional only on simplicity''; the condition is discharged in
\leanpath{TripletWpTwoZhuFinal.lean}, a later file.)
\end{status}

\begin{theorem}[The screened state is in Zhu's $O$, for every $p$]\label{l18:thm:zhuscreen}
For every $p\ge1$, $\varepsilon(-\alpha,-\alpha)^{-1}(E_{(-2p-1)}F+2H_{(-2p-1)}H+F_{(-2p-1)}E)$, which equals
$Q^2e^{-2\alpha}$ after inclusion in $V_{L_p}$, lies in $O(W(p))$.
\end{theorem}
\noindent File \leanpath{VOA/\allowbreak Examples/\allowbreak TripletWpZhuScreening.lean}.
\begin{leancode}
noncomputable def qwpSqNegTwoAlphaTriplet (p : ℕ) (hp : 1 ≤ p) :
    (tripletW p hp).carrier :=
  (latticeCocycle (latP p) (-alphaP p) (-alphaP p))⁻¹ •
    qwpSqModeTriplet p hp (-(2 * (p : ℤ)) - 1)

theorem Qwp_sq_latChargeGen_mem_O (p : ℕ) (hp : 1 ≤ p) :
    qwpSqNegTwoAlphaTriplet p hp ∈ O (tripletWVOA p hp) := by
\end{leancode}
\begin{dictionary}
\lean{qwpSqModeTriplet p hp n} is $E_{(n)}F+2H_{(n)}H+F_{(n)}E$ as an element of $W(p)$;
\lean{qwpSqNegTwoAlphaTriplet_coe} identifies its normalised $n=-2p-1$ value with $Q^2e^{-2\alpha}$. $O(V)$ is
Zhu's subspace (not $C_2$); this is the Zhu-algebra analogue of the $C_2$ membership in
Definition~\ref{l18:def:AMid}, proved independently.
\end{dictionary}
\begin{proof}[Proof on paper]
The Frenkel--Zhu residue $\Res_zY(u,z)v\,(1+z)^{\wt u}z^{-2-k}$ lies in $O(V)$ for all $k\ge0$. Take $k=2p-1$ and
sum the three residues: the expansion of $(1+z)^{2p-1}$ produces the modes $u_{(-2p-1)},u_{(-2p)},\dots$; the
leading one gives the screened state, and every later mode $n\ge-2p$ vanishes by
$Q^2(F_{(n)}F)=0$ (Theorem~\ref{l18:thm:AMrel}).
\end{proof}
\begin{physics}
In Zhu's algebra the doubly screened vertex operator is zero: this is the relation that, combined with
the Virasoro part of \eqref{l18:eq:AMid} in its Zhu form, yields the polynomial $f_2$.
\end{physics}
\begin{status}
Proved, unconditional; standard axioms.
\end{status}

\begin{warning}
\textbf{Sizes.} About $140$ of the files of this chapter (\leanpath{TripletWpTwoAM*},
\leanpath{TripletWpTwoZhuSparse*}, \leanpath{TripletWpTwoZhuFast*}, \leanpath{TripletWpTwoVirasoroCoords},
\leanpath{TripletWpTwoSector}, \leanpath{TripletWpTwoZhuCertificate}) are split pieces of the two Fock-space
computations behind Theorems~\ref{l18:thm:AMtwo} and~\ref{l18:thm:simple}; they were split for elaboration
time. Each file carries its own \texttt{\#print axioms} audit. Their public interfaces are the theorems
quoted above; the intermediate declarations are bookkeeping (ledgers of monomials, mode-by-mode partial
sums) with no independent mathematical content, and the book does not document them one by one. A referee
who wants to audit them should check (i) that no \lean{native_decide} or \lean{sorry} occurs
(\texttt{scripts/check\_repository.py} enforces this) and (ii) that the two final comparisons
(\lean{eval_hProductExpected_eq_assembly}, \lean{eval_hEntryMode0_eq_expected}) are equalities of
evaluated polynomials, not definitional unfoldings; for the first, the library deliberately keeps two
separately written copies of the $32$-term ledger and compares them.
\end{warning}

%======================================================================
\section{The Abe model: \texorpdfstring{$W(2)$}{W(2)} as the even part of symplectic fermions}\label{l18:sec:abe}
%======================================================================

\subsection{The algebra}

\begin{definition}[The Abe model of $W(2)$]\label{l18:def:tripletW2}
$\leighteenSF^+$ is the even part of the symplectic-fermion super VOA \lean{sympFermSVOA}
(Chapter~\ref{ch:L17}, Theorem~\ref{l17:thm:sfsvoa}): the span of the word states
$\chi^{a_1}_{-k_1-1}\cdots\chi^{a_n}_{-k_n-1}\leighteenvac$ with $n$ even, with the restricted vertex
operators, vacuum and conformal vector $\omega_{SF}=-\chi^0_{-1}\chi^1_{-1}\leighteenvac$.
\end{definition}
\noindent File \leanpath{VOA/\allowbreak Examples/\allowbreak TripletW2.lean}.
\begin{leancode}
theorem sympFermSVOA_integralWeights (h : ℂ) (hh : ¬ ∃ n : ℤ, (n : ℂ) = h) :
    ∀ x ∈ sympFermSVOA.parity 0, x ∈ sympFermSVOA.grading h → x = 0 := by

noncomputable def tripletW2 : VertexOperatorAlgebra (evenOdd QSF 0) :=
  sympFermSVOA.evenPartVOA sympFermSVOA_integralWeights

@[simp] theorem tripletW2_centralCharge : tripletW2.centralCharge = -2 := rfl

@[simp] theorem tripletW2_grading (n : ℤ) :
    tripletW2.grading n = (gradingSF (n : ℂ)).comap (evenOdd QSF 0).subtype := rfl

theorem evenPartSF_eq_span :
    (evenOdd QSF (0 : ZMod 2) : Submodule ℂ FockSF)
      = Submodule.span ℂ (stWordSF '' {l : SFWord | Even l.length}) := by

theorem tripletW2_grading_one : tripletW2.grading 1 = ⊥ := by
\end{leancode}
\begin{dictionary}
\lean{evenOdd QSF 0} is the even part of the Clifford (here exterior) algebra $\mathrm{Fock}_{SF}$
(\lean{FockSF := CliffordAlgebra QSF} with \lean{QSF = 0}); \lean{evenPartVOA} is the even-part
transfer of Chapter~\ref{ch:L17}, Theorem~\ref{l17:thm:evenpart}, whose only hypothesis (integral weights
in the even part) is discharged by \lean{sympFermSVOA_integralWeights}, which in fact shows that every
weight of $\mathrm{Fock}_{SF}$ is an integer. \lean{stWordSF l} is the word state of a list
$l=[(k_1,a_1),\dots]$, i.e.\ $\chi^{a_1}_{-k_1-1}\cdots\leighteenvac$, of weight $\sum(k_i+1)$. So this \emph{is}
$\leighteenSF^+$: the identifications of $Y$, vacuum, $\omega$ and grading with the ambient ones are
\lean{rfl}-lemmas (\lean{tripletW2_Y_coe}, \dots). The name ``\lean{tripletW2}'' rests on the literature
theorem $\leighteenSF^+\cong W(2)$ (Kausch; Abe, Remark~3.9); \emph{no} isomorphism with
\lean{tripletWVOA 2 hp} is proved. $W(2)_1=0$ because weight-one states are one-letter words, which are odd.
\end{dictionary}
\begin{physics}
The bosonic (fermion-number-even) sector of the $c=-2$ symplectic-fermion theory. As a chiral algebra this
is the $\Z/2$-orbifold $\leighteenSF/\langle(-1)^F\rangle$; its graded dimension is \eqref{l18:eq:char}:
$1,0,1,4,5,8,10,\dots$ The absence of weight-one states means there is no continuous symmetry (no
currents), as for $V^\natural$.
\end{physics}
\begin{status}
Proved, unconditional; \texttt{\#print axioms VOA.SympFerm.tripletW2}, \lean{tripletW2_grading_one}:
\lean{[propext, Classical.choice, Quot.sound]}.
\end{status}

\subsection{Abe's generators}

\begin{definition}[$\omega,E,H,F$ in the Abe model]\label{l18:def:abegens}
With $e=\chi^1_{-1}\leighteenvac$, $f=\chi^0_{-1}\leighteenvac$ (the normalisation for which $e_{(-1)}f=\omega_{SF}$),
put $\omega=e_{(-1)}f$, $E=e_{(-2)}e$, $F=f_{(-2)}f$, $H=\frac12(e_{(-2)}f+f_{(-2)}e)$. In word states,
$E=\chi^1_{-2}\chi^1_{-1}\leighteenvac$, $F=\chi^0_{-2}\chi^0_{-1}\leighteenvac$, $H=\frac12(\chi^1_{-2}\chi^0_{-1}+\chi^0_{-2}\chi^1_{-1})\leighteenvac$.
\end{definition}
\noindent File \leanpath{VOA/\allowbreak Examples/\allowbreak TripletW2Gens.lean}.
\begin{leancode}
noncomputable def abeSFe : FockSF := stWordSF [((0 : ℕ), (1 : Fin 2))]

noncomputable def abeSFf : FockSF := stWordSF [((0 : ℕ), (0 : Fin 2))]

noncomputable def abeGen (a b : Fin 2) : evenOdd QSF (0 : ZMod 2) :=
  ⟨stWordSF [(1, a), (0, b)], stWord_two_mem_even _ _⟩

noncomputable def abeOmega : evenOdd QSF (0 : ZMod 2) := tripletW2.conformalVector

noncomputable def abeE : evenOdd QSF (0 : ZMod 2) := abeGen 1 1

noncomputable def abeF : evenOdd QSF (0 : ZMod 2) := abeGen 0 0

noncomputable def abeH : evenOdd QSF (0 : ZMod 2) :=
  (2⁻¹ : ℂ) • (abeGen 1 0 + abeGen 0 1)
\end{leancode}
\begin{leancode}
theorem abeOmega_eq :
    ((abeOmega : evenOdd QSF (0 : ZMod 2)) : FockSF) = mode (YSF abeSFe) (-1) abeSFf :=

theorem abeE_eq :
    ((abeE : evenOdd QSF (0 : ZMod 2)) : FockSF) = mode (YSF abeSFe) (-2) abeSFe :=

theorem abeH_eq :
    ((abeH : evenOdd QSF (0 : ZMod 2)) : FockSF)
      = (2⁻¹ : ℂ) • (mode (YSF abeSFe) (-2) abeSFf + mode (YSF abeSFf) (-2) abeSFe) := by

theorem abeGens_linearIndependent :
    LinearIndependent ℂ ![abeOmega, abeE, abeH, abeF] := by

theorem abeE_eq_neg_Wminus : abeE = -Wminus := by
\end{leancode}
\begin{dictionary}
The word list \lean{[(1, a), (0, b)]} is $\chi^a_{-2}\chi^b_{-1}\leighteenvac$ (a letter $(k,a)$ is $\chi^a_{-k-1}$).
The \lean{_eq} lemmas identify the definitions with Abe's $n$-th products, using that the field of a
one-letter word is the bare fermion field. Linear independence of the four (and of the four weight-3 states
$\chi^a_{-2}\chi^b_{-1}\leighteenvac$, \lean{abeGen_linearIndependent}) is proved with an explicit dual pairing, so
none of them is zero. \lean{abeE_eq_neg_Wminus} and its siblings identify $E,F,H$ (up to sign and $\frac12$)
with the triplet $W^-,W^+,W^0$ of \leanpath{TripletTTNull.lean}, which are proved Virasoro primary
(\lean{LSF_one_WplusSF}, \dots). Normalisation: the opposite labelling $e=\chi^0$ would give
$e_{(-1)}f=-\omega_{SF}$; the file takes the one with $+$.
\end{dictionary}
\begin{physics}
$\omega$ is the stress tensor ${:}\chi^1\chi^0{:}$; $E,H,F$ are the three symmetric bilinears
${:}\partial\chi^a\chi^b{:}$, the $W^{\pm,0}$ fields of weight $3$; the antisymmetric bilinear is $\partial T$. They
form an $\mathfrak{sl}_2$-triplet under the $\mathfrak{sp}(2)=\mathfrak{sl}_2$ rotating $\chi^0,\chi^1$.
\end{physics}
\begin{status}
Proved, unconditional; standard axioms.
\end{status}

\subsection{\texorpdfstring{$\bar\omega^2\neq0$}{omega squared is nonzero}: no low-weight null vector}

\begin{theorem}[$\omega_{(-1)}\omega\notin C_2(W(2))$]\label{l18:thm:TTnull}
In the Abe model, $\bar\omega^2\ne0$ in $R_{W(2)}$; equivalently ${:}TT{:}=\omega_{(-1)}\omega$ is not in
$C_2(W(2))=L_{-1}W(2)_3$ at weight $4$.
\end{theorem}
\noindent File \leanpath{VOA/\allowbreak Examples/\allowbreak TripletTTNull.lean}.
\begin{leancode}
theorem omega_sq_not_mem_O₂ :
    poissonMul tripletW2 tripletW2.conformalVector tripletW2.conformalVector
      ∉ O₂ tripletW2 := by

theorem poissonMul_omega_omega_ne_zero_in_A₂ :
    mk₂ tripletW2 (poissonMul tripletW2 tripletW2.conformalVector tripletW2.conformalVector)
      ≠ 0 := by
\end{leancode}
\begin{dictionary}
\lean{poissonMul W u v} $=u_{(-1)}v$. The statement is about the even part's own $C_2$; in the ambient
symplectic fermions the corresponding membership is \emph{true}
($\chi^0_{-2}\chi^1_{-2}\leighteenvac=(\chi^0_{-1}\leighteenvac)_{(-2)}\chi^1_{-2}\leighteenvac$), but the first factor is odd and so not
available in $W(2)$.
\end{dictionary}
\begin{proof}[Proof on paper]
At weight $4$, $C_2(W(2))_4$ is spanned by $u_{(-2)}v$ with $\wt u+\wt v=3$; since $W(2)_1=0$ and
$\leighteenvac_{(-2)}=0$, only $u_{(-2)}\leighteenvac=Tu$, $u\in W(2)_3$, survive. In the five-dimensional weight-$4$
space with basis $e=\chi^0_{-2}\chi^1_{-2}\leighteenvac$, $w_{ab}=\chi^a_{-1}\chi^b_{-3}\leighteenvac$, one computes
$TW^+=2w_{00}$, $TW^-=2w_{11}$, $TW^0=2(w_{01}+w_{10})$, $T^2\omega=-2e+2(w_{10}-w_{01})$ and
$\omega_{(-1)}\omega=w_{10}-w_{01}$. The functional $F(e)=2$, $F(w_{10})=1$, $F(w_{01})=-1$, $F(w_{00})=F(w_{11})=0$ kills
the four translates and gives $F(\omega_{(-1)}\omega)=2$.
\end{proof}
\begin{physics}
There is no null vector at level $4$ that kills ${:}TT{:}$ in the classical limit; the triplet theory is not
``Virasoro-minimal-like'' at low level. The relations that make $R_{W(2)}$ finite appear only at weight $6$
and above.
\end{physics}
\begin{status}
Proved, unconditional; \texttt{\#print axioms VOA.SympFerm.omega\_sq\_not\_mem\_O₂}: standard three.
\end{status}

\subsection{Abe's relations and nilpotency}

\begin{theorem}[Abe's $C_2$ relations in the Abe model]\label{l18:thm:abeRel}
In $R_{W(2)}$, with $\bar\omega,\bar E,\bar H,\bar F$ the classes of the generators:
\begin{equation*}
\bar E^2=\bar F^2=0,\quad \bar E\bar H=\bar F\bar H=0,\quad 2\bar H^2+\bar E\bar F=0,\quad
2\bar H^2=\tfrac{16}{9}\bar\omega^3,\quad \bar H^4=0,\quad \bar\omega^6=0 .
\end{equation*}
Abe's sharper $\bar\omega^5=0$ follows from his (unproved here) $\bar\omega^2\bar H=0$.
\end{theorem}
\noindent File \leanpath{VOA/\allowbreak Examples/\allowbreak TripletW2AbeRelations.lean}.
\begin{leancode}
noncomputable def wbar : A₂ tripletW2 := mk₂ tripletW2 abeOmega

noncomputable def Ebar : A₂ tripletW2 := mk₂ tripletW2 abeE

noncomputable def Hbar : A₂ tripletW2 := mk₂ tripletW2 abeH

noncomputable def Fbar : A₂ tripletW2 := mk₂ tripletW2 abeF

theorem Ebar_sq : Ebar * Ebar = 0 := by

theorem Fbar_sq : Fbar * Fbar = 0 := by

theorem Ebar_mul_Hbar : Ebar * Hbar = 0 := by

theorem two_Hbar_sq_add_Ebar_Fbar : (2 : ℂ) • (Hbar * Hbar) + Ebar * Fbar = 0 := by

theorem two_Hbar_sq_eq : (2 : ℂ) • (Hbar * Hbar) = ((16 : ℂ) / 9) • (wbar * (wbar * wbar)) := by

theorem Hbar_pow_four : Hbar ^ 4 = 0 := by

theorem wbar_pow_six : wbar ^ 6 = 0 := by

def OmegaSqHbarZero : Prop := wbar * (wbar * Hbar) = 0

def OmegaPowFiveZero : Prop := wbar ^ 5 = 0

theorem omegaPowFive_of_omegaSqHbar (h : OmegaSqHbarZero) : OmegaPowFiveZero := by
\end{leancode}
\begin{dictionary}
The product \lean{*} on \lean{A₂ tripletW2} is the commutative product induced by $u_{(-1)}v$
(\lean{A₂_mul_mk}). The first two relations are proved in the stronger form $E_{(-1)}E=0$ and $F_{(-1)}F=0$
\emph{exactly} in $\mathrm{Fock}_{SF}$; the others are proved by computing both sides in the word basis and
exhibiting the difference as an explicit element of $C_2(W(2))$ (translates, plus one generator
$\omega_{(-2)}G^{10}$). \lean{OmegaSqHbarZero} and \lean{OmegaPowFiveZero} are open named \lean{Prop}s, not
assumed anywhere except as explicit hypotheses. The coefficient $\frac{16}9$ is Abe's.
\end{dictionary}
\begin{proof}[Proof on paper]
$E$ and $F$ each live in one fermion flavour; $\chi^a\chi^a$ has no contraction (the symplectic form is
antisymmetric), so $E_{(-1)}E$ is a four-letter word in one flavour of weight $6$, which would need four
distinct levels summing to $6$ --- impossible, the minimum is $1+2+3+4=10$. The mixed relations are explicit
Wick computations at weights $6$. Then $0=\bar E^2\bar F^2=(\bar E\bar F)^2=4\bar H^4$, and
$0=\bar H^4=(\tfrac89)^2\bar\omega^6$.
\end{proof}
\begin{physics}
The weight-$6$ null relations of the $c=-2$ triplet theory in the classical limit: the triplet fields square to
zero and their symmetric products are proportional to $T^3$. They make the associated scheme a fat point.
\end{physics}
\begin{status}
Proved, unconditional; \texttt{\#print axioms VOA.AbeRel.wbar\_pow\_six}, \lean{Hbar_pow_four}:
\lean{[propext, Classical.choice, Quot.sound]}. Abe's weight-$7$ relations $\bar\omega^2\bar E=\bar\omega^2\bar H=\bar\omega^2\bar F=0$
and $\bar\omega^5=0$ are \textbf{not} proved; Abe himself states the underlying four-term formula without derivation
(recorded in the file).
\end{status}

\subsection{\texorpdfstring{$C_2$}{C2}-cofiniteness of the Abe model: the one remaining gate}

\begin{theorem}[$C_2$-cofiniteness of $\leighteenSF^+$, conditional]\label{l18:thm:abeC2}
If $R_{W(2)}$ is generated as a commutative algebra by $\bar\omega,\bar E,\bar H,\bar F$, then the Abe model is
$C_2$-cofinite with $\dim R_{W(2)}\le96$. Strong generation of $W(2)$ by $\omega,E,H,F$ implies this hypothesis.
Alternatively, if $W(2)$ is spanned modulo $C_2$ by Abe's eleven states, it is $C_2$-cofinite.
\end{theorem}
\noindent Files \leanpath{VOA/\allowbreak Examples/\allowbreak TripletW2AbeAdjoin.lean},
\leanpath{TripletW2Gens.lean}.
\begin{leancode}
def AbeQuotientGeneration : Prop :=
  Algebra.adjoin ℂ ({wbar, Ebar, Hbar, Fbar} : Set (A₂ tripletW2)) = ⊤

theorem abeQuotientGeneration_of_strongGenAbe (h : StrongGenAbe) :
    AbeQuotientGeneration := by

theorem c2Cofinite_tripletW2_of_abeQuotientGeneration
    (h : AbeQuotientGeneration) : C2Cofinite tripletW2 := by

theorem finrank_A₂_tripletW2_le_96_of_abeQuotientGeneration
    (h : AbeQuotientGeneration) : Module.finrank ℂ (A₂ tripletW2) ≤ 96 := by
\end{leancode}
\begin{leancode}
noncomputable def abeGens : Set (evenOdd QSF (0 : ZMod 2)) := {abeOmega, abeE, abeH, abeF}

def StrongGenAbe : Prop := iterProdSpan tripletW2 abeGens = ⊤

def AbeSpanningEleven : Prop :=
  (⊤ : Submodule ℂ (evenOdd QSF (0 : ZMod 2))) = Submodule.span ℂ abeSpanEleven ⊔ O₂ tripletW2

noncomputable def abeSpanEleven : Set (evenOdd QSF (0 : ZMod 2)) :=
  {abeOmegaPow 0, abeOmegaPow 1, abeOmegaPow 2, abeOmegaPow 3, abeOmegaPow 4,
    abeE, abeMulOmega abeE, abeH, abeMulOmega abeH, abeF, abeMulOmega abeF}

theorem c2Cofinite_of_abeSpanningEleven (h : AbeSpanningEleven) : C2Cofinite tripletW2 := by
\end{leancode}
\begin{dictionary}
$96=6\cdot2\cdot4\cdot2$ from the exponents of Theorem~\ref{l18:thm:abeRel}. \lean{abeOmegaPow n} is
$L_{-2}^n\leighteenvac$ and \lean{abeMulOmega x} $=\omega_{(-1)}x$, so \lean{abeSpanEleven} is Abe's
$\{\leighteenvac,\omega,\dots,\omega^4,E,\omega E,H,\omega H,F,\omega F\}$ (Abe, Prop.~3.11). All three
hypotheses are statements about the \emph{even part's own} $C_2$; none is proved. The library also records
(and docstrings warn against) a tempting false shortcut: \lean{SpanModO₂Abe} ($\omega,E,H,F$ span $W(2)$ modulo
$C_2$) would give $\dim R\le4$, contradicting the lower bound $9$ below; it is named only so nobody tries to
prove it.
\end{dictionary}
\begin{proof}[Proof on paper]
Theorem~\ref{l18:thm:engine} with the nilpotency of Theorem~\ref{l18:thm:abeRel}. For the eleven-state form:
a finite set spanning modulo $C_2$ gives $\dim R<\infty$ directly (\lean{c2Cofinite_of_span_sup_O₂}).
\end{proof}
\begin{physics}
What is missing is the classical-limit form of ``$W(2)$ is a $W$-algebra with generators $T,W^a$'' --- the
statement that is easy in the lattice model (Theorem~\ref{l18:thm:mixed}) but requires, in the fermionic model,
showing that every even fermion polynomial is a normally ordered polynomial in bilinears modulo $C_2$.
\end{physics}
\begin{status}
Proved, \textbf{conditional on} \lean{AbeQuotientGeneration} (resp.\ \lean{AbeSpanningEleven}).
\texttt{\#print axioms}: standard three. \lean{C2Cofinite tripletW2} itself is \textbf{open} in the library.
\end{status}

\begin{theorem}[Further reductions of the Abe-model $C_2$ gap]\label{l18:thm:abeDirect}
(a) $V$ is $C_2$-cofinite iff $V_n\subseteq C_2(V)$ for all $n$ large. (b) The Abe model is $C_2$-cofinite as soon
as $\chi^a_{-2}v\in M+C_2(W(2))$ for all odd $v$, $a\in\{0,1\}$, and one finite-dimensional $M$. (c) Strong
generation by $\omega$ and the four bilinears together with nilpotency of their classes gives $C_2$.
\end{theorem}
\noindent Files \leanpath{VOA/\allowbreak Examples/\allowbreak TripletW2C2.lean},
\leanpath{TripletW2C2Direct.lean}, \leanpath{TripletW2Nilpotent.lean}.
\begin{leancode}
theorem C2Cofinite_iff_exists_grading_le_O₂ (W : VertexOperatorAlgebra 𝒱) :
    C2Cofinite W ↔ ∃ N : ℕ, ∀ n : ℤ, (N : ℤ) ≤ n → W.grading n ≤ O₂ W :=

noncomputable def O₂W2 : Submodule ℂ FockSF := (O₂ tripletW2).map evSF.subtype

theorem c2Cofinite_tripletW2_of_levelOne (M : Submodule ℂ FockSF) [FiniteDimensional ℂ M]
    (h : ∀ (a : Fin 2) (v : FockSF), v ∈ oddSF → sfCre 1 a * v ∈ M ⊔ O₂W2) :
    C2Cofinite tripletW2 :=

def StrongGenW2 : Prop := iterProdSpan tripletW2 (Set.range genW2) = ⊤

def NilpotentGensW2 : Prop :=
  ∀ i : GenIdx, ∃ N : ℕ, pMulPow tripletW2 (genW2 i) N ∈ O₂ tripletW2

theorem c2Cofinite_tripletW2_of_nilpotent (hgen : StrongGenW2) (hnil : NilpotentGensW2) :
    C2Cofinite tripletW2 :=
\end{leancode}
\begin{dictionary}
\lean{sfCre 1 a} $=\chi^a_{-2}\leighteenvac$ as an element of the exterior algebra, so \lean{sfCre 1 a * v} is
$\chi^a_{-2}v$; \lean{oddSF} is the odd part; \lean{O₂W2} is $C_2(W(2))$ pushed into $\mathrm{Fock}_{SF}$. In (c),
\lean{genW2} is $\omega$ together with the four states $G^{ab}=\chi^a_{-2}\chi^b_{-1}\leighteenvac$ (the antisymmetric
combination is $T\omega$). (c) is superseded by Theorem~\ref{l18:thm:abeRel} for the nilpotency half; its
generation half \lean{StrongGenW2} is open.
\end{dictionary}
\begin{warning}
\textbf{Recorded findings that bear on the attack} (\leanpath{TripletW2C2Direct.lean}, by weight counting):
$\dim W(2)_w=1,0,1,4,5,8$ for $w=0,\dots,5$, and $C_2(W(2))_w$ has dimensions $0,0,0,1,4,5$, so
$\dim R_{W(2)}\ge9$, with surviving classes $\leighteenvac,\omega,W^a,\omega^2,\omega W^a$. Hence (i) there is no
low-weight null vector that kills $R$ to dimension $3p-1=5$; (ii) the tempting ``excess filtration'' (every even
state whose creation levels sum to $\ge2$ is in $C_2$) is false, since $\omega^2$ has excess $2$; (iii) the
inclusion $C_2(W(2))\subseteq C_2(\leighteenSF)\cap W(2)$ is strict already at weight $3$ (the ambient $C_2$
contains all of $W(2)_3$, the even part's only $T\omega$), so the ambient result $\dim R_{\leighteenSF}\le4$ of
Chapter~\ref{ch:L17} does \emph{not} descend. These are paper calculations recorded in docstrings; only
$\bar\omega^2\ne0$ (Theorem~\ref{l18:thm:TTnull}) is a Lean theorem.
\end{warning}
\begin{physics}
(a) says ``finitely many classical states'' is the same as ``all sufficiently heavy states are composite modulo
derivatives''. (b) isolates the one fermionic computation still missing: that $\partial\chi^a$ times an odd state can
be rewritten, modulo composites of \emph{bosonic} fields, through finitely many states --- the precise point where the
orbifold $\leighteenSF^+$ differs from $\leighteenSF$.
\end{physics}
\begin{status}
(a) proved unconditionally (a reformulation, for any VOA); (b) and (c) proved only \textbf{conditionally} on their displayed
hypotheses, none of which is discharged. Standard axioms.
\end{status}

\paragraph{Index of the Abe-model $C_2$ reductions.} Theorems~\ref{l18:thm:abeC2}
and~\ref{l18:thm:abeDirect} give several alternative routes to \lean{C2Cofinite tripletW2}; none is
discharged. Table~\ref{l18:tab:abeC2} lists, for each theorem concluding $C_2$-cofiniteness of the Abe
model, the hypotheses it takes. Any one of the open inputs in the last column would close the gap.
\begin{table}[ht]
\centering\small
\begin{tabular}{p{7.6cm}p{3.8cm}p{2.8cm}}
\toprule
theorem & hypotheses & status of hypotheses\\
\midrule
\lean{c2Cofinite_tripletW2_of_abeQuotientGeneration} (also gives $\dim\le96$) & \lean{AbeQuotientGeneration} & open\\
\lean{abeQuotientGeneration_of_strongGenAbe} (feeds the row above) & \lean{StrongGenAbe} & open\\
\lean{c2Cofinite_of_abeSpanningEleven} & \lean{AbeSpanningEleven} & open\\
\lean{c2Cofinite_tripletW2_of_levelOne} & a finite-dimensional $M$ and the level-one condition \lean{h} & open\\
\lean{c2Cofinite_tripletW2_of_nilpotent} & \lean{StrongGenW2}, \lean{NilpotentGensW2} & first open; for the second see Thm~\ref{l18:thm:abeRel} and the dictionary above\\
\lean{C2Cofinite_iff_exists_grading_le_O₂} & none (a reformulation, any VOA) & ---\\
\bottomrule
\end{tabular}
\caption{The Abe-model $C_2$ reductions and their named hypotheses. All are \texttt{\#print axioms}-clean
(standard three axioms); \lean{C2Cofinite tripletW2} itself is not proved.}
\label{l18:tab:abeC2}
\end{table}

\subsection{Zhu's algebra and logarithmic modules}

\begin{definition}[$\omega$-Jordan condition]\label{l18:def:omegaJordan}
For a VOA $V$ and $h\in\C$ let $\omega_h=[\omega]-h[\leighteenvac]\in A(V)$. \lean{OmegaJordanAt V h} is the statement
$\omega_h\notin A(V)\,\omega_h^2$ (the left ideal generated by $\omega_h^2$).
\end{definition}
\noindent File \leanpath{VOA/\allowbreak Examples/\allowbreak TripletW2Log.lean}.
\begin{leancode}
noncomputable def omegaShift (h : ℂ) : A W :=
  mk W W.conformalVector - h • mk W W.vacuum

def OmegaJordanAt (W : VertexOperatorAlgebra 𝒱) (h : ℂ) : Prop :=
  omegaShift W h ∉
    zhuLeftMul W (Amul (zhuRingData_holds W) (omegaShift W h) (omegaShift W h))
\end{leancode}
\begin{dictionary}
\lean{zhuLeftMul W x} is the left ideal $A(V)x$ (\lean{mem_zhuLeftMul}); \lean{Amul} is Zhu's product on
$A(V)$. Since $[\omega]$ is central, the left ideal is two-sided. Equivalent forms, all proved:
\lean{omegaJordanAt_iff_exists_amod} ($\iff$ some $A(V)$-module has $\rho(\omega_h)^2=0\ne\rho(\omega_h)$) and
\lean{omegaJordanAt_of_zhuMul} (it suffices to find $v\in V$ with $(\omega-h\leighteenvac)*v\notin O(V)$ and
$(\omega-h\leighteenvac)*((\omega-h\leighteenvac)*v)\in O(V)$).
\end{dictionary}
\begin{physics}
$A(V)$ acts on ground states, $[\omega]$ by $L_0$; \lean{OmegaJordanAt V h} says that $L_0$ can have a Jordan block of
size two on a space of ground states of weight $h$. For $W(2)$ this is expected at $h=0$ (the double root of
$f_2$, Definition~\ref{l18:def:ftwo}).
\end{physics}
\begin{status}
Definition. \textbf{Not proved for any VOA and any $h$} in the library.
\end{status}

\begin{theorem}[A logarithmic module from a Zhu-algebra Jordan block]\label{l18:thm:logmodule}
Let $V$ be a VOA with $A(V)$ and all $A_{n,0}(V)$ finite-dimensional. For every $h$ there is a logarithmic module
$M_h=M(U_h)$ (the Dong--Jiang induced module on $U_h=A(V)/A(V)\omega_h^2$) on which $L_0=\rho([\omega])+n$ on level
$n$. If moreover the Dong--Jiang gate $A_{0,0}(V)=A(V)$ holds and \lean{OmegaJordanAt V h}, then
$L_0^{\mathrm{nil}}\ne0$ on $M_h$: there are $x,y$ of generalized weight $h$ with $y\ne0$ and $L_0x=hx+y$.
For the Abe model the two finiteness hypotheses follow from $C_2$-cofiniteness.
\end{theorem}
\noindent Files \leanpath{VOA/\allowbreak Examples/\allowbreak TripletW2Log.lean},
\leanpath{TripletW2C2.lean}.
\begin{leancode}
theorem djLogVModule_L0nil_ne_zero (W : VertexOperatorAlgebra 𝒱) (U : AmMod W 0) (h : ℂ) (k : ℕ)
    (hnil : amModOmegaNil U h ^ k = 0)
    (hfin : ∀ n : ℕ, FiniteDimensional ℂ (DJPiece W U n))
    (hgate : DJAnnEqAn W 0) (hne : amModOmegaNil U h ≠ 0) :
    (djLogVModule W U h k hnil hfin).L0nil ≠ 0 := by

theorem omegaJordan_nil_sq (h : ℂ) :
    amModOmegaNil (omegaJordanAMod W h).toAmModZero h ^ 2 = 0 := by

noncomputable def tripletW2Log (h : ℂ)
    (hAnm : ∀ n : ℕ, FiniteDimensional ℂ (AnmFull tripletW2 n 0))
    (hA : FiniteDimensional ℂ (A tripletW2)) :
    LogVModule tripletW2
      (DJCarrier tripletW2 (omegaJordanAMod tripletW2 h).toAmModZero) :=

theorem tripletW2Log_jordan_pair (h : ℂ)
    (hAnm : ∀ n : ℕ, FiniteDimensional ℂ (AnmFull tripletW2 n 0))
    (hA : FiniteDimensional ℂ (A tripletW2)) (hgate : DJAnnEqAn tripletW2 0)
    (hnd : OmegaJordanAt tripletW2 h) :
    ∃ x y : DJCarrier tripletW2 (omegaJordanAMod tripletW2 h).toAmModZero,
      x ∈ (tripletW2Log h hAnm hA).weight h ∧ y ∈ (tripletW2Log h hAnm hA).weight h ∧
        y ≠ 0 ∧ (tripletW2Log h hAnm hA).L0op x = h • x + y :=
\end{leancode}
\begin{leancode}
theorem tripletW2_zhuFiniteness_of_C2Cofinite (hC2 : C2Cofinite tripletW2) :
    (∀ n : ℕ, FiniteDimensional ℂ (AnmFull tripletW2 n 0))
      ∧ FiniteDimensional ℂ (A tripletW2) :=
\end{leancode}
\begin{dictionary}
\lean{AmMod W 0} is an $A_0(V)=A(V)$-module, \lean{DJCarrier W U} and \lean{DJPiece W U n} are the Dong--Jiang
induced module $M(U)=\bigoplus_nA_{n,0}(V)\otimes_{A(V)}U$ and its level-$n$ piece (Chapter~\ref{ch:L13},
Definition~\ref{l13:def:djcarrier}); \lean{DJAnnEqAn W 0} is the gate of Definition~\ref{l13:def:djann}, which
identifies $M(U)(0)$ with $U$ (at project level it is discharged by the \lean{BottomProjection} datum, which
the theorem \lean{VOA.pbw} supplies for every $V$; here it is carried as a hypothesis). \lean{amModOmegaNil U h} is
$\rho([\omega])-h$ on $U$. \lean{LogVModule} and its \lean{L0nil}, \lean{L0op}, \lean{weight} (generalized
eigenspaces) are Definition~\ref{l12:def:log}. \lean{omegaJordanAMod W h} is the quotient $A(V)/A(V)\omega_h^2$;
on it $(\rho(\omega_h))^2=0$ automatically, because $[\omega]$ is central. The Abe-model theorem therefore has
four hypotheses (\lean{hAnm}, \lean{hA}, \lean{hgate}, \lean{hnd}), which amount to three substantive ones: finiteness (implied by $C_2$, itself open --- the docstring of
\leanpath{TripletW2Log.lean} says that $C_2$ does \emph{not} imply \lean{hAnm} ``in the form needed'', but
\lean{tripletW2_zhuFiniteness_of_C2Cofinite} proves that it does: stale, documentation only), the gate, and
\lean{OmegaJordanAt tripletW2 h}.
\end{dictionary}
\begin{proof}[Proof on paper]
On $M(U)(n)$, $L_0([v]\otimes w)=[v]\otimes\rho([\omega])w+n[v]\otimes w$ (\lean{djModeZero_omega_tmul}): the level-$n$
collapse $[\omega*^n_0v]=[v*^n_0\omega]+n[v]$ in $A_{n,0}(V)$ moves $\omega$ onto $U$. So $(L_0-h-n)^k=1\otimes(\rho([\omega])-h)^k$,
and finite-dimensional pieces lying in generalized eigenspaces give a \lean{LogVModule}. If the bottom of
$M(U)$ is $U$ (the gate) and $\rho(\omega_h)\ne0$ on $U$, then $L_0^{\mathrm{nil}}\ne0$. On $U_h$,
$\rho(\omega_h)[1]=\omega_h\ne0$ exactly when $\omega_h\notin A(V)\omega_h^2$.
\end{proof}
\begin{physics}
The ``staggered'' or logarithmic module of the $c=-2$ theory: inducing from a two-dimensional space of ground
states on which $L_0$ is a Jordan block gives a module whose $L_0$ is non-diagonalisable --- the
$\mathcal R_0$ projective cover of the vacuum in the physics literature. The construction is algebraic; the
only representation-theoretic input is that such a Jordan block exists in Zhu's algebra.
\end{physics}
\begin{status}
Proved, \textbf{conditional} on four hypotheses: \lean{hAnm}, \lean{hA} (or $C_2$, itself open for the Abe model), \lean{hgate}
and \lean{hnd}. Standard axioms, pbw-free. The key hypothesis \lean{hnd : OmegaJordanAt tripletW2 h} literally asserts that
$[\omega]-h$ is not semisimple in $A(V)$; the Jordan pair in the conclusion is a formal transcription of it, and it is proved for
no VOA and no $h$. \textbf{No logarithmic module of any concrete VOA is constructed unconditionally} in the library; ``$W(2)$ has
a logarithmic module'' is not a theorem of the library.
\end{status}

\begin{theorem}[The obvious certificate cannot work]\label{l18:thm:obstruction}
If $[\omega]$ acts on an $A(V)$-module by a scalar, the module has no Jordan vector; in particular the
top level of any ordinary lowest-weight module never certifies \lean{OmegaJordanAt}.
\end{theorem}
\noindent File \leanpath{VOA/\allowbreak Examples/\allowbreak ZhuTripletW2Alt.lean}.
\begin{leancode}
def IsOmegaJordanVector (M : AMod W) (h : ℂ) (m : M.carrier) : Prop :=
  M.ρ (omegaShift W h) (M.ρ (omegaShift W h) m) = 0 ∧ M.ρ (omegaShift W h) m ≠ 0

theorem not_isOmegaJordanVector_lowestAMod {P : OrdVModule W Mod} {h₀ : ℂ}
    (hcw : IsLowestWeight P h₀) (h : ℂ) (x : (lowestAMod hcw).carrier) :
    ¬ IsOmegaJordanVector (lowestAMod hcw) h x :=

theorem omegaJordanAt_of_zhuMul {h : ℂ} (v : 𝒱)
    (hsq : zhuMul W (omegaSubV W h) (zhuMul W (omegaSubV W h) v) ∈ O W)
    (hne : zhuMul W (omegaSubV W h) v ∉ O W) :
    OmegaJordanAt W h := by
\end{leancode}
\begin{dictionary}
\lean{lowestAMod hcw} is the $A(V)$-module on the lowest-weight space of an ordinary module
(Chapter~\ref{ch:L12}, Theorem~\ref{l12:thm:moduleMh}); on it $[\omega]$ acts by $h_0$, so
$\rho(\omega_h)$ is the scalar $h_0-h$, which cannot be nonzero and square to zero. \lean{omegaSubV W h} $=\omega-h\leighteenvac$.
\end{dictionary}
\begin{physics}
Ordinary modules have diagonalisable $L_0$ on their ground states, so they cannot witness a Jordan block;
a witness must come from a genuinely logarithmic source (projective covers, or a free-field module with a
non-diagonal action). This is a no-go result for a natural but doomed strategy, not a statement about $W(2)$.
\end{physics}
\begin{status}
Proved, unconditional; standard axioms. Also proved (\leanpath{ZhuTripletW2.lean}): $A(W(2))\ne0$
(\lean{tripletW2_one_ne_zero}), and the formulas $[\omega][v]=[\omega_{(-1)}v]-k[v]$ for $v\in V_k$
(\lean{Amul_omega_mk}).
\end{status}

\subsection{Irrationality of the Abe model, conditional on the literature}

\begin{theorem}[$A(W(2))$ is not semisimple; $W(2)$ is not rational --- conditional]\label{l18:thm:abeirr}
Assume: (C2) the Abe model is $C_2$-cofinite; (blocks) every Artin--Wedderburn decomposition of $A(W(2))$ has
exactly four blocks; (five) the space of genus-one blocks of $W(2)$ has dimension $\ge5$. Then $A(W(2))$ is not
semisimple. Alternatively, (blocks) together with $\dim A(W(2))=11$ implies the same. If moreover
(Zhu) admissible rationality implies semisimplicity of $A(V)$, then $W(2)$ is not admissibly rational. (The
hypothesis (Zhu) holds for every VOA: \lean{zhuAdmissibleSemisimple_holds}, Chapter~\ref{ch:L12}.)
\end{theorem}
\noindent File \leanpath{VOA/\allowbreak Examples/\allowbreak TripletIrrational.lean}.
\begin{leancode}
def WedderburnBlockCount (n : ℕ) : Prop :=
  ∀ (m : ℕ) (d : Fin m → ℕ), (∀ i, NeZero (d i)) →
    Nonempty (A ≃ₐ[K] ∀ i : Fin m, Matrix (Fin (d i)) (Fin (d i)) K) → m = n

theorem tripletW2_not_isSemisimpleRing_A
    (hC2 : C2Cofinite tripletW2)
    (hblocks : letI := instRingA tripletW2; letI := instAlgebraA tripletW2;
      WedderburnBlockCount ℂ (A tripletW2) 4)
    (hfive : 5 ≤ finrank ℂ (Submodule.map restrictImGT (analyticBlocks tripletW2))) :
    letI := instRingA tripletW2; ¬ IsSemisimpleRing (A tripletW2) :=

theorem tripletW2_not_isSemisimpleRing_A_of_finrank_eleven
    (hfin : FiniteDimensional ℂ (A tripletW2))
    (hblocks : letI := instRingA tripletW2; letI := instAlgebraA tripletW2;
      WedderburnBlockCount ℂ (A tripletW2) 4)
    (hdim : finrank ℂ (A tripletW2) = 11) :
    letI := instRingA tripletW2; ¬ IsSemisimpleRing (A tripletW2) := by

theorem not_admissiblyRational_tripletW2
    (hC2 : C2Cofinite tripletW2)
    (hblocks : letI := instRingA tripletW2; letI := instAlgebraA tripletW2;
      WedderburnBlockCount ℂ (A tripletW2) 4)
    (hfive : 5 ≤ finrank ℂ (Submodule.map restrictImGT (analyticBlocks tripletW2)))
    (hzhu : ZhuAdmissibleSemisimple tripletW2) :
    ¬ AdmissiblyRational tripletW2 :=

theorem not_eleven_sum_four_pos_sq {a b c e : ℕ} (ha : 0 < a) (hb : 0 < b) (hc : 0 < c)
    (he : 0 < e) : a * a + b * b + c * c + e * e ≠ 11 := by
\end{leancode}
\begin{leancode}
def ZhuAdmissibleSemisimple (W : VertexOperatorAlgebra 𝒱) : Prop :=
  AdmissiblyRational W → letI := instRingA W; IsSemisimpleRing (A W)
\end{leancode}
\noindent (The last block is from \leanpath{VOA/\allowbreak Modules/\allowbreak AdmissibleRationality.lean}.)
\begin{dictionary}
\lean{instRingA}, \lean{instAlgebraA} are the ring and algebra structures on $A(V)$ (Chapter~\ref{ch:L12},
Theorem~\ref{l12:thm:zhuring}). \lean{IsSemisimpleRing} is Mathlib's. \lean{analyticBlocks W} is the space of
analytic genus-one blocks (torus one-point functions satisfying Zhu's recursion; Chapter~\ref{ch:L15},
Definition~\ref{l15:def:analyticblock}) and \lean{restrictImGT} restricts to the upper half plane; its dimension
is bounded above by $\dim\mathrm{SLF}(A(V))$, the space of symmetric linear functionals on $A(V)$, for every
$C_2$-cofinite $V$ (the Arike--Nagatomo bound, \leanpath{Modularity/BlockBoundC2.lean}). For a semisimple algebra
over $\C$ that dimension is the number of blocks. So the first theorem is: $4<5\le\dim\mathcal C_1\le\dim\mathrm{SLF}$,
impossible if semisimple. Hypotheses, with the citations the file gives:
\begin{itemize}
\item \lean{hC2}: Abe 2007, Thm.~3.9 --- not proved in the library (Theorem~\ref{l18:thm:abeC2}).
\item \lean{hblocks}: Abe 2007, Thm.~4.2 ($W(2)$ has four simple modules) with Zhu's correspondence.
  \emph{Note:} \lean{WedderburnBlockCount ℂ A n} says only ``\emph{if} $A\cong\prod_iM_{d_i}(\C)$ then $m=n$''. For a
  non-semisimple $A$ (which $A(W(2))$ is) no such isomorphism exists, so the hypothesis is \emph{vacuously true};
  it is therefore a true statement about the real $A(W(2))$, and assuming it is harmless. Its role is to fix the
  count in the hypothetical semisimple case.
\item \lean{hfive}: Arike--Nagatomo 2013, Thm.~6.3.2(4)--(5), the lower bound by explicit pseudo-trace
  construction (not their Thm.~6.3.4, which would be circular). The docstring flags that the library's block space
  differs from theirs in two conventions (equivalent on paper, not proved equivalent in Lean).
\item \lean{hdim}: Adamovi\'c--Milas 2011, $\dim A(W(p))=6p-1$. \lean{hfin} is implied by \lean{hdim}
  ($\mathrm{finrank}\ne0$ forces finite dimension), so it is redundant (moderate, minor; the docstring acknowledges it).
\item \lean{hzhu}: Dong--Li--Mason 1998, Thm.~4.10(a). \textbf{Discharged}: \lean{zhuAdmissibleSemisimple_holds} proves
  \lean{ZhuAdmissibleSemisimple W} for every \lean{W} (Chapter~\ref{ch:L12}, Theorem~\ref{l12:thm:litproved}, item~4).
\end{itemize}
\end{dictionary}
\begin{proof}[Proof on paper]
For a finite-dimensional semisimple algebra over $\C$, $A\cong\prod_{i=1}^nM_{d_i}(\C)$, symmetric functionals are
determined on $A/[A,A]\cong\C^n$, so $\dim\mathrm{SLF}(A)=n$ (\lean{finrank_symmetricFunctionals_eq_of_wedderburnBlockCount}).
Combined with the bound $\dim\mathcal C_1\le\dim\mathrm{SLF}(A(V))$ this gives $5\le4$. Route two: $\dim A=\sum d_i^2$ with
four positive $d_i$, and $11$ is not a sum of four positive squares ($1+1+1+1=4$, $4+1+1+1=7$, $4+4+1+1=10$,
$4+4+4+1=13$, $9+1+1+1=12$).
\end{proof}
\begin{physics}
The $c=-2$ theory has four irreducible representations but five independent torus amplitudes: the extra one is
the pseudo-trace on the projective cover of the vacuum, whose modular transformation produces the
$\tau$-linear term in \eqref{l18:eq:Schi}. A rational theory has as many torus amplitudes as irreducibles.
\end{physics}
\begin{status}
Proved, \textbf{conditional on} \lean{hC2}, \lean{hblocks}, \lean{hfive}, \lean{hzhu} (resp.\ \lean{hfin},
\lean{hblocks}, \lean{hdim}, \lean{hzhu}). \texttt{\#print axioms}: standard three. \lean{hzhu} is discharged by
\lean{zhuAdmissibleSemisimple_holds} (Chapter~\ref{ch:L12}); \lean{hC2} and \lean{hfive} (resp.\ \lean{hfin}, \lean{hdim})
remain, together with \lean{hblocks}, which is vacuously true for the non-semisimple $A(W(2))$ but is not discharged in
Lean. So non-rationality of the Abe model is \textbf{conditional}, not proved. \lean{¬ GannonRational tripletW2}
is \textbf{not} proved for the Abe model (contrast Theorem~\ref{l18:thm:irrtwo}, which is about the lattice model).
\end{status}

\subsection{The modular route: a \texorpdfstring{$\log q$}{log q} forbids rationality}\label{l18:sec:logmodular}

\begin{definition}[$T$-eigenfunctions, the log-free span, $\log q$ terms]\label{l18:def:logterm}
On functions $f:\mathbb H\to\C$ let $(Tf)(\tau)=f(\tau+1)$. $f$ is a \emph{$T$-eigenfunction} with eigenvalue $\mu$ if
$Tf=\mu f$. The \emph{log-free span} is the $\C$-span of all $T$-eigenfunctions. $f$ \emph{has a log term} if
$f=A+\tau B$ with $A$ in the log-free span and $B\ne0$ a $T$-eigenfunction with nonzero eigenvalue.
\end{definition}
\noindent File \leanpath{VOA/\allowbreak Examples/\allowbreak TripletLogModular.lean}.
\begin{leancode}
noncomputable def Tsh : Module.End ℂ (UpperHalfPlane → ℂ) where
  toFun f := fun τ => f ((1 : ℝ) +ᵥ τ)
  map_add' _ _ := rfl
  map_smul' _ _ := rfl

def IsTEigen (μ : ℂ) (f : UpperHalfPlane → ℂ) : Prop := Tsh f = μ • f

noncomputable def logFreeSpan : Submodule ℂ (UpperHalfPlane → ℂ) :=
  Submodule.span ℂ {f : UpperHalfPlane → ℂ | ∃ μ : ℂ, IsTEigen μ f}

def linTerm (B : UpperHalfPlane → ℂ) : UpperHalfPlane → ℂ := fun τ => (τ : ℂ) * B τ

def HasLogTerm (f : UpperHalfPlane → ℂ) : Prop :=
  ∃ (A B : UpperHalfPlane → ℂ) (ν : ℂ),
    A ∈ logFreeSpan ∧ ν ≠ 0 ∧ IsTEigen ν B ∧ B ≠ 0 ∧ f = A + linTerm B

theorem linTerm_notMem_logFreeSpan {ν : ℂ} {B : UpperHalfPlane → ℂ} (hν : ν ≠ 0)
    (hB : IsTEigen ν B) (hB0 : B ≠ 0) : linTerm B ∉ logFreeSpan := by

theorem hasLogTerm_coe : HasLogTerm (fun τ : UpperHalfPlane => (τ : ℂ)) := by
\end{leancode}
\begin{dictionary}
\lean{(1 : ℝ) +ᵥ τ} is $\tau+1$ (the additive action of $\mathbb{R}$ on $\mathbb H$). The log-free span is a span, so $A$ may
mix eigenvalues --- necessary, since in \eqref{l18:eq:Schi} $\chi_{-1/8}$ and $\chi_{3/8}$ have different
$T$-eigenvalues. \lean{hasLogTerm_coe} is a non-vacuity guard ($\tau=0+\tau\cdot1$). Writing $\tau=\log q/2\pi i$,
\lean{linTerm B} is the $\log q$ term.
\end{dictionary}
\begin{proof}[Proof on paper]
Since $(T-\nu)(\tau B)=(\tau+1)\nu B-\nu\tau B=\nu B$ and $(T-\nu)B=0$, the operator $T$ acts on
$\mathrm{span}\{\tau B,B\}$ as $\nu+N$ with $N(\tau B)=\nu B\ne0$, $N(B)=0$; in particular $\tau B$ and $B$ are linearly
independent. Suppose $\tau B=\sum_jc_jf_j$ with $Tf_j=\mu_jf_j$, and let $P$ be the product of $(X-\mu)$ over the
\emph{distinct} $\mu_j$, a polynomial without repeated roots. Then $P(T)$ kills the right side, while on the left
$P(T)(\tau B)=P(\nu)\tau B+P'(\nu)N(\tau B)=P(\nu)\tau B+\nu P'(\nu)B$. Independence forces $P(\nu)=0$ and (as
$\nu\ne0$) $P'(\nu)=0$: a double root, a contradiction.
\end{proof}
\begin{physics}
Characters of ordinary modules, $q^{h-c/24}\sum a_nq^n$, are $T$-eigenfunctions; any finite combination of them is
log-free. A term $\tau B(\tau)$ is precisely what a non-diagonalisable $L_0$ produces in a pseudo-trace
($q^{L_0}=q^he^{2\pi i\tau L_0^{\mathrm{nil}}}$). So ``no $\log q$'' is the modular shadow of ``$L_0$ diagonalisable''.
\end{physics}
\begin{status}
Proved, unconditional; \texttt{\#print axioms VOA.LogModular.linTerm\_notMem\_logFreeSpan}: standard three.
\end{status}

\begin{theorem}[A $\log q$ in $S\chi_0$ refutes rationality --- conditional]\label{l18:thm:logS}
If Zhu's modularity theorem holds for $W(2)$ in the form ``rational and $C_2$-cofinite $\Rightarrow$ the vacuum
character belongs to a modular datum satisfying Zhu's analytic law'', if $W(2)$ is $C_2$-cofinite, and if
$\chi_0(S\tau)$ has a log term, then $W(2)$ is not admissibly rational. Similarly, a logarithmic module refutes
rationality given ``rational $\Rightarrow$ every logarithmic module is $L_0$-semisimple''.
\end{theorem}
\noindent File \leanpath{VOA/\allowbreak Examples/\allowbreak TripletLogModular.lean}.
\begin{leancode}
noncomputable def vacChar (W : VertexOperatorAlgebra 𝒱) : UpperHalfPlane → ℂ :=
  qEvalAt (gradedDim (adjointModule W)) (0 - W.centralCharge / 24)

def ZhuVacuumModularity (W : VertexOperatorAlgebra 𝒱) : Prop :=
  AdmissiblyRational W → C2Cofinite W → ∃ (n : ℕ) (D : ModularData W (Fin n)) (i : Fin n),
    ZhuModularityAnalytic D ∧ charOf D i = vacChar W

theorem not_admissiblyRational_tripletW2_of_logS
    (hzhu : ZhuVacuumModularity tripletW2) (hC2 : C2Cofinite tripletW2)
    (hlog : HasLogTerm (fun τ => vacChar tripletW2 (ModularGroup.S • τ))) :
    ¬ AdmissiblyRational tripletW2 :=

def AllModulesOrdinary (W : VertexOperatorAlgebra 𝒱) : Prop :=
  ∀ {M : Type u} [AddCommGroup M] [Module ℂ M] (N : LogVModule W M), N.IsL0Semisimple

def RationalImpliesOrdinary (W : VertexOperatorAlgebra 𝒱) : Prop :=
  AdmissiblyRational W → AllModulesOrdinary W

theorem not_admissiblyRational_tripletW2_of_logModule
    (hro : RationalImpliesOrdinary tripletW2) (h : ℂ)
    (hAnm : ∀ n : ℕ, FiniteDimensional ℂ (AnmFull tripletW2 n 0))
    (hA : FiniteDimensional ℂ (A tripletW2)) (hgate : DJAnnEqAn tripletW2 0)
    (hnd : OmegaJordanAt tripletW2 h) :
    ¬ AdmissiblyRational tripletW2 :=
\end{leancode}
\begin{dictionary}
\lean{qEvalAt χ h τ} $=q^h\sum_{n\ge0}\chi(h+n)q^n$ (Chapter~\ref{ch:L14}); with $h=-c/24=\frac1{12}$ and the graded dimension
of the adjoint module, \lean{vacChar tripletW2} is $\chi_0$ of \eqref{l18:eq:char} as a function on $\mathbb H$
(\lean{w2VacChar_low} checks its first two coefficients $1,0$). \lean{ModularData}, \lean{ZhuModularityAnalytic} and
\lean{charOf} are Chapter~\ref{ch:L15}, Definition~\ref{l15:def:modulardata}: a finite family of \emph{ordinary}
irreducible modules whose characters transform under $SL_2(\Z)$ by a matrix. Since every \lean{charOf D i} is a
$T$-eigenfunction (\lean{charOf_isTEigen}), Zhu's law puts $\chi_0(S\tau)$ in the log-free span, contradicting
\lean{hlog}. The three hypotheses of the first theorem are all open in the library: \lean{hzhu} would follow from the
rational headline of Chapter~\ref{ch:L15} (\lean{zhuModularity_canonical_of_C2Cofinite_admissiblyRational}) once the Abe
model is proved simple, so that its vacuum character is one member of the canonical family; neither the simplicity nor
the composition is in the library. \lean{hC2} is
Theorem~\ref{l18:thm:abeC2}, and \lean{hlog} is the explicit transformation \eqref{l18:eq:Schi}, which needs the
product formula for the character (the library has no basis of the symplectic-fermion Fock space).
\lean{IsL0Semisimple} means \lean{L0nil = 0}.
\end{dictionary}
\begin{physics}
The physicist's argument ``rational theories have no $\log q$ in their modular transformations'', formalised:
in a rational theory the torus partition function lies in the span of characters of ordinary modules; the $c=-2$
vacuum character's $S$-transform contains $-\frac{i\tau}2\eta(\tau)^2$, which is not of that form.
\end{physics}
\begin{status}
Proved, \textbf{conditional on} \lean{hzhu}, \lean{hC2}, \lean{hlog} (resp.\ \lean{hro} and the four inputs of
Theorem~\ref{l18:thm:logmodule}); none is discharged. Standard axioms.
\end{status}

%======================================================================
\section{Pseudo-traces and logarithmic \texorpdfstring{$q$}{q}-series}\label{l18:sec:pseudo}
%======================================================================

In a rational theory the torus one-point functions are traces $\mathrm{tr}_M\,o(v)q^{L_0-c/24}$ over the irreducible
modules $M$. In a logarithmic theory two things break: $L_0$ is not diagonalisable, so $q^{L_0}$ has polynomial
dependence on $\tau$; and traces over irreducibles no longer span the space of torus amplitudes. Miyamoto's
answer (2004) is to take, for a projective module $P$ with a large commutant algebra $A$, a symmetric linear
functional $\varphi$ on $A$ and form the \emph{pseudo-trace} $\mathrm{tr}^\varphi_P$. The library formalises the
algebraic core in its most general form (Arike's): the Hattori--Stallings trace.

\subsection{Symmetric functionals and coordinate systems}

\begin{definition}[Symmetric functionals; commutator span]\label{l18:def:slf}
For a $K$-algebra $A$, $\mathrm{SLF}(A)=\{\varphi:A\to K\text{ linear}\mid\varphi(ab)=\varphi(ba)\}$ and
$[A,A]=\mathrm{span}_K\{ab-ba\}$. A symmetric functional is the same as a functional on $A/[A,A]$.
\end{definition}
\noindent File \leanpath{VOA/\allowbreak Modularity/\allowbreak SemisimpleTrace.lean} (assigned to Chapter~\ref{ch:L15};
quoted here for self-containment).
\begin{leancode}
def commSpan : Submodule K A := Submodule.span K {x : A | ∃ a b : A, x = a * b - b * a}

def symmetricFunctionals : Submodule K (A →ₗ[K] K) where
  carrier := {φ : A →ₗ[K] K | ∀ a b : A, φ (a * b) = φ (b * a)}
\end{leancode}
\begin{dictionary}
Literally the textbook definitions. \lean{symmetricLift φ} (same file) is the induced map $A/[A,A]\to K$.
\end{dictionary}
\begin{physics}
Symmetric functionals on Zhu's algebra are the ``torus amplitudes at the level of ground states''; for semisimple
$A$ they are combinations of traces of simple modules, and their number is the number of irreducibles.
\end{physics}
\begin{status}
Definitions.
\end{status}

\begin{definition}[Finite coordinate system; raw trace]\label{l18:def:fincoords}
A \emph{finite coordinate system} on a left $A$-module $W$ is a family $(u_i,\alpha_i)_{i<n}$ with $u_i\in W$,
$\alpha_i\in\mathrm{Hom}_A(W,A)$ and $\sum_i\alpha_i(w)u_i=w$ for all $w$. For $f\in\End_A(W)$ its \emph{raw trace} is
$\sum_i\alpha_i(f(u_i))\in A$.
\end{definition}
\noindent File \leanpath{VOA/\allowbreak Modularity/\allowbreak PseudoTrace.lean}.
\begin{leancode}
structure FinCoords (A : Type*) [Ring A] (W : Type*) [AddCommGroup W] [Module A W] where
  /-- The number of generators. -/
  card : ℕ
  /-- The generators. -/
  gen : Fin card → W
  /-- The coordinate functionals. -/
  coord : Fin card → (W →ₗ[A] A)
  /-- The reproducing identity. -/
  total : ∀ w : W, ∑ i, coord i w • gen i = w

theorem nonempty_of_finite_projective (A : Type*) [Ring A] (W : Type*) [AddCommGroup W]
    [Module A W] [Module.Finite A W] [Module.Projective A W] : Nonempty (FinCoords A W) := by

def rawTrace (c : FinCoords A W) (f : Module.End A W) : A :=
  ∑ i, c.coord i (f (c.gen i))
\end{leancode}
\begin{dictionary}
A finite coordinate system is the dual-basis lemma: it exists exactly for finitely generated projective modules
(\lean{nonempty_of_finite_projective}; conversely \lean{FinCoords.finite}, \lean{FinCoords.projective}). This
replaces the tensor product $W\otimes_A\mathrm{Hom}_A(W,A)$, which Mathlib lacks over noncommutative rings. Mathlib's
modules are left modules; Miyamoto and Arike use right modules (a harmless side convention).
\end{dictionary}
\begin{physics}
A projective module over the commutant is the ``multiplicity space'' of a projective (logarithmic) VOA-module;
the coordinate system plays the role of a basis in which the trace is taken.
\end{physics}
\begin{status}
Definition and theorem, standard axioms.
\end{status}

\subsection{The Hattori--Stallings trace and the pseudo-trace}

\begin{theorem}[Hattori--Stallings trace]\label{l18:thm:HS}
For a finitely generated projective $A$-module $W$ the class of the raw trace in $A/[A,A]$ is independent of the
coordinate system, defining a $K$-linear $T_W:\End_A(W)\to A/[A,A]$, and $T_W(f\circ g)=T_X(g\circ f)$ for
$g:W\to X$, $f:X\to W$.
\end{theorem}
\noindent File \leanpath{VOA/\allowbreak Modularity/\allowbreak PseudoTrace.lean}.
\begin{leancode}
noncomputable def hattoriTrace : Module.End A W →ₗ[K] (A ⧸ commSpan K A) :=

theorem hattoriTrace_apply_of_coords (c : FinCoords A W) (f : Module.End A W) :
    hattoriTrace K A W f = (commSpan K A).mkQ (c.rawTrace f) := by

theorem hattoriTrace_comp_comm {X : Type*} [AddCommGroup X] [Module A X] [Module K X]
    [IsScalarTower K A X] [Module.Finite A X] [Module.Projective A X]
    (f : X →ₗ[A] W) (g : W →ₗ[A] X) :
    hattoriTrace K A W (f ∘ₗ g) = hattoriTrace K A X (g ∘ₗ f) := by
\end{leancode}
\begin{dictionary}
The section variables (not shown) are a field $K$, a $K$-algebra $A$, and $W$ a finitely generated projective left
$A$-module with compatible $K$-structure. \lean{(commSpan K A).mkQ} is the quotient map to $A/[A,A]$. The theorem
\lean{hattoriTrace_apply_of_coords} says the trace can be computed in \emph{any} coordinate system.
\end{dictionary}
\begin{proof}[Proof on paper]
For coordinate systems $(u_i,\alpha_i)$ on $W$ and $(v_j,\beta_j)$ on $X$:
$\sum_i\alpha_i(fg u_i)=\sum_{i,j}\alpha_i(f(\beta_j(gu_i)v_j))=\sum_{i,j}\beta_j(gu_i)\alpha_i(fv_j)$ and
$\sum_j\beta_j(gfv_j)=\sum_{i,j}\alpha_i(fv_j)\beta_j(gu_i)$; the two sums differ by commutators. Taking $X=W$, $g=\mathrm{id}$
gives independence; one coordinate system gives cyclicity.
\end{proof}
\begin{physics}
Cyclicity of the trace is what makes torus amplitudes well defined (the torus has no preferred starting point
for the time evolution); the Hattori--Stallings trace is the most general cyclic trace on endomorphisms of
projective modules.
\end{physics}
\begin{status}
Proved, unconditional; standard axioms.
\end{status}

\begin{theorem}[Pseudo-trace; it is the trace; it realises every symmetric functional]\label{l18:thm:pseudo}
For $\varphi\in\mathrm{SLF}(A)$ define $\mathrm{tr}^\varphi_W=\bar\varphi\circ T_W:\End_A(W)\to K$. Then $\mathrm{tr}^\varphi_W$ is a symmetric
functional on $\End_A(W)$; for $A=K$, $\varphi=\mathrm{id}$, it is the ordinary trace; and on the regular module,
$\mathrm{tr}^\varphi_A(x\mapsto xa)=\varphi(a)$, so $\varphi\mapsto\mathrm{tr}^\varphi_A$ is injective.
\end{theorem}
\noindent File \leanpath{VOA/\allowbreak Modularity/\allowbreak PseudoTrace.lean}.
\begin{leancode}
noncomputable def pseudoTrace (φ : symmetricFunctionals K A) : Module.End A W →ₗ[K] K :=

theorem pseudoTrace_mem_symmetricFunctionals (φ : symmetricFunctionals K A) :
    pseudoTrace W φ ∈ symmetricFunctionals K (Module.End A W) := by

theorem pseudoTrace_idSymmetric_eq_trace (f : Module.End K W) :
    pseudoTrace W (idSymmetric K) f = LinearMap.trace K W f := by

theorem pseudoTrace_rightMulEnd (φ : symmetricFunctionals K A) (a : A) :
    pseudoTrace A φ (rightMulEnd a) = (φ : A →ₗ[K] K) a := by

theorem pseudoTraceMap_regular_injective :
    Function.Injective (pseudoTraceMap (K := K) (A := A) A) := by
\end{leancode}
\begin{dictionary}
\lean{idSymmetric K} is the identity functional on $K$; \lean{LinearMap.trace K W f} is Mathlib's trace of an
endomorphism of a finite-dimensional space; \lean{rightMulEnd a} is $x\mapsto xa$ (an $A$-linear map of the left
regular module). The second statement is the decisive validation: on vector spaces the definition reduces to the
trace, so it is not some other functional. This is Arike's definition (Arike 2010, \S2), shown by him to agree with
Miyamoto's basis-dependent Definition~3.6; that comparison is \emph{not} formalised (not needed).
\end{dictionary}
\begin{proof}[Proof on paper]
Symmetry is cyclicity of $T_W$. For $A=K$ any basis is a coordinate system and the raw trace is the sum of diagonal
entries. For the regular module take the coordinate system $(1,\mathrm{id})$: the raw trace of $x\mapsto xa$ is $a$.
\end{proof}
\begin{physics}
Miyamoto's theorem (not formalised) says that for $C_2$-cofinite $V$ the torus one-point functions are spanned by
pseudo-traces $\mathrm{tr}^\varphi_P\,o(v)q^{L_0-c/24}$ over projective modules; the regular-module statement here is the
algebraic fact that every symmetric functional on the commutant is realised.
\end{physics}
\begin{status}
Proved, unconditional; \texttt{\#print axioms VOA.pseudoTrace\_idSymmetric\_eq\_trace},
\lean{pseudoTrace_mem_symmetricFunctionals}: standard three.
\end{status}

\subsection{Polynomial-valued \texorpdfstring{$q$}{q}-series}

\begin{definition}[Logarithmic $q$-series; the pseudo-trace exponential polynomial]\label{l18:def:logq}
A \emph{logarithmic formal $q$-series} is a function from exponents $e\in\C$ to polynomials in $\tau$. For a finite
projective $A$-module $P$, endomorphisms $S,N$ and a cutoff $r$,
\begin{equation*}
\mathcal P_\varphi(S,N;r)(\tau)=\sum_{k<r}\frac{(2\pi i\tau)^k}{k!}\,\mathrm{tr}^\varphi_P(S N^k),
\end{equation*}
which for $N$ nilpotent with $N^r=0$ is $\mathrm{tr}^\varphi_P\bigl(S\,e^{2\pi i\tau N}\bigr)$.
\end{definition}
\noindent File \leanpath{VOA/\allowbreak Modularity/\allowbreak LogarithmicQSeries.lean}.
\begin{leancode}
abbrev LogFormalQSeries := ℂ → Polynomial ℂ

noncomputable def logExpCoeff (k : ℕ) : ℂ :=
  (2 * (Real.pi : ℂ) * Complex.I) ^ k / (k.factorial : ℂ)

noncomputable def pseudoTraceExpPolynomial (phi : symmetricFunctionals ℂ A)
    (S N : Module.End A P) (r : ℕ) : Polynomial ℂ :=
  ∑ k ∈ Finset.range r,
    Polynomial.monomial k (logExpCoeff k * pseudoTrace P phi (S * N ^ k))

theorem pseudoTraceExpPolynomial_zero_eq_trace (S : Module.End ℂ P) :
    pseudoTraceExpPolynomial P (idSymmetric ℂ) S 0 1 =
      Polynomial.C (LinearMap.trace ℂ P S) := by
\end{leancode}
\begin{dictionary}
\lean{Polynomial.monomial k c} is $c\,X^k$, with $X$ standing for $\tau$. The value at exponent $e$ of a
\lean{LogFormalQSeries} is the coefficient polynomial of $q^e$, so the series is $\sum_eP_e(\tau)q^e$.
\lean{pseudoTraceExpPolynomial_stable} shows that once $N^r=0$ the cutoff does not matter, and the last theorem
that with $N=0$, $A=\C$ one recovers the ordinary trace. The file states explicitly that it asserts no VOA origin,
no recursion, no convergence and no modularity.
\end{dictionary}
\begin{physics}
On a generalized $L_0$-eigenspace of eigenvalue $h$, $q^{L_0}=q^h\,e^{2\pi i\tau L_0^{\mathrm{nil}}}$; the pseudo-trace of
$o(v)q^{L_0}$ is therefore $q^h$ times a polynomial in $\tau$. These polynomial coefficients are the $\tau$-linear
(and higher) terms of the logarithmic torus amplitudes --- the same terms that appear in \eqref{l18:eq:Schi}.
\end{physics}
\begin{status}
Proved; standard axioms.
\end{status}

\subsection{The genus-one layer, on hypothetical commutant data}

\begin{definition}[Commutant data on a logarithmic module]\label{l18:def:commutant}
For a logarithmic module $N$ of $V$ and a $\C$-algebra $A$, \emph{commutant data} consist of: $A$
finite-dimensional; an $A$-module structure on every generalized weight space $N_{[\alpha]}$, finitely generated and
projective; a symmetric functional $\varphi$ on $A$; the requirement that every vertex-operator mode intertwines these
$A$-actions; and an $A$-linear lift of each zero mode $o(v)$ to $N_{[\alpha]}$. From such data one forms the
one-point function $S(v)=\sum_e\mathcal P_\varphi(o(v),L_0^{\mathrm{nil}};r_\alpha)\,q^e$, $\alpha=e+c/24$.
\end{definition}
\noindent File \leanpath{VOA/\allowbreak Modularity/\allowbreak LogarithmicPseudoTraceRealization.lean}.
\begin{leancode}
structure LogPseudoTraceCommutantData where
  /-- The auxiliary algebra is finite-dimensional over `ℂ`. -/
  finite_algebra : FiniteDimensional ℂ A
  /-- Every generalized weight piece is finitely generated over the auxiliary algebra. -/
  finite_piece : ∀ α : ℂ, Module.Finite A (N.weight α)
  /-- Every generalized weight piece is projective over the auxiliary algebra. -/
  projective_piece : ∀ α : ℂ, Module.Projective A (N.weight α)
  /-- The symmetric functional used by the Hattori--Stallings pseudo-trace. -/
  phi : symmetricFunctionals ℂ A
  /-- The auxiliary action commutes with every actual vertex-operator mode between generalized
  weight pieces. -/
  mode_commutes : ∀ (a : A) (α β : ℂ) (v : 𝒱) (n : ℤ) (x : N.weight α)
    (hx : mode (N.YM v) n (x : M) ∈ N.weight β),
    mode (N.YM v) n ((a • x : N.weight α) : M) =
      ((a • (⟨mode (N.YM v) n (x : M), hx⟩ : N.weight β) : N.weight β) : M)
  /-- The genuine VOA zero mode, lifted to an `A`-linear endomorphism of every generalized
  weight piece and linear in the inserted state. -/
  zeroModeA : ∀ α : ℂ, 𝒱 →ₗ[ℂ] Module.End A (N.weight α)
  /-- The lift has the actual VOA zero mode as its underlying map. -/
  zeroModeA_coe : ∀ (α : ℂ) (v : 𝒱) (x : N.weight α),
    ((zeroModeA α v x : N.weight α) : M) = logOMod N v (x : M)
\end{leancode}
\begin{leancode}
noncomputable def logPseudoTraceOnePoint (D : LogPseudoTraceCommutantData N A) (v : 𝒱) :
    LogFormalQSeries :=
  fun e =>
    let α := e + W.centralCharge / 24
    letI : Module.Finite A (N.weight α) := D.finite_piece α
    letI : Module.Projective A (N.weight α) := D.projective_piece α
    pseudoTraceExpPolynomial (N.weight α) D.phi (D.zeroModeA α v)
      (D.nilpotentPiece α) (D.nilpotentCutoff α)
\end{leancode}
\begin{dictionary}
\lean{N.weight α} is the generalized $L_0$-eigenspace of a \lean{LogVModule} (Definition~\ref{l12:def:log});
\lean{logOMod N v} is the zero mode $o(v)$ (the homogeneous component $v_{(\wt v-1)}$, extended linearly).
\lean{nilpotentPiece D α} is \emph{derived}: it is $o(\omega)-\alpha$, i.e.\ $L_0^{\mathrm{nil}}$, lifted; \lean{nilpotentCutoff}
is its nilpotency exponent. The structure has \emph{no instance in the library} (it is used only in the three
\leanpath{LogarithmicPseudoTrace*.lean} files). It is not unsatisfiable: for any logarithmic module with
finite-dimensional generalized weight spaces one can take $A=\C$, $\varphi=\mathrm{id}$ (every field is then
automatic, since scalar multiplication commutes with linear maps), and $S(v)$ becomes the ordinary trace of
$o(v)e^{2\pi i\tau L_0^{\mathrm{nil}}}$ (paper argument; not formalised). The interesting instances, with $A$ the
endomorphism algebra of a projective generator, come from the Huang--Miyamoto existence theorem, which is not in the
library.
\end{dictionary}
\begin{physics}
$A$ is the algebra of ``hidden symmetries'' of a projective module that commute with the chiral algebra (for
$W(2)$ on the projective cover of the vacuum, a two-dimensional algebra $\C[\varepsilon]/\varepsilon^2$); pseudo-traces
twisted by $\varphi$ are the new torus amplitudes of the logarithmic theory.
\end{physics}
\begin{status}
Definitions; standard axioms. \textbf{No instance}, hence no concrete application.
\end{status}

\begin{theorem}[The cyclic block relations; the recursion bridge]\label{l18:thm:logblock}
For any commutant data, the one-point function $S$ is linear in $v$, its vacuum value is the pseudo-character,
$S(L_{[-1]}a)=0$, and $S(a[0]b)=0$ (Zhu's zero-mode relation, with square-bracket modes). If moreover the logarithmic
$[-1]$-recursion \lean{LogPseudoTraceRecursionData} holds, then $S$ also satisfies the $[-2]$ relation and the
logarithmic $\omega$-recursion with $q\frac d{dq}$ acting on polynomial coefficients as
$P_e\mapsto eP_e+(2\pi i)^{-1}P'_e$; i.e.\ $S$ is a logarithmic formal genus-one block.
\end{theorem}
\noindent Files \leanpath{VOA/\allowbreak Modularity/\allowbreak LogarithmicPseudoTraceBlock.lean},
\leanpath{LogarithmicPseudoTraceRecursion.lean}.
\begin{leancode}
structure IsLogPseudoTraceCyclicBlock
    (S : 𝒱 →ₗ[ℂ] LogFormalQSeries) (character : LogFormalQSeries) : Prop where
  vacuum : S W.vacuum = character
  translation : ∀ a : 𝒱, S (Lsq W (-1) a) = 0
  zhuZero : ∀ (k : ℕ), 1 ≤ k → ∀ a ∈ W.grading (k : ℤ), ∀ b : 𝒱,
    S (vSqZ W a k 0 b) = 0
\end{leancode}
\begin{leancode}
noncomputable def logQDeriv (χ : LogFormalQSeries) : LogFormalQSeries :=
  fun e => e • χ e + (logDerivativeScale)⁻¹ • (χ e).derivative

theorem logPseudoTraceTwoPoint_conformalVector
    (D : LogPseudoTraceCommutantData N A) (v : 𝒱) :
    D.logPseudoTraceTwoPoint W.conformalVector v =
      logQDeriv (D.logPseudoTraceOnePoint v) +
        (W.centralCharge / 24) • D.logPseudoTraceOnePoint v := by

structure LogPseudoTraceRecursionData (D : LogPseudoTraceCommutantData N A) : Prop where
  zhu : ∀ (ka : ℕ), 1 ≤ ka → ∀ a ∈ W.grading (ka : ℤ), ∀ v : 𝒱,
    D.logPseudoTraceOnePointHom (vSqZ W a ka (-1) v) =
      D.logPseudoTraceTwoPoint a v +
        ∑ᶠ m : ℕ, logQConv (zhuG (m + 1))
          (D.logPseudoTraceOnePointHom
            (vSqZ W a ka (2 * (m + 1) - 1 : ℤ) v))
\end{leancode}
\begin{dictionary}
\lean{Lsq W n} and \lean{vSqZ W a k n b} are the square-bracket Virasoro modes $L_{[n]}$ and the square-bracket
products $a[n]b$ of Chapter~\ref{ch:L15}, Definition~\ref{l15:def:sqmode}; \lean{zhuG (m+1)} are Zhu's Eisenstein
series $G_{2m+2}$ and \lean{logQConv} their convolution with $q$-coefficients (Definition~\ref{l15:def:zhuG}).
\lean{logDerivativeScale} $=2\pi i$. \lean{logPseudoTraceTwoPoint a v} is the pseudo-trace of $o(a)o(v)q^{L_0}$. The
$\omega$-insertion formula is the logarithmic version of $\mathrm{tr}\,o(\omega)o(v)q^{L_0-c/24}=(q\frac d{dq}+\frac c{24})\mathrm{tr}\,o(v)q^{L_0-c/24}$:
the extra $(2\pi i)^{-1}P'$ term is what differentiating $e^{2\pi i\tau L_0^{\mathrm{nil}}}$ produces. The recursion data
is a hypothesis (one cross-weight identity), not proved for any $D$; the conclusion
\lean{IsLogPseudoTraceFormalBlock} is then a theorem (\lean{isLogPseudoTraceFormalBlock}).
\end{dictionary}
\begin{proof}[Proof on paper]
Translation: $o(L_{[-1]}a)=0$ is already true as an operator. Zero-mode relation: $o(a[0]b)=[o(a),o(b)]$ up to the
$L_0$-shift, and pseudo-traces kill commutators, also when multiplied by powers of $L_0^{\mathrm{nil}}$ which commute with
zero modes (\lean{pseudoTrace_commutator_mul_pow_eq_zero}). The $\omega$-formula: $o(\omega)=L_0=\alpha+L_0^{\mathrm{nil}}$ on $N_{[\alpha]}$,
and $\frac d{d\tau}\sum_k\frac{(2\pi i\tau)^k}{k!}\mathrm{tr}(SN^k)=2\pi i\sum_k\frac{(2\pi i\tau)^{k}}{k!}\mathrm{tr}(SN^{k+1})$.
\end{proof}
\begin{physics}
These are the genus-one Ward identities for logarithmic torus amplitudes: invariance under translations on the
torus, the zero-mode (commutator) relation, and the modular heat equation $\langle T\rangle_\tau\propto\partial_\tau$ --- with the
twist that $\partial_\tau$ now also hits the explicit $\tau$-dependence of the pseudo-trace.
\end{physics}
\begin{status}
Proved; the formal-block theorem conditional on \lean{LogPseudoTraceRecursionData}. \texttt{\#print axioms
VOA.LogPseudoTraceCommutantData.isLogPseudoTraceCyclicBlock}, \lean{isLogPseudoTraceFormalBlock}: standard three.
Applies to no concrete VOA (no instance of the commutant data). No convergence, spanning or modularity is
claimed.
\end{status}

%======================================================================
\section{The doublet algebra}\label{l18:sec:doublet}
%======================================================================

\subsection{The generators on an abstract lattice}

\begin{theorem}[Doublet generators]\label{l18:thm:doubletgens}
Let $L$ be positive definite, $g\in L$ with $2\langle g,g\rangle=p$ (so $\gamma=2g$ has $\langle\gamma,\gamma\rangle=2p$), and
$\langle\lambda,g\rangle=\frac{p-1}2$. Put $Q=(e^{2g})_{(0)}$, $a^-=e^{-g}$, $a^+=Qa^-$. Then $a^\pm$ have deformed weight
$\frac{3p-2}4$, $a^+\ne0$, $Qa^+=0$, and $e^{+g}$ has weight $\frac{2-p}4$, negative for $p>2$.
\end{theorem}
\noindent File \leanpath{VOA/\allowbreak Examples/\allowbreak DoubletGenerators.lean}.
\begin{leancode}
theorem LBG_zero_aMinus (hPos : L.PosDef) (lam : Fin (rank L) → ℂ) (g : L.carrier) (p : ℕ)
    (hgg : 2 * L.form g g = (p : ℤ)) (hlam : bgPair L lam g = ((p : ℂ) - 1) / 2) :
    LBG L lam 0 (aMinus L g) = ((3 * (p : ℂ) - 2) / 4) • aMinus L g := by

theorem LBG_zero_aPlus (hPos : L.PosDef) (lam : Fin (rank L) → ℂ) (g : L.carrier) (p : ℕ)
    (hgg : 2 * L.form g g = (p : ℤ)) (hlam : bgPair L lam g = ((p : ℂ) - 1) / 2) :
    LBG L lam 0 (aPlus L g) = ((3 * (p : ℂ) - 2) / 4) • aPlus L g := by

theorem aPlus_ne_zero (hPos : L.PosDef) (g : L.carrier) (p : ℕ) (hp : 1 ≤ p)
    (hgg : 2 * L.form g g = (p : ℤ)) : aPlus L g ≠ 0 := by

theorem Qdb_aPlus (g : L.carrier) (p : ℕ) (hp : 1 ≤ p) (hgg : 2 * L.form g g = (p : ℤ)) :
    Qdb L g (aPlus L g) = 0 := by

theorem LBG_zero_eG (hPos : L.PosDef) (lam : Fin (rank L) → ℂ) (g : L.carrier) (p : ℕ)
    (hgg : 2 * L.form g g = (p : ℤ)) (hlam : bgPair L lam g = ((p : ℂ) - 1) / 2) :
    LBG L lam 0 (latChargeGen L g) = ((2 - (p : ℂ)) / 4) • latChargeGen L g := by
\end{leancode}
\begin{dictionary}
\lean{Qdb L g} $=$ \lean{Screen.Qscr L (g + g)}, \lean{aMinus L g} $=e^{-g}$, \lean{aPlus L g} $=Q\,e^{-g}$. The lattice is
abstract: these are weight computations, true in any lattice satisfying the two hypotheses. Since the Gram form is
integer-valued, $2\langle g,g\rangle=p$ forces $p$ even (\lean{p_even}). Weights: $\frac12\cdot\frac p2+\frac{p-1}2=\frac{3p-2}4$ for $e^{-g}$,
$\frac p4-\frac{p-1}2=\frac{2-p}4$ for $e^{g}$. $a^+=\varepsilon\,S_{p-1}(2g)\otimes e^{g}$ (\lean{aPlus_eq}), nonzero by
\lean{vSchur_ne_zero}.
\end{dictionary}
\begin{physics}
The doublet fields $a^\pm$ of Adamovi\'c--Milas: the ``square roots'' of the triplet, with half the charge.
Their OPE $a^+(z)a^-(w)$ produces the triplet; the negative-weight state $e^{g}$ shows that the ambient half-lattice
theory is not of CFT type, so the doublet must be cut out carefully.
\end{physics}
\begin{status}
Proved, unconditional; standard axioms.
\end{status}

\subsection{The half lattice}

\begin{definition}[The half lattice $\frac12L_p$]\label{l18:def:hlat}
\lean{hlatP p} is the rank-one lattice $\Z g$ with $\langle mg,ng\rangle=\frac p2mn$ (natural-number division), and
$q_H(m)=\frac{pm^2}4-\frac{(p-1)m}2$ is the deformed weight of $e^{mg}$.
\end{definition}
\noindent File \leanpath{VOA/\allowbreak Examples/\allowbreak DoubletHalfLattice.lean}.
\begin{leancode}
noncomputable def hlatP (p : ℕ) : IntegralLattice where
  carrier := ℤ
  form := ((p / 2 : ℕ) : ℤ) • LinearMap.mul ℤ ℤ

noncomputable def qH (p : ℕ) (m : ℤ) : ℚ :=
  ((p : ℚ) * (m : ℚ) ^ 2) / 4 - (((p : ℚ) - 1) * (m : ℚ)) / 2

theorem hlatP_even_iff (p k : ℕ) (hp2 : 2 * k = p) : (hlatP p).Even ↔ 4 ∣ p := by

theorem not_posQuad_qrH (p k : ℕ) (hp2 : 2 * k = p) (hp : 2 < p) :
    ¬ PosQuad (hlatP p) (qrH p k hp2) := by

theorem qH_two_mul (p : ℕ) (m : ℤ) : qH p (2 * m) = ((Triplet.qZP p m : ℤ) : ℚ) := by
\end{leancode}
\begin{dictionary}
\lean{p / 2} is natural-number division, so for odd $p$ the form is $\lfloor p/2\rfloor mn$; every theorem that matters
carries \lean{hp2 : 2 * k = p}, so only even $p$ is used. \lean{qH_two_mul} shows the doubling $m\mapsto2m$ embeds the
triplet lattice $L_p$ isometrically with $q_H(2m)=q_p(m)$, as it must. The half lattice is even iff $4\mid p$; for $p\equiv2$
mod $4$ it is odd and its lattice vertex algebra is a super algebra, which the library cannot build from the bosonic
\lean{YfullGen}. \lean{not_posQuad_qrH}: the refinement is not positive for $p>2$ ($q_H(1)<0$), but it is bounded below
(\lean{qH_bddBelow}).
\end{dictionary}
\begin{physics}
Halving the compactification radius doubles the charge lattice; the extra states $e^{\pm g}$ are the doublet's
building blocks. For $4\mid p$ they have half-integral weight but bosonic statistics (a spin--statistics
violation), for $p\equiv2$ mod $4$ integral weight but fermionic statistics.
\end{physics}
\begin{status}
Proved; standard axioms.
\end{status}

\subsection{The doublet as a generated super subalgebra}

\begin{definition}[The doublet $\leighteenA(p)$]\label{l18:def:doublet}
For a super near-VOA $W$ and states $a^-,a^+$, \lean{doubletA W am ap} is the super subalgebra generated by
$\{\omega,a^-,a^+\}$. A \emph{deformed ambient} for $(L,\lambda)$ is a super near-VOA on $V_L$ whose $L_0$ is $L^\lambda_0$.
The concrete doublet \lean{doubletApC p h4 hp} is \lean{doubletA} in the concrete ambient \lean{ambientC p h4 hp} on the
half lattice, for $4\mid p$, with $a^\pm$ as above.
\end{definition}
\noindent Files \leanpath{VOA/\allowbreak Examples/\allowbreak DoubletAp.lean},
\leanpath{DoubletApConcrete.lean}.
\begin{leancode}
def doubletA (am ap : 𝒱) : SuperVOASubalgebra W :=
  SuperVOASubalgebra.generated W {W.conformalVector, am, ap}

structure DeformedAmbient (lam : Fin (rank L) → ℂ) where
  /-- The ambient super near-VOA on the half-lattice module. -/
  W : SuperNearVOA (LatModuleGen L)
  /-- Its Virasoro zero mode is the background-charge deformed one. -/
  L0_eq : W.L 0 = LBG L lam 0
\end{leancode}
\begin{leancode}
noncomputable def ambientC (p : ℕ) (h4 : 4 ∣ p) (hp : 0 < p) :
    SuperNearVOA (LatModuleGen (latD p)) :=

noncomputable def doubletApC (p : ℕ) (h4 : 4 ∣ p) (hp : 0 < p) :
    SuperVOASubalgebra (ambientC p h4 hp) :=
  DoubletAp.doubletAp (deformedAmbientC p h4 hp) (gD p)
\end{leancode}
\begin{dictionary}
\lean{SuperNearVOA} (Chapter~\ref{ch:L17}, Definition~\ref{l17:def:snvoa}) is a super vertex algebra with a conformal
vector and a $\C$-grading but \emph{without} positivity (bounded-belowness is a separate predicate
\lean{BoundedBelow}); this is forced here because the ambient has negative weights. \lean{latD p} $=$
\lean{hlatP p}, \lean{gD p} its generator; \lean{ambientC} is built from the general super ambient of
\leanpath{LatticeSuperAmbientB.lean} (Chapter~\ref{ch:L17}, Section~\ref{l17:sec:lattice}) with background charge
\lean{lamD p} $=\frac{p-1}p\,g$. \lean{DoubletAp.doubletAp da g} is \lean{doubletA} with $a^\pm$ of
Theorem~\ref{l18:thm:doubletgens}. In the abstract file \leanpath{DoubletAp.lean} nothing is instantiated; the concrete
file closes that gap for $4\mid p$.
\end{dictionary}
\begin{physics}
$\leighteenA(p)$ is the chiral algebra generated by $T$ and the doublet $a^\pm$; its charge-even part is $W(p)$. In
the language of extensions, $\leighteenA(p)$ is a simple-current extension of $W(p)$ by its doublet module.
\end{physics}
\begin{status}
Definitions; standard axioms.
\end{status}

\begin{theorem}[What is proved about the concrete doublet]\label{l18:thm:doublet}
For $4\mid p$: $\leighteenA(p)$ is bounded below; its ambient is bounded below but not positive (it has the nonzero state $e^{g}$ of negative weight);
its central charge is $c_{1,p}$, that of $W(p)$; its parity is trivial, so $\leighteenA(p)_{\bar0}=\leighteenA(p)$; and the
integrality hypothesis needed to form the even part as a bosonic VOA is false.
\end{theorem}
\noindent File \leanpath{VOA/\allowbreak Examples/\allowbreak DoubletApConcrete.lean}.
\begin{leancode}
theorem doubletApC_boundedBelow (p : ℕ) (h4 : 4 ∣ p) (hp : 0 < p) :
    (doubletApC p h4 hp).toSuperNearVOA.BoundedBelow :=

theorem ambientC_not_positivity (p : ℕ) (h4 : 4 ∣ p) (hp : 0 < p) :
    ¬ (∀ h : ℂ, h.re < 0 → (ambientC p h4 hp).grading h = ⊥) :=

theorem doubletApC_centralCharge_eq_tripletW (p : ℕ) (h4 : 4 ∣ p) (hp : 0 < p) :
    (doubletApC p h4 hp).toSuperNearVOA.centralCharge = (Wp.tripletWVOA p hp).centralCharge := by

theorem doubletApC_parity_zero (p : ℕ) (h4 : 4 ∣ p) (hp : 0 < p) :
    (doubletApC p h4 hp).toSuperNearVOA.parity 0 = ⊤ := by

theorem hint_false_of_four_dvd (p : ℕ) (h4 : 4 ∣ p) (hp : 0 < p) :
    ¬ (∀ h : ℂ, (¬ ∃ n : ℤ, (n : ℂ) = h) →
        ∀ x ∈ (doubletApC p h4 hp).toSuperNearVOA.parity 0,
          x ∈ (doubletApC p h4 hp).toSuperNearVOA.grading h → x = 0) := by
\end{leancode}
\begin{dictionary}
\lean{hp : 0 < p} is passed as \lean{1 ≤ p} to \lean{tripletWVOA}; the central-charge comparison is again a
comparison of numbers. The parity statement shows the super structure of the concrete ambient is trivial (everything
even, because for $4\mid p$ the half lattice is even and the ambient is a bosonic vertex algebra with half-integral
weights). So the even-part construction \lean{doubletAEvenVOA} of \leanpath{DoubletAp.lean} --- which would be the
natural candidate for $W(p)$ inside $\leighteenA(p)$ --- cannot be applied, and \lean{hint_false_of_four_dvd} proves it.
In the literature the relevant $\Z/2$ is the \emph{charge parity} $e^{mg}\mapsto(-1)^me^{mg}$, not the super parity;
the library does not construct it.
\end{dictionary}
\begin{warning}
\textbf{What is not proved about the doublet} (all recorded in the source): positivity of $\leighteenA(p)$ (equivalently
$e^{g}\notin\leighteenA(p)$ --- only reduced to exhibiting an intermediate subalgebra,
\lean{doubletA_positivity_of_subalgebra}); $\leighteenA(p)_{\bar0}\cong W(p)$ or any relation to \lean{tripletWVOA};
anything at $p\equiv2$ mod $4$ (including $p=2$, where $\leighteenA(2)$ should be the symplectic fermions); any
representation theory. The concrete doublet is therefore, at present, a well-defined bounded-below subalgebra with the
right central charge and nonzero generators, and nothing more. Is it degenerate? It contains $a^-\ne0$ and $a^+\ne0$
(\lean{aMinusC_ne_zero}, \lean{aPlusC_ne_zero}) and $\omega\ne0$, so it is not zero; whether it is a proper subalgebra is
open.
\end{warning}
\begin{physics}
At $4\mid p$ the doublet theory has bosonic fields of half-integral weight $(3p-2)/4$ (e.g.\ $\frac52$ at $p=4$): it is a
vertex algebra but not a VOA in the $\Z$-graded sense, the same phenomenon as for the $\beta\gamma$ system at
half-integral twist. Bounded-belowness, not positivity, is the right finiteness property for such an ambient.
\end{physics}
\begin{status}
Proved, unconditional (only $4\mid p$, $p>0$); standard axioms. What is proved is only what is displayed: a \emph{bounded-below}
generated super near-VOA subalgebra exists for $4\mid p$ (with the stated central charge, trivial parity and failing even-part
hypothesis). It is \textbf{not} proved to be the doublet $\leighteenA(p)$ of the literature, nor that
$\leighteenA(p)^{\mathrm{even}}=W(p)$, and nothing exists for $4\nmid p$.
\end{status}

%======================================================================
\section{The remaining files of this chapter}\label{l18:sec:others}
%======================================================================

For completeness, the files not discussed theorem-by-theorem above, and what they contain.
\begin{itemize}
\item \leanpath{TripletWpAmbient.lean}: besides Definition~\ref{l18:def:AMid}, the transport dictionary
  (\lean{O₂amb}, \lean{mem_O₂amb_iff}, \lean{L_coe}), \lean{strongGenWp_iff} and \lean{SmulClosedWp}
  (strong generation $\iff$ the ambient word span is closed under all products), the lattice product
  \lean{smul_latChargeGen_bottom} ($(e^\alpha)_{(-\langle\alpha,\beta\rangle-1)}e^\beta=\varepsilon(\alpha,\beta)e^{\alpha+\beta}$), and the weight
  bookkeeping forcing $m=3p-1$.
\item \leanpath{TripletWpTwoSector.lean}, \leanpath{TripletWpTwoVirasoroCoords.lean}: finite expansions of the
  charge dressing and of the weight-$10$ Virasoro words in Fock coordinates (no PBW: the Fock space is a polynomial
  ring with its monomial basis).
\item \leanpath{TripletWpTwoMixedAssembly.lean}, \leanpath{TripletWpTwoMixedClose.lean}: assembly of
  Theorem~\ref{l18:thm:mixed}.
\item \leanpath{TripletWpTwoAM\{Assembly,Close,Coords,ExpectedSparse,Proof,Schur,Sparse*\}.lean}: the pieces of
  Theorem~\ref{l18:thm:AMtwo}.
\item \leanpath{TripletWpTwoZhu\{Certificate,Fast*,Sparse*\}.lean} (129 files, counting \leanpath{TripletWpTwoZhuFinal.lean}): the pieces of the $f_2$
  certificate (Theorem~\ref{l18:thm:simple}); \lean{fTwoZhuRepresentative_}\allowbreak\lean{eq_}\allowbreak\lean{intrinsicCertificateRhs_}\allowbreak\lean{of_}\allowbreak\lean{sparse}
  is the interface.
\item \leanpath{TripletWpTwoFelderCarrier.lean}: charged Fock carriers for a future Felder-complex analysis; no
  exactness theorem.
\item \leanpath{TripletW2C2.lean}: the general criteria $C_2\iff$ eventual vanishing modulo $O_2$
  (Theorem~\ref{l18:thm:abeDirect}(a)), Abe's Prop.~2.2 (strong generation $\Rightarrow$ ring generation),
  \lean{vacuumGradeZero_of_VOA} (\lean{VacuumGradeZero} holds for every VOA, so it is redundant wherever it is carried,
  for instance in the cusp theorems of Chapter~\ref{ch:L15}), and Abe's Lemma~3.4 (\lean{BW2_reduce}).
\item \leanpath{ZhuTripletW2.lean}: general Zhu-product calculus on homogeneous left factors
  (\lean{zhuMul_of_mem_grading}: $u*v=\sum_{m\le k}\binom km u_{(m-1)}v$), $[L_{-1}u]=-k[u]$, and the $W(2)$
  specialisations.
\end{itemize}

%======================================================================
\section{Verification record}\label{l18:sec:verification}
%======================================================================

The following was run for this chapter at an earlier commit (not re-run for this edition), with the library's existing build products, on a
scratch file importing the relevant modules (no full \texttt{lake build} was run for this chapter):
\begin{plaincode}
lake env lean -DmaxSynthPendingDepth=3 Axioms.lean
\end{plaincode}
Output (abridged to the declaration names; every line reads
\texttt{depends on axioms: [propext, Classical.choice, Quot.sound]}):
\begin{plaincode}
VOA.LatGen.Triplet.tripletLatticeVOA
VOA.LatGen.Screen.LBG_zero_comm_Qscr          VOA.LatGen.Screen.Hscr_ne_zero
VOA.LatGen.Screen.Qscr_cube_Fscr              VOA.LatGen.Wp.tripletWVOA
VOA.LatGen.Wp.tripletW_carrier_ne_bot         VOA.LatGen.WpC2.mk₂_genH_pow_four
VOA.LatGen.WpC2.c2Cofinite_tripletWVOA_of_omegaPowMem
VOA.LatGen.WpC2.c2Cofinite_tripletWVOA_of_amCor44
VOA.LatGen.WpC2.c2Cofinite_tripletWVOA_of_lowering_AM
VOA.LatGen.WpC2.screenGens_primary            VOA.LatGen.WpC2.threeCoordinateIdentitiesWp_two
VOA.LatGen.WpTwoAMFinal.amOmegaIdentity_two
VOA.LatGen.WpC2.tripletWVOA_two_c2Cofinite_and_finrank_le
VOA.LatGen.Wp.strongGenWp_two
VOA.LatGen.Wp.not_admissiblyRational_tripletWVOA_two
VOA.LatGen.Wp.not_gannonRational_tripletWVOA_two
VOA.LatGen.Wp.tripletWVOA_two_isSimple
VOA.LatGen.Wp.tripletWVOA_two_stronglyFinite_and_irrational
VOA.LatGen.Wp.not_stronglyRational_tripletWVOA_two
VOA.LatGen.WpC2.Qwp_sq_latChargeGen_mem_O
VOA.SympFerm.tripletW2                         VOA.SympFerm.tripletW2_grading_one
VOA.AbeAdjoin.c2Cofinite_tripletW2_of_abeQuotientGeneration
VOA.AbeAdjoin.finrank_A₂_tripletW2_le_96_of_abeQuotientGeneration
VOA.Irrationality.not_admissiblyRational_tripletW2
VOA.Irrationality.tripletW2_not_isSemisimpleRing_A
VOA.SympFerm.tripletW2Log_L0nil_ne_zero
VOA.LogModular.not_admissiblyRational_tripletW2_of_logS
VOA.LogModular.linTerm_notMem_logFreeSpan
VOA.AbeRel.wbar_pow_six                         VOA.AbeRel.Hbar_pow_four
VOA.SympFerm.omega_sq_not_mem_O₂               VOA.AbeGens.c2Cofinite_of_abeSpanningEleven
VOA.AbeGens.abeGens_linearIndependent          VOA.TripletNilp.c2Cofinite_tripletW2
VOA.TripletDirect.c2Cofinite_tripletW2_of_levelOne
VOA.ZhuJordanAlt.not_isOmegaJordanVector_lowestAMod
VOA.pseudoTrace_idSymmetric_eq_trace           VOA.pseudoTrace_mem_symmetricFunctionals
VOA.LogPseudoTraceCommutantData.isLogPseudoTraceCyclicBlock
VOA.LogPseudoTraceCommutantData.isLogPseudoTraceFormalBlock
VOA.DoubletC.doubletApC_boundedBelow           VOA.DoubletC.hint_false_of_four_dvd
VOA.LatGen.Doublet.aPlus_ne_zero
\end{plaincode}
None of these depends on \lean{VOA.pbw}. The Lean checks cited in this chapter (not part of this release) were
compiled the same way and depend only on the same three axioms.
A clean \texttt{\#print axioms} does not report explicit hypotheses: every theorem marked ``conditional'' in this chapter
is \texttt{\#print axioms}-clean \emph{and} conditional on the named \lean{Prop} arguments listed in its status box.

\paragraph{Release audit 2026-10-27: RELEASE\_AUDIT\_L17\_L18\_L23.md.} The independent release audit
\texttt{Notes/\allowbreak Textbook/\allowbreak referee/\allowbreak RELEASE\_AUDIT\_L17\_L18\_L23.md} re-elaborated the headline
statements of this chapter at an earlier commit (branch \texttt{deligne-lean}; \leanpath{VOA/} unchanged since
the audit's start; not re-run for this edition) from the existing build products, with \texttt{lake env lean -DmaxSynthPendingDepth=3} (\texttt{lake build} not
run), re-ran \texttt{\#print axioms} (standard three for \lean{tripletLatticeVOA}, \lean{tripletWVOA}, the $p=2$ lattice theorems,
\lean{c2Cofinite_tripletWVOA}, \lean{tripletW2}, \lean{c2Cofinite_tripletW2_of_abeQuotientGeneration},
\lean{tripletW2Log_jordan_pair}, both \lean{not_admissiblyRational_tripletW2*} and \lean{doubletApC_boundedBelow}), and recompiled
the Lean checks cited in this chapter. Verdicts:
\emph{lattice $W(2)$: announceable with caveats} (lattice model only; state the rationality notions, findings M4, m5);
\emph{general $p$, the Abe model, the logarithmic module, the doublet and pseudo-traces: not announceable as theorems}, only as
conditional reductions or constructions (M4, M5, m7, m8). This chapter's status boxes have been brought into line: every result
on the Abe model, for general $p$, on the logarithmic module and on the doublet is marked conditional (with its hypotheses) or
restricted to what is displayed. The audit also records (M7) that the aggregate build is not certified for the release commit and
(m6) that \lean{tripletW p} is the generated subalgebra, never proved equal to $\ker\widetilde Q$, with $E\ne0$ and
$W(p)\ne V_{L_p}$ proved only at $p=2$, outside the library.

%======================================================================
\section{Recorded traps, dead ends and Lean defects}\label{l18:sec:traps}
%======================================================================

\begin{warning}
\textbf{Two $W(2)$'s.} \lean{tripletWVOA 2 hp} (lattice) and \lean{SympFerm.tripletW2} (Abe) live on different carrier
types; no isomorphism, no homomorphism and no identification of their graded dimensions is proved. A theorem about
one is not a theorem about the other. In particular ``$W(2)$ is $C_2$-cofinite'', ``simple'', ``strongly finite'', ``not
Gannon-rational'' are proved for the lattice model only; ``$A(W(2))$ is not semisimple'' and ``a logarithmic
$W(2)$-module exists'' are (conditionally) about the Abe model only. Constructing the isomorphism (bosonization of
symplectic fermions) would transfer everything.
\end{warning}

\begin{warning}
\textbf{Dead ends recorded in the source} (useful to anyone continuing): (i) Adamovi\'c--Milas' printed proof of
$\bar\omega$-nilpotency (Section~\ref{l18:sec:c2general}); (ii) $\bar\omega^2=0$ in $R_{W(2)}$ is false
(Theorem~\ref{l18:thm:TTnull}); (iii) $\omega,E,H,F$ do \emph{not} span $W(2)$ modulo $C_2$ (\lean{SpanModO₂Abe} would give
$\dim R\le4<9$); (iv) the ambient symplectic-fermion $C_2$ result does not descend to the even part; (v) the top level of
an ordinary module can never certify \lean{OmegaJordanAt} (Theorem~\ref{l18:thm:obstruction}); (vi) ``rational $\Rightarrow$
$\mathrm{tr}\,L_0^{\mathrm{nil}}=0$'' is no obstruction ($L_0^{\mathrm{nil}}$ is nilpotent, so all its traces vanish anyway) --- the
reason pseudo-traces are needed; (vii) the doublet's even part cannot be formed by the super even-part construction for
$4\mid p$ (Theorem~\ref{l18:thm:doublet}).
\end{warning}

\begin{warning}
\textbf{Lean defects found by this chapter}:
\begin{enumerate}
\item (documentation only) \leanpath{TripletWp.lean}, \leanpath{TripletWpC2.lean} docstrings: ``$E\ne0$ not proved'', ``$W(p)\ne V_{L_p}$ not
  proved'' are stale at $p=2$ (Theorem~\ref{l18:thm:Ene0}).
\item (documentation only) \leanpath{TripletScreening.lean}, \leanpath{TripletWp.lean}: $Q=e^\alpha_{(0)}$ called the ``short'' screening; it is
  the long one.
\item (documentation only) \leanpath{TripletWpPairClosure.lean} ``What remains open'' and \leanpath{TripletWpStronglyFinite.lean} (``conditional
  on simplicity'') are stale: the mixed coordinates, the AM identity and simplicity are proved in later files.
\item (documentation only) \leanpath{TripletW2Log.lean}, carried hypothesis 2: says \lean{hAnm} is not implied by $C_2$ in the form needed;
  \lean{tripletW2_zhuFiniteness_of_C2Cofinite} proves that it is.
\item (moderate, minor) \lean{tripletW2_not_isSemisimpleRing_A_of_finrank_eleven},
  \lean{not_admissiblyRational_tripletW2_of_finrank_eleven}: \lean{hfin} is implied by \lean{hdim}.
\item (moderate, harmless) \lean{WedderburnBlockCount ℂ (A tripletW2) 4} is vacuously true for a non-semisimple algebra;
  the hypothesis \lean{hblocks} therefore costs nothing but also carries no information about the real $A(W(2))$.
\item (open gap) \lean{LogPseudoTraceCommutantData} has no instance; the logarithmic genus-one layer applies to no concrete
  VOA.
\item (open gap) The \textsc{Library Map} summary ``general $p$ conditional on \lean{OmegaPowMem}'' omits the second hypothesis
  (\lean{AdjoinGenWp}, or \lean{LoweringPairsWp}); and \lean{tripletW p hp} is never related to $\ker\widetilde Q$.
\item (documentation only) Chapter~\ref{ch:L04}'s remark that the triplet screening fits \lean{Screening.screening} only at $\lambda=0$ is
  superseded by \lean{Qwp}.
\end{enumerate}
\end{warning}

%======================================================================
\section{Referee's checklist}\label{l18:sec:checklist}
%======================================================================

\begingroup\small
\begin{longtable}{>{\raggedright\arraybackslash}p{4.6cm}>{\raggedright\arraybackslash}p{4.4cm}>{\raggedright\arraybackslash}p{3.4cm}p{1.6cm}}
\toprule
Lean declaration & Intended meaning & Hypotheses & Check\\
\midrule\endhead
\lean{tripletLatticeVOA} & $V_{L_p}$ with Feigin--Fuchs $\omega$, $c=c_{1,p}$ & \lean{hp} & Thm~\ref{l18:thm:latVOA}\\
\lean{LBG_zero_comm_Qscr} & $[L^\lambda_0,Q]=0$ & $\wt e^\alpha=1$ & Thm~\ref{l18:thm:screen}\\
\lean{Hscr_ne_zero}, \lean{Qscr_cube_Fscr} & $H\ne0$, $Q^3F=0$ & lattice data & Thm~\ref{l18:thm:screen}\\
\lean{tripletWVOA} & subalgebra generated by $\omega,E,H,F$ & \lean{hp} & Def~\ref{l18:def:tripletW}\\
(Lean check, not in this release) & $E\ne0$, $W(2)\ne V_{L_2}$ & none & Thm~\ref{l18:thm:Ene0}\\
\lean{mk₂_genH_pow_four} & $\bar H^4=0$ in $R_{W(p)}$ & none & Thm~\ref{l18:thm:AMrel}\\
\lean{c2Cofinite_}\allowbreak\lean{tripletWVOA_}\allowbreak\lean{of_}\allowbreak\lean{omegaPowMem} & $W(p)$ is $C_2$-cofinite, $\dim R\le16m$ (conditional) & \lean{AdjoinGenWp}, \lean{OmegaPowMem} & Thm~\ref{l18:thm:c2general}\\
\lean{c2Cofinite_}\allowbreak\lean{tripletWVOA_}\allowbreak\lean{of_}\allowbreak\lean{lowering_}\allowbreak\lean{AM} & same & \lean{LoweringPairsWp}, \lean{AMOmegaIdentity} & Thm~\ref{l18:thm:pairs}\\
\lean{screenGens_primary} & $E,H,F$ primary & none & Thm~\ref{l18:thm:pairs}\\
\lean{strongGenWp_two} & $W(2)$ strongly generated & none & Thm~\ref{l18:thm:mixed}\\
\lean{amOmegaIdentity_two} & weight-10 AM identity, $B_2=\frac{1024}{1575}$ & none & Thm~\ref{l18:thm:AMtwo}\\
\lean{tripletWVOA_}\allowbreak\lean{two_}\allowbreak\lean{c2Cofinite_}\allowbreak\lean{and_}\allowbreak\lean{finrank_}\allowbreak\lean{le} & lattice $W(2)$ $C_2$-cofinite, $\dim R\le80$ & none (lattice model only) & Thm~\ref{l18:thm:c2two}\\
\lean{not_}\allowbreak\lean{gannonRational_}\allowbreak\lean{tripletWVOA_}\allowbreak\lean{two}, \lean{not_}\allowbreak\lean{admissiblyRational_}\allowbreak\lean{tripletWVOA_}\allowbreak\lean{two} & lattice $W(2)$ not rational & none & Thm~\ref{l18:thm:irrtwo}\\
\lean{tripletWVOA_}\allowbreak\lean{two_}\allowbreak\lean{isSimple} & lattice $W(2)$ simple & none & Thm~\ref{l18:thm:simple}\\
\lean{tripletWVOA_}\allowbreak\lean{two_}\allowbreak\lean{stronglyFinite_}\allowbreak\lean{and_}\allowbreak\lean{irrational} & strongly finite, not rational & none & Thm~\ref{l18:thm:simple}\\
\lean{Qwp_}\allowbreak\lean{sq_}\allowbreak\lean{latChargeGen_}\allowbreak\lean{mem_}\allowbreak\lean{O} & $Q^2e^{-2\alpha}\in O(W(p))$ & none & Thm~\ref{l18:thm:zhuscreen}\\
\lean{tripletW2} & $\leighteenSF^+$, $c=-2$ & none & Def~\ref{l18:def:tripletW2}\\
\lean{omega_sq_not_mem_O₂} & $\bar\omega^2\ne0$ (Abe model) & none & Thm~\ref{l18:thm:TTnull}\\
\lean{wbar_pow_six}, \lean{Hbar_pow_four} & nilpotency in $R_{\leighteenSF^+}$ & none & Thm~\ref{l18:thm:abeRel}\\
\lean{c2Cofinite_}\allowbreak\lean{tripletW2_}\allowbreak\lean{of_}\allowbreak\lean{abeQuotientGeneration} & Abe model $C_2$-cofinite & \lean{AbeQuotientGeneration} & Thm~\ref{l18:thm:abeC2}\\
\lean{tripletW2Log_}\allowbreak\lean{jordan_}\allowbreak\lean{pair} & logarithmic $W(2)$-module (conditional; \lean{hnd} essentially assumes the Jordan block) & \lean{hAnm}, \lean{hA}, \lean{hgate}, \lean{hnd} & Thm~\ref{l18:thm:logmodule}\\
\lean{not_}\allowbreak\lean{isOmegaJordanVector_}\allowbreak\lean{lowestAMod} & ordinary tops never certify & none & Thm~\ref{l18:thm:obstruction}\\
\lean{tripletW2_}\allowbreak\lean{not_}\allowbreak\lean{isSemisimpleRing_}\allowbreak\lean{A} & $A(W(2))$ not semisimple & \lean{hC2}, \lean{hblocks}, \lean{hfive} & Thm~\ref{l18:thm:abeirr}\\
\lean{not_}\allowbreak\lean{admissiblyRational_}\allowbreak\lean{tripletW2} & Abe model not rational (conditional) & + \lean{hzhu} (discharged by \lean{zhuAdmissibleSemisimple_holds}) & Thm~\ref{l18:thm:abeirr}\\
\lean{linTerm_}\allowbreak\lean{notMem_}\allowbreak\lean{logFreeSpan} & $\tau B$ not a sum of $T$-eigenfunctions & $B\ne0$, $\nu\ne0$ & Def~\ref{l18:def:logterm}\\
\lean{not_}\allowbreak\lean{admissiblyRational_}\allowbreak\lean{tripletW2_}\allowbreak\lean{of_}\allowbreak\lean{logS} & $\log q$ refutes rationality & \lean{hzhu}, \lean{hC2}, \lean{hlog} & Thm~\ref{l18:thm:logS}\\
\lean{pseudoTrace_}\allowbreak\lean{idSymmetric_}\allowbreak\lean{eq_}\allowbreak\lean{trace} & pseudo-trace $=$ trace over a field & none & Thm~\ref{l18:thm:pseudo}\\
\lean{isLogPseudoTraceFormalBlock} & logarithmic formal block & commutant data (no instance), recursion & Thm~\ref{l18:thm:logblock}\\
\lean{doubletApC_}\allowbreak\lean{boundedBelow} & a generated super subalgebra (not identified with $\leighteenA(p)$) bounded below & $4\mid p$ & Thm~\ref{l18:thm:doublet}\\
\lean{hint_}\allowbreak\lean{false_}\allowbreak\lean{of_}\allowbreak\lean{four_}\allowbreak\lean{dvd} & even part not formable & $4\mid p$ & Thm~\ref{l18:thm:doublet}\\
\bottomrule
\end{longtable}
\endgroup

%======================================================================
\section{Notes and references}\label{l18:sec:notes}
%======================================================================

\begin{itemize}
\item H.~G.~Kausch, \emph{Extended conformal algebras generated by a multiplet of primary fields}, Phys.\ Lett.\ B
  \textbf{259} (1991) 448--455 --- the triplet $W$-algebras. H.~G.~Kausch, \emph{Symplectic fermions}, Nucl.\ Phys.\ B
  \textbf{583} (2000) 513--541, arXiv:hep-th/0003029.
\item M.~R.~Gaberdiel and H.~G.~Kausch, \emph{A rational logarithmic conformal field theory}, Phys.\ Lett.\ B \textbf{386}
  (1996) 131--137, arXiv:hep-th/9606050; and \emph{A local logarithmic conformal field theory}, Nucl.\ Phys.\ B \textbf{538}
  (1999) 631--658.
\item V.~Gurarie, \emph{Logarithmic operators in conformal field theory}, Nucl.\ Phys.\ B \textbf{410} (1993) 535--549 ---
  the origin of logarithmic CFT.
\item B.~L.~Feigin, A.~M.~Gainutdinov, A.~M.~Semikhatov, I.~Yu.~Tipunin, \emph{Modular group representations and fusion in
  logarithmic conformal field theories and in the quantum group center}, Commun.\ Math.\ Phys.\ \textbf{265} (2006) 47--93;
  \emph{Logarithmic extensions of minimal models: characters and modular transformations}, Nucl.\ Phys.\ B \textbf{757}
  (2006) 303--343.
\item J.~Fuchs, S.~Hwang, A.~M.~Semikhatov, I.~Yu.~Tipunin, \emph{Nonsemisimple fusion algebras and the Verlinde formula},
  Commun.\ Math.\ Phys.\ \textbf{247} (2004) 713--742, arXiv:hep-th/0306274 (the $\tau$-term in \eqref{l18:eq:Schi}).
\item M.~Flohr and M.~R.~Gaberdiel, \emph{Logarithmic torus amplitudes}, J.\ Phys.\ A \textbf{39} (2006) 1955--1968,
  arXiv:hep-th/0509075.
\item T.~Abe, \emph{A $\Z_2$-orbifold model of the symplectic fermionic vertex operator superalgebra}, Math.\ Z.\
  \textbf{255} (2007) 755--792, arXiv:math/0503472 --- $C_2$-cofiniteness of $\leighteenSF^+$ (Thm.~3.9, Prop.~3.11), its
  four simple modules (Thm.~4.2).
\item D.~Adamovi\'c and A.~Milas, \emph{On the triplet vertex algebra $W(p)$}, Adv.\ Math.\ \textbf{217} (2008)
  2664--2699, arXiv:0707.1857 --- Prop.~2.3, Thm.~3.1, Cor.~4.6, Thm.~6.1. D.~Adamovi\'c, A.~Milas, \emph{The doublet vertex
  operator superalgebras $\mathcal A(p)$ and $\mathcal A_{2,p}$}, Contemp.\ Math.\ \textbf{602} (2013)
  (the doublet). D.~Adamovi\'c, A.~Milas, \emph{The structure of Zhu's algebras for certain $W$-algebras}, Adv.\ Math.\
  \textbf{227} (2011) 2425--2456, arXiv:1006.5134 ($\dim A(W(p))=6p-1$). D.~Adamovi\'c, \emph{Classification of
  irreducible modules of certain subalgebras of free boson vertex algebra}, J.\ Algebra \textbf{270} (2003) 115--132.
\item N.~Carqueville and M.~Flohr, \emph{Nonmeromorphic operator product expansion and $C_2$-cofiniteness for a family of
  $W$-algebras}, J.\ Phys.\ A \textbf{39} (2006) 951--966, arXiv:math-ph/0508015 --- the null-vector route to $C_2$ (not
  used by the library).
\item M.~Miyamoto, \emph{Modular invariance of vertex operator algebras satisfying $C_2$-cofiniteness}, Duke Math.\ J.\
  \textbf{122} (2004) 51--91 --- pseudo-traces.
\item Y.~Arike, \emph{Some remarks on symmetric linear functions and pseudotrace maps}, Proc.\ Japan Acad.\ \textbf{86}
  (2010), arXiv:1001.2696; Y.~Arike and K.~Nagatomo, \emph{Some remarks on pseudo-trace functions for orbifold models
  associated with symplectic fermions}, Internat.\ J.\ Math.\ \textbf{24} (2013) 1350008, arXiv:1104.0068.
\item K.~Nagatomo and A.~Tsuchiya, \emph{The triplet vertex operator algebra $W(p)$ and the restricted quantum group
  $\bar U_q(sl_2)$ at $q=e^{\pi i/p}$}, Adv.\ Stud.\ Pure Math.\ \textbf{61} (2011), arXiv:0902.4607.
\item T.~Creutzig and D.~Ridout, \emph{Logarithmic conformal field theory: beyond an introduction}, J.\ Phys.\ A
  \textbf{46} (2013) 494006, arXiv:1303.0847 --- the review this chapter's physics sections follow.
\item T.~Creutzig and T.~Gannon, \emph{Logarithmic conformal field theory, log-modular tensor categories and modular forms},
  J.\ Phys.\ A \textbf{50} (2017) 404004; R.~McRae, \emph{On rationality for $C_2$-cofinite vertex operator algebras},
  arXiv:2108.01898 --- strong finiteness.
\item Y.~Zhu, \emph{Modular invariance of characters of vertex operator algebras}, J.\ Amer.\ Math.\ Soc.\ \textbf{9}
  (1996) 237--302; C.~Dong, H.~Li, G.~Mason, \emph{Vertex operator algebras and associative algebras}, J.\ Algebra
  \textbf{206} (1998) 67--96; C.~Dong and C.~Jiang, \emph{Bimodules associated to vertex operator algebras}, Math.\ Z.\
  \textbf{259} (2008) 799--826.
\item Y.-Z.~Huang, J.~Lepowsky, L.~Zhang, \emph{Logarithmic tensor category theory for generalized modules for a conformal
  vertex algebra}, arXiv:1012.4193 --- logarithmic modules.
\end{itemize}

\endgroup
